\documentclass[
  doc,
  floatsintext,
  longtable,
  a4paper,
  nolmodern,
  notxfonts,
  notimes,
  12pt,
  colorlinks=true,linkcolor=blue,citecolor=blue,urlcolor=blue]{apa7}

\usepackage{amsmath}
\usepackage{amssymb}

\geometry{inner=1in, outer=1in}
\fancyhfoffset[LE,RO]{0cm}


\usepackage[bidi=default]{babel}
\babelprovide[main,import]{spanish}


% get rid of language-specific shorthands (see #6817):
\let\LanguageShortHands\languageshorthands
\def\languageshorthands#1{}

\RequirePackage{longtable}
\RequirePackage{threeparttablex}

\makeatletter
\renewcommand{\paragraph}{\@startsection{paragraph}{4}{\parindent}%
	{0\baselineskip \@plus 0.2ex \@minus 0.2ex}%
	{-.5em}%
	{\normalfont\normalsize\bfseries\typesectitle}}

\renewcommand{\subparagraph}[1]{\@startsection{subparagraph}{5}{0.5em}%
	{0\baselineskip \@plus 0.2ex \@minus 0.2ex}%
	{-\z@\relax}%
	{\normalfont\normalsize\bfseries\itshape\hspace{\parindent}{#1}\textit{\addperi}}{\relax}}
\makeatother




\usepackage{longtable, booktabs, multirow, multicol, colortbl, hhline, caption, array, float, xpatch}
\usepackage{subcaption}


\renewcommand\thesubfigure{\Alph{subfigure}}
\setcounter{topnumber}{2}
\setcounter{bottomnumber}{2}
\setcounter{totalnumber}{4}
\renewcommand{\topfraction}{0.85}
\renewcommand{\bottomfraction}{0.85}
\renewcommand{\textfraction}{0.15}
\renewcommand{\floatpagefraction}{0.7}

\usepackage{tcolorbox}
\tcbuselibrary{listings,theorems, breakable, skins}
\usepackage{fontawesome5}

\definecolor{quarto-callout-color}{HTML}{909090}
\definecolor{quarto-callout-note-color}{HTML}{0758E5}
\definecolor{quarto-callout-important-color}{HTML}{CC1914}
\definecolor{quarto-callout-warning-color}{HTML}{EB9113}
\definecolor{quarto-callout-tip-color}{HTML}{00A047}
\definecolor{quarto-callout-caution-color}{HTML}{FC5300}
\definecolor{quarto-callout-color-frame}{HTML}{ACACAC}
\definecolor{quarto-callout-note-color-frame}{HTML}{4582EC}
\definecolor{quarto-callout-important-color-frame}{HTML}{D9534F}
\definecolor{quarto-callout-warning-color-frame}{HTML}{F0AD4E}
\definecolor{quarto-callout-tip-color-frame}{HTML}{02B875}
\definecolor{quarto-callout-caution-color-frame}{HTML}{FD7E14}

%\newlength\Oldarrayrulewidth
%\newlength\Oldtabcolsep


\usepackage{hyperref}



\usepackage{color}
\usepackage{fancyvrb}
\newcommand{\VerbBar}{|}
\newcommand{\VERB}{\Verb[commandchars=\\\{\}]}
\DefineVerbatimEnvironment{Highlighting}{Verbatim}{commandchars=\\\{\}}
% Add ',fontsize=\small' for more characters per line
\usepackage{framed}
\definecolor{shadecolor}{RGB}{241,243,245}
\newenvironment{Shaded}{\begin{snugshade}}{\end{snugshade}}
\newcommand{\AlertTok}[1]{\textcolor[rgb]{0.68,0.00,0.00}{#1}}
\newcommand{\AnnotationTok}[1]{\textcolor[rgb]{0.37,0.37,0.37}{#1}}
\newcommand{\AttributeTok}[1]{\textcolor[rgb]{0.40,0.45,0.13}{#1}}
\newcommand{\BaseNTok}[1]{\textcolor[rgb]{0.68,0.00,0.00}{#1}}
\newcommand{\BuiltInTok}[1]{\textcolor[rgb]{0.00,0.23,0.31}{#1}}
\newcommand{\CharTok}[1]{\textcolor[rgb]{0.13,0.47,0.30}{#1}}
\newcommand{\CommentTok}[1]{\textcolor[rgb]{0.37,0.37,0.37}{#1}}
\newcommand{\CommentVarTok}[1]{\textcolor[rgb]{0.37,0.37,0.37}{\textit{#1}}}
\newcommand{\ConstantTok}[1]{\textcolor[rgb]{0.56,0.35,0.01}{#1}}
\newcommand{\ControlFlowTok}[1]{\textcolor[rgb]{0.00,0.23,0.31}{\textbf{#1}}}
\newcommand{\DataTypeTok}[1]{\textcolor[rgb]{0.68,0.00,0.00}{#1}}
\newcommand{\DecValTok}[1]{\textcolor[rgb]{0.68,0.00,0.00}{#1}}
\newcommand{\DocumentationTok}[1]{\textcolor[rgb]{0.37,0.37,0.37}{\textit{#1}}}
\newcommand{\ErrorTok}[1]{\textcolor[rgb]{0.68,0.00,0.00}{#1}}
\newcommand{\ExtensionTok}[1]{\textcolor[rgb]{0.00,0.23,0.31}{#1}}
\newcommand{\FloatTok}[1]{\textcolor[rgb]{0.68,0.00,0.00}{#1}}
\newcommand{\FunctionTok}[1]{\textcolor[rgb]{0.28,0.35,0.67}{#1}}
\newcommand{\ImportTok}[1]{\textcolor[rgb]{0.00,0.46,0.62}{#1}}
\newcommand{\InformationTok}[1]{\textcolor[rgb]{0.37,0.37,0.37}{#1}}
\newcommand{\KeywordTok}[1]{\textcolor[rgb]{0.00,0.23,0.31}{\textbf{#1}}}
\newcommand{\NormalTok}[1]{\textcolor[rgb]{0.00,0.23,0.31}{#1}}
\newcommand{\OperatorTok}[1]{\textcolor[rgb]{0.37,0.37,0.37}{#1}}
\newcommand{\OtherTok}[1]{\textcolor[rgb]{0.00,0.23,0.31}{#1}}
\newcommand{\PreprocessorTok}[1]{\textcolor[rgb]{0.68,0.00,0.00}{#1}}
\newcommand{\RegionMarkerTok}[1]{\textcolor[rgb]{0.00,0.23,0.31}{#1}}
\newcommand{\SpecialCharTok}[1]{\textcolor[rgb]{0.37,0.37,0.37}{#1}}
\newcommand{\SpecialStringTok}[1]{\textcolor[rgb]{0.13,0.47,0.30}{#1}}
\newcommand{\StringTok}[1]{\textcolor[rgb]{0.13,0.47,0.30}{#1}}
\newcommand{\VariableTok}[1]{\textcolor[rgb]{0.07,0.07,0.07}{#1}}
\newcommand{\VerbatimStringTok}[1]{\textcolor[rgb]{0.13,0.47,0.30}{#1}}
\newcommand{\WarningTok}[1]{\textcolor[rgb]{0.37,0.37,0.37}{\textit{#1}}}

\providecommand{\tightlist}{%
  \setlength{\itemsep}{0pt}\setlength{\parskip}{0pt}}
\usepackage{longtable,booktabs,array}
\newcounter{none} % for unnumbered tables
\usepackage{calc} % for calculating minipage widths
% Correct order of tables after \paragraph or \subparagraph
\usepackage{etoolbox}
\makeatletter
\patchcmd\longtable{\par}{\if@noskipsec\mbox{}\fi\par}{}{}
\makeatother
% Allow footnotes in longtable head/foot
\IfFileExists{footnotehyper.sty}{\usepackage{footnotehyper}}{\usepackage{footnote}}
\makesavenoteenv{longtable}

\usepackage{graphicx}
\makeatletter
\newsavebox\pandoc@box
\newcommand*\pandocbounded[1]{% scales image to fit in text height/width
  \sbox\pandoc@box{#1}%
  \Gscale@div\@tempa{\textheight}{\dimexpr\ht\pandoc@box+\dp\pandoc@box\relax}%
  \Gscale@div\@tempb{\linewidth}{\wd\pandoc@box}%
  \ifdim\@tempb\p@<\@tempa\p@\let\@tempa\@tempb\fi% select the smaller of both
  \ifdim\@tempa\p@<\p@\scalebox{\@tempa}{\usebox\pandoc@box}%
  \else\usebox{\pandoc@box}%
  \fi%
}
% Set default figure placement to htbp
\def\fps@figure{htbp}
\makeatother







\usepackage{newtx}

\defaultfontfeatures{Scale=MatchLowercase}
\defaultfontfeatures[\rmfamily]{Ligatures=TeX,Scale=1}





\title{Métodos de Machine Learning para Clasificación}


\shorttitle{ML PARA CLASIFICACIÓN}


\usepackage{etoolbox}



\ccoppy{\textcopyright~2021}



\author{Edison Achalma}



\affiliation{
{Escuela Profesional de Economía, Universidad Nacional de San Cristóbal
de Huamanga}}




\leftheader{Achalma}

\date{2021-12-06}


\abstract{Este abstract será actualizado una vez que se complete el
contenido final del artículo. }

\keywords{keyword1, keyword2}

\authornote{\par{\addORCIDlink{Edison Achalma}{0000-0001-6996-3364}} 
\par{ }
\par{   El autor no tiene conflictos de interés que revelar.    Los
roles de autor se clasificaron utilizando la taxonomía de roles de
colaborador (CRediT; https://credit.niso.org/) de la siguiente
manera:  Edison Achalma:   conceptualización, metodología, análisis
formal, investigación, recursos, curación de
datos, redacción, visualización, supervisión, administración del
proyecto}
\par{La correspondencia relativa a este artículo debe dirigirse a Edison
Achalma, Escuela Profesional de Economía, Universidad Nacional de San
Cristóbal de Huamanga, Portal Independencia N°
57, Ayacucho, AYA 5001, Perú, Email: \href{mailto:elmer.achalma.09@unsch.edu.pe}{elmer.achalma.09@unsch.edu.pe}}
}

\makeatletter
\let\endoldlt\endlongtable
\def\endlongtable{
\hline
\endoldlt
}
\makeatother

\urlstyle{same}



\hypersetup{
  pdftitle={M\'{e}todos de Machine Learning para Clasificaci\'{o}n},
  pdfauthor={Autor}
}
\makeatletter
\@ifpackageloaded{caption}{}{\usepackage{caption}}
\AtBeginDocument{%
\ifdefined\contentsname
  \renewcommand*\contentsname{Tabla de contenidos}
\else
  \newcommand\contentsname{Tabla de contenidos}
\fi
\ifdefined\listfigurename
  \renewcommand*\listfigurename{Lista de Figuras}
\else
  \newcommand\listfigurename{Lista de Figuras}
\fi
\ifdefined\listtablename
  \renewcommand*\listtablename{Lista de Tablas}
\else
  \newcommand\listtablename{Lista de Tablas}
\fi
\ifdefined\figurename
  \renewcommand*\figurename{Figura}
\else
  \newcommand\figurename{Figura}
\fi
\ifdefined\tablename
  \renewcommand*\tablename{Tabla}
\else
  \newcommand\tablename{Tabla}
\fi
}
\@ifpackageloaded{float}{}{\usepackage{float}}
\floatstyle{ruled}
\@ifundefined{c@chapter}{\newfloat{codelisting}{h}{lop}}{\newfloat{codelisting}{h}{lop}[chapter]}
\floatname{codelisting}{Listado}
\newcommand*\listoflistings{\listof{codelisting}{Lista de Listados}}
\makeatother
\makeatletter
\makeatother
\makeatletter
\@ifpackageloaded{caption}{}{\usepackage{caption}}
\@ifpackageloaded{subcaption}{}{\usepackage{subcaption}}
\makeatother
\makeatletter
\@ifpackageloaded{fontawesome5}{}{\usepackage{fontawesome5}}
\makeatother

% From https://tex.stackexchange.com/a/645996/211326
%%% apa7 doesn't want to add appendix section titles in the toc
%%% let's make it do it
\makeatletter
\xpatchcmd{\appendix}
  {\par}
  {\addcontentsline{toc}{section}{\@currentlabelname}\par}
  {}{}
\makeatother

%% Disable longtable counter
%% https://tex.stackexchange.com/a/248395/211326

\usepackage{etoolbox}

\makeatletter
\patchcmd{\LT@caption}
  {\bgroup}
  {\bgroup\global\LTpatch@captiontrue}
  {}{}
\patchcmd{\longtable}
  {\par}
  {\par\global\LTpatch@captionfalse}
  {}{}
\apptocmd{\endlongtable}
  {\ifLTpatch@caption\else\addtocounter{table}{-1}\fi}
  {}{}
\newif\ifLTpatch@caption
\makeatother

\begin{document}

\maketitle


\hypertarget{toc}{}
\tableofcontents
\newpage
\section[Introduction]{Métodos de Machine Learning para Clasificación}

\setcounter{secnumdepth}{3}

\setlength\LTleft{0pt}




En esta novena guía exploraremos los principales algoritmos de machine
learning para problemas de clasificación. Estos métodos son
fundamentales para resolver problemas económicos como clasificación de
riesgo crediticio, predicción de comportamiento de consumidores,
segmentación de mercados, entre otros.

\section{Introducción al aprendizaje
supervisado}\label{introducciuxf3n-al-aprendizaje-supervisado}

\subsection{Concepto fundamental}\label{concepto-fundamental}

El aprendizaje supervisado es un tipo de machine learning donde el
modelo aprende de datos etiquetados (con respuestas conocidas) para
predecir etiquetas de nuevos datos.

\begin{Shaded}
\begin{Highlighting}[]
\ImportTok{import}\NormalTok{ pandas }\ImportTok{as}\NormalTok{ pd}
\ImportTok{import}\NormalTok{ numpy }\ImportTok{as}\NormalTok{ np}
\ImportTok{import}\NormalTok{ matplotlib.pyplot }\ImportTok{as}\NormalTok{ plt}
\ImportTok{import}\NormalTok{ seaborn }\ImportTok{as}\NormalTok{ sns}
\ImportTok{from}\NormalTok{ matplotlib.colors }\ImportTok{import}\NormalTok{ ListedColormap}

\CommentTok{\# Scikit{-}learn}
\ImportTok{from}\NormalTok{ sklearn.datasets }\ImportTok{import}\NormalTok{ load\_iris, make\_classification}
\ImportTok{from}\NormalTok{ sklearn.model\_selection }\ImportTok{import}\NormalTok{ train\_test\_split, cross\_val\_score}
\ImportTok{from}\NormalTok{ sklearn.preprocessing }\ImportTok{import}\NormalTok{ StandardScaler}
\ImportTok{from}\NormalTok{ sklearn.metrics }\ImportTok{import}\NormalTok{ (confusion\_matrix, classification\_report, }
\NormalTok{                             roc\_curve, auc, accuracy\_score)}

\CommentTok{\# Modelos}
\ImportTok{from}\NormalTok{ sklearn.linear\_model }\ImportTok{import}\NormalTok{ LogisticRegression}
\ImportTok{from}\NormalTok{ sklearn.neighbors }\ImportTok{import}\NormalTok{ KNeighborsClassifier}
\ImportTok{from}\NormalTok{ sklearn.tree }\ImportTok{import}\NormalTok{ DecisionTreeClassifier}
\ImportTok{from}\NormalTok{ sklearn.ensemble }\ImportTok{import}\NormalTok{ RandomForestClassifier}

\CommentTok{\# Visualización de árboles}
\ImportTok{from}\NormalTok{ sklearn }\ImportTok{import}\NormalTok{ tree}

\CommentTok{\# Configuración}
\NormalTok{plt.style.use(}\StringTok{\textquotesingle{}seaborn{-}v0\_8{-}darkgrid\textquotesingle{}}\NormalTok{)}
\NormalTok{plt.rcParams[}\StringTok{\textquotesingle{}figure.figsize\textquotesingle{}}\NormalTok{] }\OperatorTok{=}\NormalTok{ (}\DecValTok{12}\NormalTok{, }\DecValTok{6}\NormalTok{)}
\NormalTok{np.random.seed(}\DecValTok{42}\NormalTok{)}

\BuiltInTok{print}\NormalTok{(}\StringTok{"Librerías importadas exitosamente"}\NormalTok{)}
\end{Highlighting}
\end{Shaded}

\subsection{Tipos de problemas de
clasificación}\label{tipos-de-problemas-de-clasificaciuxf3n}

\begin{Shaded}
\begin{Highlighting}[]
\BuiltInTok{print}\NormalTok{(}\StringTok{"""}
\StringTok{TIPOS DE PROBLEMAS DE CLASIFICACIÓN}

\StringTok{1. CLASIFICACIÓN BINARIA}
\StringTok{   {-} Dos clases posibles}
\StringTok{   {-} Ejemplos económicos:}
\StringTok{     * Aprobado/Rechazado (crédito)}
\StringTok{     * Paga/No paga (incumplimiento)}
\StringTok{     * Éxito/Fracaso (empresa)}
\StringTok{     * Compra/No compra (marketing)}

\StringTok{2. CLASIFICACIÓN MULTI{-}CLASE}
\StringTok{   {-} Más de dos clases}
\StringTok{   {-} Ejemplos económicos:}
\StringTok{     * Clasificación de riesgo (AAA, AA, A, BBB, etc.)}
\StringTok{     * Segmentación de clientes (bajo, medio, alto valor)}
\StringTok{     * Categorías de productos}
\StringTok{     * Niveles socioeconómicos (A, B, C, D, E)}

\StringTok{3. CLASIFICACIÓN MULTI{-}ETIQUETA}
\StringTok{   {-} Múltiples etiquetas simultáneas}
\StringTok{   {-} Ejemplos económicos:}
\StringTok{     * Características de consumidor (múltiples intereses)}
\StringTok{     * Atributos de producto (múltiples categorías)}
\StringTok{"""}\NormalTok{)}
\end{Highlighting}
\end{Shaded}

\subsection{Dataset de ejemplo: Iris}\label{dataset-de-ejemplo-iris}

\begin{Shaded}
\begin{Highlighting}[]
\CommentTok{\# Cargar dataset clásico}
\NormalTok{iris }\OperatorTok{=}\NormalTok{ load\_iris()}

\CommentTok{\# Crear DataFrame}
\NormalTok{df\_iris }\OperatorTok{=}\NormalTok{ pd.DataFrame(}
\NormalTok{    data}\OperatorTok{=}\NormalTok{iris.data,}
\NormalTok{    columns}\OperatorTok{=}\NormalTok{iris.feature\_names}
\NormalTok{)}
\NormalTok{df\_iris[}\StringTok{\textquotesingle{}especie\textquotesingle{}}\NormalTok{] }\OperatorTok{=}\NormalTok{ iris.target}
\NormalTok{df\_iris[}\StringTok{\textquotesingle{}especie\_nombre\textquotesingle{}}\NormalTok{] }\OperatorTok{=}\NormalTok{ df\_iris[}\StringTok{\textquotesingle{}especie\textquotesingle{}}\NormalTok{].}\BuiltInTok{map}\NormalTok{(\{}
    \DecValTok{0}\NormalTok{: }\StringTok{\textquotesingle{}Setosa\textquotesingle{}}\NormalTok{,}
    \DecValTok{1}\NormalTok{: }\StringTok{\textquotesingle{}Versicolor\textquotesingle{}}\NormalTok{, }
    \DecValTok{2}\NormalTok{: }\StringTok{\textquotesingle{}Virginica\textquotesingle{}}
\NormalTok{\})}

\BuiltInTok{print}\NormalTok{(}\StringTok{"Dataset Iris:"}\NormalTok{)}
\BuiltInTok{print}\NormalTok{(df\_iris.head())}
\BuiltInTok{print}\NormalTok{(}\SpecialStringTok{f"}\CharTok{\textbackslash{}n}\SpecialStringTok{Forma: }\SpecialCharTok{\{}\NormalTok{df\_iris}\SpecialCharTok{.}\NormalTok{shape}\SpecialCharTok{\}}\SpecialStringTok{"}\NormalTok{)}
\BuiltInTok{print}\NormalTok{(}\SpecialStringTok{f"}\CharTok{\textbackslash{}n}\SpecialStringTok{Distribución de especies:"}\NormalTok{)}
\BuiltInTok{print}\NormalTok{(df\_iris[}\StringTok{\textquotesingle{}especie\_nombre\textquotesingle{}}\NormalTok{].value\_counts())}

\CommentTok{\# Visualizar (usar solo 2 características)}
\NormalTok{plt.figure(figsize}\OperatorTok{=}\NormalTok{(}\DecValTok{10}\NormalTok{, }\DecValTok{6}\NormalTok{))}
\ControlFlowTok{for}\NormalTok{ especie }\KeywordTok{in}\NormalTok{ df\_iris[}\StringTok{\textquotesingle{}especie\textquotesingle{}}\NormalTok{].unique():}
\NormalTok{    mask }\OperatorTok{=}\NormalTok{ df\_iris[}\StringTok{\textquotesingle{}especie\textquotesingle{}}\NormalTok{] }\OperatorTok{==}\NormalTok{ especie}
\NormalTok{    plt.scatter(}
\NormalTok{        df\_iris.loc[mask, }\StringTok{\textquotesingle{}sepal length (cm)\textquotesingle{}}\NormalTok{],}
\NormalTok{        df\_iris.loc[mask, }\StringTok{\textquotesingle{}sepal width (cm)\textquotesingle{}}\NormalTok{],}
\NormalTok{        label}\OperatorTok{=}\NormalTok{df\_iris.loc[mask, }\StringTok{\textquotesingle{}especie\_nombre\textquotesingle{}}\NormalTok{].iloc[}\DecValTok{0}\NormalTok{],}
\NormalTok{        alpha}\OperatorTok{=}\FloatTok{0.7}\NormalTok{,}
\NormalTok{        s}\OperatorTok{=}\DecValTok{100}
\NormalTok{    )}

\NormalTok{plt.xlabel(}\StringTok{\textquotesingle{}Largo del sépalo (cm)\textquotesingle{}}\NormalTok{)}
\NormalTok{plt.ylabel(}\StringTok{\textquotesingle{}Ancho del sépalo (cm)\textquotesingle{}}\NormalTok{)}
\NormalTok{plt.title(}\StringTok{\textquotesingle{}Dataset Iris {-} Clasificación de especies\textquotesingle{}}\NormalTok{)}
\NormalTok{plt.legend()}
\NormalTok{plt.grid(}\VariableTok{True}\NormalTok{, alpha}\OperatorTok{=}\FloatTok{0.3}\NormalTok{)}
\NormalTok{plt.savefig(}\StringTok{\textquotesingle{}iris\_scatter.png\textquotesingle{}}\NormalTok{, dpi}\OperatorTok{=}\DecValTok{300}\NormalTok{, bbox\_inches}\OperatorTok{=}\StringTok{\textquotesingle{}tight\textquotesingle{}}\NormalTok{)}
\BuiltInTok{print}\NormalTok{(}\StringTok{"}\CharTok{\textbackslash{}n}\StringTok{Gráfico guardado como \textquotesingle{}iris\_scatter.png\textquotesingle{}"}\NormalTok{)}
\end{Highlighting}
\end{Shaded}

\section{Regresión logística}\label{regresiuxf3n-loguxedstica}

\subsection{Fundamento matemático}\label{fundamento-matemuxe1tico}

La regresión logística es uno de los métodos más fundamentales y
ampliamente utilizados para problemas de clasificación binaria. A pesar
de su nombre que incluye ``regresión'', es un algoritmo de clasificación
que predice probabilidades de pertenencia a clases mediante una
transformación no lineal de una combinación lineal de características.

\subsubsection{Modelo de probabilidad}\label{modelo-de-probabilidad}

Para clasificación binaria donde \(Y \in \{0, 1\}\), la regresión
logística modela la probabilidad de que la observación pertenezca a la
clase positiva (\(Y=1\)):

\[
P(Y=1 | \mathbf{X}) = p(\mathbf{X}) = \frac{1}{1 + e^{-(\beta_0 + \beta_1 X_1 + \cdots + \beta_p X_p)}}
\]

Definiendo la combinación lineal: \[
z = \beta_0 + \beta_1 X_1 + \beta_2 X_2 + \cdots + \beta_p X_p = \beta_0 + \mathbf{X}^\top\boldsymbol{\beta}
\]

La probabilidad se expresa como: \[
p(\mathbf{X}) = \frac{1}{1 + e^{-z}} = \frac{e^z}{1 + e^z}
\]

Esta función se conoce como \textbf{función logística} o
\textbf{sigmoide}.

\subsubsection{Función sigmoide}\label{funciuxf3n-sigmoide}

La función sigmoide tiene la forma: \[
\sigma(z) = \frac{1}{1 + e^{-z}}
\]

\textbf{Propiedades clave}:

\begin{enumerate}
\def\labelenumi{\arabic{enumi}.}
\item
  \textbf{Rango}: \(\sigma(z) \in (0, 1)\) para todo
  \(z \in \mathbb{R}\)

  \begin{itemize}
  \tightlist
  \item
    Garantiza que las predicciones sean probabilidades válidas
  \end{itemize}
\item
  \textbf{Simetría}: \(\sigma(-z) = 1 - \sigma(z)\)
\item
  \textbf{Límites asintóticos}: \[
  \lim_{z \to \infty} \sigma(z) = 1, \quad \lim_{z \to -\infty} \sigma(z) = 0
  \]
\item
  \textbf{Punto medio}: \(\sigma(0) = 0.5\)
\item
  \textbf{Derivada}: \[
  \frac{d\sigma(z)}{dz} = \sigma(z)(1 - \sigma(z))
  \] Esta propiedad simplifica el cálculo de gradientes.
\end{enumerate}

\subsubsection{Odds y log-odds (logit)}\label{odds-y-log-odds-logit}

\textbf{Odds (momios)}:

La probabilidad de éxito sobre la probabilidad de fracaso: \[
\text{Odds} = \frac{p(\mathbf{X})}{1 - p(\mathbf{X})} = e^z
\]

\textbf{Log-odds (logit)}:

Tomando logaritmo natural: \[
\log\left(\frac{p(\mathbf{X})}{1-p(\mathbf{X})}\right) = z = \beta_0 + \beta_1 X_1 + \cdots + \beta_p X_p
\]

Esta es la \textbf{transformación logit} que convierte probabilidades
(0,1) al rango completo de números reales \((-\infty, \infty)\).

\textbf{Interpretación crucial}: El modelo es lineal en el log-odds, no
en la probabilidad misma.

\subsubsection{Interpretación de
coeficientes}\label{interpretaciuxf3n-de-coeficientes}

Un incremento unitario en \(X_j\) (manteniendo otras variables
constantes):

\textbf{En log-odds}: \[
\Delta \log(\text{Odds}) = \beta_j
\]

\textbf{En odds}: \[
\text{Odds ratio} = e^{\beta_j}
\]

\begin{itemize}
\tightlist
\item
  Si \(\beta_j > 0\): \(e^{\beta_j} > 1\) → La variable incrementa el
  odds de \(Y=1\)
\item
  Si \(\beta_j < 0\): \(e^{\beta_j} < 1\) → La variable reduce el odds
  de \(Y=1\)
\item
  Si \(\beta_j = 0\): \(e^{\beta_j} = 1\) → La variable no afecta el
  odds
\end{itemize}

\textbf{Efecto marginal en probabilidad}: \[
\frac{\partial p}{\partial X_j} = \beta_j \cdot p(\mathbf{X}) \cdot (1 - p(\mathbf{X}))
\]

El efecto no es constante; depende del nivel de probabilidad actual.

\subsubsection{Estimación: Maximum
Likelihood}\label{estimaciuxf3n-maximum-likelihood}

Los coeficientes se estiman mediante \textbf{máxima verosimilitud}
(MLE).

\textbf{Log-verosimilitud}: \[
\ell(\boldsymbol{\beta}) = \sum_{i=1}^{n} \left[ Y_i \log p(\mathbf{X}_i) + (1-Y_i) \log(1-p(\mathbf{X}_i)) \right]
\]

\textbf{Objetivo}: Maximizar \(\ell(\boldsymbol{\beta})\) mediante
métodos iterativos (Newton-Raphson, gradient descent).

\textbf{Función de pérdida equivalente} (Binary Cross-Entropy): \[
\text{BCE}(\boldsymbol{\beta}) = -\frac{1}{n}\sum_{i=1}^{n} \left[ Y_i \log \hat{p}_i + (1-Y_i) \log(1-\hat{p}_i) \right]
\]

\subsubsection{Frontera de decisión}\label{frontera-de-decisiuxf3n}

La frontera de decisión es donde \(p(\mathbf{X}) = 0.5\): \[
\beta_0 + \beta_1 X_1 + \cdots + \beta_p X_p = 0
\]

En el espacio de características, esta es una \textbf{hipersuperficie
lineal} (hiperplano).

\subsubsection{Ventajas}\label{ventajas}

\begin{itemize}
\tightlist
\item
  Salidas probabilísticas interpretables
\item
  Eficiencia computacional
\item
  Pocos hiperparámetros
\item
  Funciona bien como baseline
\item
  Extensible a clasificación multi-clase
\end{itemize}

\subsubsection{Limitaciones}\label{limitaciones}

\begin{itemize}
\tightlist
\item
  Asume linealidad en log-odds
\item
  No captura relaciones no lineales complejas
\item
  Sensible a outliers
\item
  Requiere features relevantes
\end{itemize}

\subsubsection{Aplicaciones
económicas}\label{aplicaciones-econuxf3micas}

\begin{itemize}
\tightlist
\item
  Predicción de incumplimiento crediticio
\item
  Participación laboral
\item
  Decisiones de compra
\item
  Migración
\item
  Adopción de tecnología
\item
  Éxito de startups
\end{itemize}

\subsection{Implementación}\label{implementaciuxf3n}

\begin{Shaded}
\begin{Highlighting}[]
\CommentTok{\# Preparar datos (clasificación binaria: Virginica vs resto)}
\NormalTok{X }\OperatorTok{=}\NormalTok{ iris.data[:, :}\DecValTok{2}\NormalTok{]  }\CommentTok{\# Solo usar 2 características para visualización}
\NormalTok{y }\OperatorTok{=}\NormalTok{ (iris.target }\OperatorTok{==} \DecValTok{2}\NormalTok{).astype(}\BuiltInTok{int}\NormalTok{)  }\CommentTok{\# 1 si es Virginica, 0 si no}

\CommentTok{\# Dividir datos}
\NormalTok{X\_train, X\_test, y\_train, y\_test }\OperatorTok{=}\NormalTok{ train\_test\_split(}
\NormalTok{    X, y, test\_size}\OperatorTok{=}\FloatTok{0.3}\NormalTok{, random\_state}\OperatorTok{=}\DecValTok{42}
\NormalTok{)}

\BuiltInTok{print}\NormalTok{(}\SpecialStringTok{f"Conjunto de entrenamiento: }\SpecialCharTok{\{}\NormalTok{X\_train}\SpecialCharTok{.}\NormalTok{shape[}\DecValTok{0}\NormalTok{]}\SpecialCharTok{\}}\SpecialStringTok{ muestras"}\NormalTok{)}
\BuiltInTok{print}\NormalTok{(}\SpecialStringTok{f"Conjunto de prueba: }\SpecialCharTok{\{}\NormalTok{X\_test}\SpecialCharTok{.}\NormalTok{shape[}\DecValTok{0}\NormalTok{]}\SpecialCharTok{\}}\SpecialStringTok{ muestras"}\NormalTok{)}
\BuiltInTok{print}\NormalTok{(}\SpecialStringTok{f"}\CharTok{\textbackslash{}n}\SpecialStringTok{Distribución en entrenamiento:"}\NormalTok{)}
\BuiltInTok{print}\NormalTok{(}\SpecialStringTok{f"  Clase 0 (No Virginica): }\SpecialCharTok{\{}\NormalTok{(y\_train }\OperatorTok{==} \DecValTok{0}\NormalTok{)}\SpecialCharTok{.}\BuiltInTok{sum}\NormalTok{()}\SpecialCharTok{\}}\SpecialStringTok{"}\NormalTok{)}
\BuiltInTok{print}\NormalTok{(}\SpecialStringTok{f"  Clase 1 (Virginica): }\SpecialCharTok{\{}\NormalTok{(y\_train }\OperatorTok{==} \DecValTok{1}\NormalTok{)}\SpecialCharTok{.}\BuiltInTok{sum}\NormalTok{()}\SpecialCharTok{\}}\SpecialStringTok{"}\NormalTok{)}

\CommentTok{\# Entrenar modelo}
\NormalTok{modelo\_logit }\OperatorTok{=}\NormalTok{ LogisticRegression(random\_state}\OperatorTok{=}\DecValTok{42}\NormalTok{)}
\NormalTok{modelo\_logit.fit(X\_train, y\_train)}

\BuiltInTok{print}\NormalTok{(}\SpecialStringTok{f"}\CharTok{\textbackslash{}n\textbackslash{}n}\SpecialStringTok{MODELO ENTRENADO"}\NormalTok{)}
\BuiltInTok{print}\NormalTok{(}\StringTok{"="} \OperatorTok{*} \DecValTok{70}\NormalTok{)}
\BuiltInTok{print}\NormalTok{(}\SpecialStringTok{f"Intercepto (β₀): }\SpecialCharTok{\{}\NormalTok{modelo\_logit}\SpecialCharTok{.}\NormalTok{intercept\_[}\DecValTok{0}\NormalTok{]}\SpecialCharTok{:.4f\}}\SpecialStringTok{"}\NormalTok{)}
\BuiltInTok{print}\NormalTok{(}\SpecialStringTok{f"Coeficientes:"}\NormalTok{)}
\BuiltInTok{print}\NormalTok{(}\SpecialStringTok{f"  β₁ (largo sépalo): }\SpecialCharTok{\{}\NormalTok{modelo\_logit}\SpecialCharTok{.}\NormalTok{coef\_[}\DecValTok{0}\NormalTok{, }\DecValTok{0}\NormalTok{]}\SpecialCharTok{:.4f\}}\SpecialStringTok{"}\NormalTok{)}
\BuiltInTok{print}\NormalTok{(}\SpecialStringTok{f"  β₂ (ancho sépalo): }\SpecialCharTok{\{}\NormalTok{modelo\_logit}\SpecialCharTok{.}\NormalTok{coef\_[}\DecValTok{0}\NormalTok{, }\DecValTok{1}\NormalTok{]}\SpecialCharTok{:.4f\}}\SpecialStringTok{"}\NormalTok{)}
\end{Highlighting}
\end{Shaded}

\subsection{Predicción e
interpretación}\label{predicciuxf3n-e-interpretaciuxf3n}

\begin{Shaded}
\begin{Highlighting}[]
\CommentTok{\# Predecir probabilidades}
\NormalTok{probs\_train }\OperatorTok{=}\NormalTok{ modelo\_logit.predict\_proba(X\_train)[:, }\DecValTok{1}\NormalTok{]}
\NormalTok{probs\_test }\OperatorTok{=}\NormalTok{ modelo\_logit.predict\_proba(X\_test)[:, }\DecValTok{1}\NormalTok{]}

\CommentTok{\# Predecir clases}
\NormalTok{y\_pred }\OperatorTok{=}\NormalTok{ modelo\_logit.predict(X\_test)}

\CommentTok{\# Evaluar}
\NormalTok{accuracy }\OperatorTok{=}\NormalTok{ accuracy\_score(y\_test, y\_pred)}
\BuiltInTok{print}\NormalTok{(}\SpecialStringTok{f"}\CharTok{\textbackslash{}n}\SpecialStringTok{Accuracy en test: }\SpecialCharTok{\{}\NormalTok{accuracy}\SpecialCharTok{:.4f\}}\SpecialStringTok{ (}\SpecialCharTok{\{}\NormalTok{accuracy}\OperatorTok{*}\DecValTok{100}\SpecialCharTok{:.1f\}}\SpecialStringTok{\%)"}\NormalTok{)}

\CommentTok{\# Matriz de confusión}
\NormalTok{cm }\OperatorTok{=}\NormalTok{ confusion\_matrix(y\_test, y\_pred)}
\BuiltInTok{print}\NormalTok{(}\SpecialStringTok{f"}\CharTok{\textbackslash{}n}\SpecialStringTok{Matriz de confusión:"}\NormalTok{)}
\BuiltInTok{print}\NormalTok{(cm)}

\CommentTok{\# Reporte de clasificación}
\BuiltInTok{print}\NormalTok{(}\SpecialStringTok{f"}\CharTok{\textbackslash{}n}\SpecialStringTok{Reporte de clasificación:"}\NormalTok{)}
\BuiltInTok{print}\NormalTok{(classification\_report(y\_test, y\_pred, }
\NormalTok{                          target\_names}\OperatorTok{=}\NormalTok{[}\StringTok{\textquotesingle{}No Virginica\textquotesingle{}}\NormalTok{, }\StringTok{\textquotesingle{}Virginica\textquotesingle{}}\NormalTok{]))}

\CommentTok{\# Ejemplo de predicción}
\NormalTok{ejemplo }\OperatorTok{=}\NormalTok{ np.array([[}\FloatTok{6.5}\NormalTok{, }\FloatTok{3.0}\NormalTok{]])  }\CommentTok{\# Nuevo ejemplo}
\NormalTok{prob }\OperatorTok{=}\NormalTok{ modelo\_logit.predict\_proba(ejemplo)[}\DecValTok{0}\NormalTok{, }\DecValTok{1}\NormalTok{]}
\NormalTok{pred }\OperatorTok{=}\NormalTok{ modelo\_logit.predict(ejemplo)[}\DecValTok{0}\NormalTok{]}

\BuiltInTok{print}\NormalTok{(}\SpecialStringTok{f"}\CharTok{\textbackslash{}n\textbackslash{}n}\SpecialStringTok{EJEMPLO DE PREDICCIÓN"}\NormalTok{)}
\BuiltInTok{print}\NormalTok{(}\SpecialStringTok{f"Características: Largo=}\SpecialCharTok{\{}\NormalTok{ejemplo[}\DecValTok{0}\NormalTok{, }\DecValTok{0}\NormalTok{]}\SpecialCharTok{\}}\SpecialStringTok{ cm, Ancho=}\SpecialCharTok{\{}\NormalTok{ejemplo[}\DecValTok{0}\NormalTok{, }\DecValTok{1}\NormalTok{]}\SpecialCharTok{\}}\SpecialStringTok{ cm"}\NormalTok{)}
\BuiltInTok{print}\NormalTok{(}\SpecialStringTok{f"Probabilidad de ser Virginica: }\SpecialCharTok{\{}\NormalTok{prob}\SpecialCharTok{:.4f\}}\SpecialStringTok{ (}\SpecialCharTok{\{}\NormalTok{prob}\OperatorTok{*}\DecValTok{100}\SpecialCharTok{:.1f\}}\SpecialStringTok{\%)"}\NormalTok{)}
\BuiltInTok{print}\NormalTok{(}\SpecialStringTok{f"Predicción: }\SpecialCharTok{\{}\StringTok{\textquotesingle{}Virginica\textquotesingle{}} \ControlFlowTok{if}\NormalTok{ pred }\OperatorTok{==} \DecValTok{1} \ControlFlowTok{else} \StringTok{\textquotesingle{}No Virginica\textquotesingle{}}\SpecialCharTok{\}}\SpecialStringTok{"}\NormalTok{)}
\end{Highlighting}
\end{Shaded}

\subsection{Visualización del modelo}\label{visualizaciuxf3n-del-modelo}

\begin{Shaded}
\begin{Highlighting}[]
\KeywordTok{def}\NormalTok{ plot\_decision\_boundary(modelo, X, y, title, filename):}
    \CommentTok{"""Visualiza la frontera de decisión del modelo"""}
    
\NormalTok{    h }\OperatorTok{=} \FloatTok{0.02}  \CommentTok{\# Tamaño del paso en la malla}
    
    \CommentTok{\# Crear malla}
\NormalTok{    x\_min, x\_max }\OperatorTok{=}\NormalTok{ X[:, }\DecValTok{0}\NormalTok{].}\BuiltInTok{min}\NormalTok{() }\OperatorTok{{-}} \FloatTok{0.5}\NormalTok{, X[:, }\DecValTok{0}\NormalTok{].}\BuiltInTok{max}\NormalTok{() }\OperatorTok{+} \FloatTok{0.5}
\NormalTok{    y\_min, y\_max }\OperatorTok{=}\NormalTok{ X[:, }\DecValTok{1}\NormalTok{].}\BuiltInTok{min}\NormalTok{() }\OperatorTok{{-}} \FloatTok{0.5}\NormalTok{, X[:, }\DecValTok{1}\NormalTok{].}\BuiltInTok{max}\NormalTok{() }\OperatorTok{+} \FloatTok{0.5}
\NormalTok{    xx, yy }\OperatorTok{=}\NormalTok{ np.meshgrid(np.arange(x\_min, x\_max, h),}
\NormalTok{                         np.arange(y\_min, y\_max, h))}
    
    \CommentTok{\# Predecir para cada punto en la malla}
\NormalTok{    Z }\OperatorTok{=}\NormalTok{ modelo.predict\_proba(np.c\_[xx.ravel(), yy.ravel()])[:, }\DecValTok{1}\NormalTok{]}
\NormalTok{    Z }\OperatorTok{=}\NormalTok{ Z.reshape(xx.shape)}
    
    \CommentTok{\# Graficar}
\NormalTok{    plt.figure(figsize}\OperatorTok{=}\NormalTok{(}\DecValTok{12}\NormalTok{, }\DecValTok{8}\NormalTok{))}
    
    \CommentTok{\# Contorno de probabilidades}
\NormalTok{    contour }\OperatorTok{=}\NormalTok{ plt.contourf(xx, yy, Z, levels}\OperatorTok{=}\DecValTok{20}\NormalTok{, cmap}\OperatorTok{=}\StringTok{\textquotesingle{}RdYlBu\_r\textquotesingle{}}\NormalTok{, alpha}\OperatorTok{=}\FloatTok{0.6}\NormalTok{)}
\NormalTok{    plt.colorbar(contour, label}\OperatorTok{=}\StringTok{\textquotesingle{}Probabilidad clase positiva\textquotesingle{}}\NormalTok{)}
    
    \CommentTok{\# Frontera de decisión (probabilidad = 0.5)}
\NormalTok{    plt.contour(xx, yy, Z, levels}\OperatorTok{=}\NormalTok{[}\FloatTok{0.5}\NormalTok{], colors}\OperatorTok{=}\StringTok{\textquotesingle{}black\textquotesingle{}}\NormalTok{, linewidths}\OperatorTok{=}\DecValTok{2}\NormalTok{)}
    
    \CommentTok{\# Puntos de datos}
\NormalTok{    scatter }\OperatorTok{=}\NormalTok{ plt.scatter(X[:, }\DecValTok{0}\NormalTok{], X[:, }\DecValTok{1}\NormalTok{], c}\OperatorTok{=}\NormalTok{y, }
\NormalTok{                         cmap}\OperatorTok{=}\StringTok{\textquotesingle{}RdYlBu\_r\textquotesingle{}}\NormalTok{, edgecolors}\OperatorTok{=}\StringTok{\textquotesingle{}black\textquotesingle{}}\NormalTok{, }
\NormalTok{                         s}\OperatorTok{=}\DecValTok{100}\NormalTok{, alpha}\OperatorTok{=}\FloatTok{0.9}\NormalTok{)}
    
\NormalTok{    plt.xlabel(}\StringTok{\textquotesingle{}Largo del sépalo (cm)\textquotesingle{}}\NormalTok{, fontsize}\OperatorTok{=}\DecValTok{12}\NormalTok{)}
\NormalTok{    plt.ylabel(}\StringTok{\textquotesingle{}Ancho del sépalo (cm)\textquotesingle{}}\NormalTok{, fontsize}\OperatorTok{=}\DecValTok{12}\NormalTok{)}
\NormalTok{    plt.title(title, fontsize}\OperatorTok{=}\DecValTok{14}\NormalTok{, fontweight}\OperatorTok{=}\StringTok{\textquotesingle{}bold\textquotesingle{}}\NormalTok{)}
\NormalTok{    plt.grid(}\VariableTok{True}\NormalTok{, alpha}\OperatorTok{=}\FloatTok{0.3}\NormalTok{)}
    
\NormalTok{    plt.tight\_layout()}
\NormalTok{    plt.savefig(filename, dpi}\OperatorTok{=}\DecValTok{300}\NormalTok{, bbox\_inches}\OperatorTok{=}\StringTok{\textquotesingle{}tight\textquotesingle{}}\NormalTok{)}
    \BuiltInTok{print}\NormalTok{(}\SpecialStringTok{f"Gráfico guardado como \textquotesingle{}}\SpecialCharTok{\{}\NormalTok{filename}\SpecialCharTok{\}}\SpecialStringTok{\textquotesingle{}"}\NormalTok{)}

\CommentTok{\# Visualizar}
\NormalTok{plot\_decision\_boundary(}
\NormalTok{    modelo\_logit, X\_train, y\_train,}
    \StringTok{\textquotesingle{}Regresión Logística {-} Frontera de Decisión\textquotesingle{}}\NormalTok{,}
    \StringTok{\textquotesingle{}logistic\_regression\_boundary.png\textquotesingle{}}
\NormalTok{)}
\end{Highlighting}
\end{Shaded}

\section{K-Nearest Neighbors (KNN)}\label{k-nearest-neighbors-knn}

\subsection{Fundamento del algoritmo}\label{fundamento-del-algoritmo}

K-Nearest Neighbors es un algoritmo de clasificación no paramétrico
basado en la premisa de que observaciones similares (vecinas) tienden a
pertenecer a la misma clase. Es uno de los algoritmos más simples
conceptualmente, pero sorprendentemente efectivo en muchas aplicaciones.

\subsubsection{Principio fundamental}\label{principio-fundamental}

\textbf{Idea central}: ``Dime quiénes son tus vecinos y te diré quién
eres''

Una nueva observación se clasifica según la clase más común entre sus
\(k\) vecinos más cercanos en el espacio de características.

\subsubsection{Definición formal}\label{definiciuxf3n-formal}

Dado:

\begin{itemize}
\tightlist
\item
  Conjunto de entrenamiento
  \(\mathcal{D} = \{(\mathbf{X}_i, Y_i)\}_{i=1}^{n}\)
\item
  Nuevo punto a clasificar \(\mathbf{X}_0\)
\item
  Número de vecinos \(k\)
\end{itemize}

El algoritmo:

\begin{enumerate}
\def\labelenumi{\arabic{enumi}.}
\item
  \textbf{Calcular distancias} a todos los puntos de entrenamiento: \[
  d_i = d(\mathbf{X}_0, \mathbf{X}_i) \quad \text{para } i = 1, \ldots, n
  \]
\item
  \textbf{Identificar} los \(k\) puntos con menores distancias: \[
  \mathcal{N}_k(\mathbf{X}_0) = \{k \text{ vecinos más cercanos}\}
  \]
\item
  \textbf{Votar por mayoría}: \[
  \hat{Y}_0 = \arg\max_{c} \sum_{i \in \mathcal{N}_k(\mathbf{X}_0)} \mathbb{1}(Y_i = c)
  \]

  donde \(\mathbb{1}(\cdot)\) es la función indicadora y \(c\) recorre
  todas las clases posibles.
\end{enumerate}

\subsubsection{Métricas de distancia}\label{muxe9tricas-de-distancia}

La elección de la métrica de distancia es fundamental en KNN:

\textbf{1. Distancia Euclidiana} (la más común): \[
d(\mathbf{X}_i, \mathbf{X}_j) = \sqrt{\sum_{m=1}^{p} (X_{im} - X_{jm})^2} = \|\mathbf{X}_i - \mathbf{X}_j\|_2
\]

\textbf{2. Distancia Manhattan} (city-block): \[
d(\mathbf{X}_i, \mathbf{X}_j) = \sum_{m=1}^{p} |X_{im} - X_{jm}| = \|\mathbf{X}_i - \mathbf{X}_j\|_1
\]

\textbf{3. Distancia Minkowski} (generalización): \[
d(\mathbf{X}_i, \mathbf{X}_j) = \left(\sum_{m=1}^{p} |X_{im} - X_{jm}|^q\right)^{1/q}
\]

donde:

\begin{itemize}
\tightlist
\item
  \(q = 1\): Manhattan
\item
  \(q = 2\): Euclidiana
\item
  \(q \to \infty\): Chebyshev
\end{itemize}

\textbf{4. Distancia de Mahalanobis} (considera correlaciones): \[
d(\mathbf{X}_i, \mathbf{X}_j) = \sqrt{(\mathbf{X}_i - \mathbf{X}_j)^\top \mathbf{\Sigma}^{-1} (\mathbf{X}_i - \mathbf{X}_j)}
\]

\subsubsection{Regla de decisión
probabilística}\label{regla-de-decisiuxf3n-probabiluxedstica}

KNN puede proporcionar probabilidades estimadas:

\[
P(Y = c | \mathbf{X}_0) = \frac{1}{k} \sum_{i \in \mathcal{N}_k(\mathbf{X}_0)} \mathbb{1}(Y_i = c)
\]

Es decir, la proporción de vecinos que pertenecen a la clase \(c\).

\subsubsection{Versión ponderada}\label{versiuxf3n-ponderada}

En lugar de votos iguales, se pueden usar \textbf{pesos inversamente
proporcionales a la distancia}:

\[
\hat{Y}_0 = \arg\max_{c} \sum_{i \in \mathcal{N}_k(\mathbf{X}_0)} w_i \cdot \mathbb{1}(Y_i = c)
\]

donde: \[
w_i = \frac{1}{d(\mathbf{X}_0, \mathbf{X}_i) + \epsilon}
\]

\(\epsilon > 0\) evita división por cero. Vecinos más cercanos tienen
más influencia.

\subsubsection{Parámetro k: Trade-off
sesgo-varianza}\label{paruxe1metro-k-trade-off-sesgo-varianza}

\textbf{k pequeño} (ej: k=1):

\begin{itemize}
\tightlist
\item
  \textbf{Baja bias}: Frontera de decisión flexible, se adapta a
  detalles locales
\item
  \textbf{Alta varianza}: Sensible a ruido, predicciones inestables
\item
  Riesgo de \textbf{overfitting}
\item
  Fronteras irregulares
\end{itemize}

\textbf{k grande} (ej: k=50):

\begin{itemize}
\tightlist
\item
  \textbf{Alta bias}: Frontera de decisión suave, ignora detalles
  locales
\item
  \textbf{Baja varianza}: Predicciones estables
\item
  Riesgo de \textbf{underfitting}
\item
  Fronteras más regulares
\end{itemize}

\textbf{k óptimo}: Balance entre sesgo y varianza, determinado
típicamente por validación cruzada.

\subsubsection{Complejidad
computacional}\label{complejidad-computacional}

\textbf{Búsqueda naive}:

\begin{itemize}
\tightlist
\item
  \textbf{Entrenamiento}: \(O(1)\) (solo almacenar datos)
\item
  \textbf{Predicción}: \(O(np)\) por cada nueva observación

  \begin{itemize}
  \tightlist
  \item
    Calcular \(n\) distancias de dimensión \(p\)
  \item
    Encontrar los \(k\) mínimos: \(O(n \log k)\)
  \end{itemize}
\end{itemize}

\textbf{Total por predicción}: \(O(np + n \log k) \approx O(np)\)

\textbf{Estructuras eficientes} (KD-trees, Ball trees):

\begin{itemize}
\tightlist
\item
  \textbf{Construcción}: \(O(np \log n)\)
\item
  \textbf{Predicción}: \(O(\log n)\) en promedio (mejor caso)
\end{itemize}

Sin embargo, en alta dimensionalidad (\(p\) grande), estas estructuras
pierden eficiencia.

\subsubsection{Maldición de la
dimensionalidad}\label{maldiciuxf3n-de-la-dimensionalidad}

En espacios de alta dimensión, KNN sufre severamente:

\textbf{Problema 1}: Las distancias se vuelven uniformes \[
\lim_{p \to \infty} \frac{d_{\max} - d_{\min}}{d_{\min}} \to 0
\]

Todos los puntos parecen equidistantes, haciendo que el concepto de
``vecino cercano'' pierda significado.

\textbf{Problema 2}: Volumen crece exponencialmente

Para mantener densidad constante de vecinos, el volumen requerido crece
como \(r^p\).

\textbf{Consecuencia}: KNN funciona mejor cuando \(p < 10-20\).

\subsubsection{Sensibilidad a la escala}\label{sensibilidad-a-la-escala}

KNN es \textbf{extremadamente sensible} a la escala de las variables.

\textbf{Problema}: Una variable con rango {[}0, 1000{]} dominará una con
rango {[}0, 1{]} en el cálculo de distancias.

\textbf{Solución obligatoria}: Estandarización \[
\tilde{X}_{jm} = \frac{X_{jm} - \mu_m}{\sigma_m}
\]

Sin estandarización, el algoritmo es inválido.

\subsubsection{Fronteras de decisión}\label{fronteras-de-decisiuxf3n}

Las fronteras de decisión de KNN son:

\begin{itemize}
\tightlist
\item
  \textbf{No lineales}: Pueden ser arbitrariamente complejas
\item
  \textbf{No paramétricas}: No asumen forma funcional
\item
  \textbf{Locales}: Definidas solo por vecinos cercanos
\item
  \textbf{Irregulares}: Especialmente con k pequeño
\end{itemize}

Para clasificación binaria con k=1, la frontera es la \textbf{diagrama
de Voronoi} del conjunto de entrenamiento.

\subsubsection{Propiedades teóricas}\label{propiedades-teuxf3ricas}

\textbf{Teorema (Cover \& Hart, 1967)}:

Para k=1 y \(n \to \infty\), el error de KNN está acotado por: \[
\text{Error}_{1-NN} \leq 2 \cdot \text{Error}_{\text{Bayes}}
\]

donde Error\(_{\text{Bayes}}\) es el error del clasificador óptimo.

\textbf{Consistencia}: KNN es consistente si \(k \to \infty\) y
\(k/n \to 0\) cuando \(n \to \infty\).

\subsubsection{Ventajas}\label{ventajas-1}

\begin{enumerate}
\def\labelenumi{\arabic{enumi}.}
\tightlist
\item
  \textbf{Simplicidad conceptual}: Muy intuitivo
\item
  \textbf{No paramétrico}: No asume distribución de datos
\item
  \textbf{Flexibilidad}: Captura fronteras complejas
\item
  \textbf{Lazy learning}: No hay fase de entrenamiento
\item
  \textbf{Actualización trivial}: Fácil incorporar nuevos datos
\item
  \textbf{Múltiples clases}: Extiende naturalmente a clasificación
  multi-clase
\end{enumerate}

\subsubsection{Limitaciones}\label{limitaciones-1}

\begin{enumerate}
\def\labelenumi{\arabic{enumi}.}
\tightlist
\item
  \textbf{Costo computacional alto en predicción}: \(O(np)\) por cada
  predicción
\item
  \textbf{Sensibilidad a escala}: Requiere estandarización obligatoria
\item
  \textbf{Maldición dimensionalidad}: Falla en espacios de alta
  dimensión
\item
  \textbf{Memoria}: Debe almacenar todo el conjunto de entrenamiento
\item
  \textbf{Sensibilidad a features irrelevantes}: Variables ruidosas
  contaminan distancias
\item
  \textbf{No interpretable}: No hay modelo explícito para interpretar
\item
  \textbf{Elección de k}: Requiere validación cruzada
\end{enumerate}

\subsubsection{Selección del parámetro
k}\label{selecciuxf3n-del-paruxe1metro-k}

\textbf{Validación cruzada}: \[
k^* = \arg\min_{k} \frac{1}{K} \sum_{j=1}^{K} \text{Error}_j(k)
\]

\textbf{Regla empírica}: \[
k \approx \sqrt{n}
\]

\textbf{Recomendación}: Probar valores impares para evitar empates en
clasificación binaria.

\subsubsection{Extensiones y variantes}\label{extensiones-y-variantes}

\textbf{1. Weighted KNN}: Pesos basados en distancia o kernel.

\textbf{2. Distance-weighted KNN}: \[
w_i = \exp\left(-\frac{d_i^2}{2\sigma^2}\right)
\]

\textbf{3. Adaptive KNN}: Diferentes valores de \(k\) para diferentes
regiones del espacio.

\textbf{4. Locally Weighted Learning}: Combina KNN con regresión local.

\subsubsection{Aplicaciones en
economía}\label{aplicaciones-en-economuxeda}

\begin{enumerate}
\def\labelenumi{\arabic{enumi}.}
\item
  \textbf{Sistemas de recomendación}: Recomendar productos basados en
  consumidores similares
\item
  \textbf{Credit scoring}: Clasificar aplicantes basándose en aplicantes
  históricos similares
\item
  \textbf{Detección de fraude}: Identificar transacciones anómalas
  (diferentes de vecinos)
\item
  \textbf{Valuación inmobiliaria}: Precio de casas basado en propiedades
  similares
\item
  \textbf{Segmentación de clientes}: Agrupar clientes con
  características similares
\item
  \textbf{Predicción de quiebra}: Clasificar empresas basándose en
  empresas similares del pasado
\end{enumerate}

\subsubsection{Comparación con métodos
paramétricos}\label{comparaciuxf3n-con-muxe9todos-paramuxe9tricos}

{\def\LTcaptype{none} % do not increment counter
\begin{longtable}[]{@{}lll@{}}
\toprule\noalign{}
Aspecto & KNN & Métodos Paramétricos \\
\midrule\noalign{}
\endhead
\bottomrule\noalign{}
\endlastfoot
Supuestos & Mínimos & Forma funcional específica \\
Interpretabilidad & Baja & Alta \\
Flexibilidad & Alta & Media-Baja \\
Costo computacional & Alto (predicción) & Bajo (predicción) \\
Dimensionalidad & Sufre con p alto & Mejor con regularización \\
Fronteras & No lineales complejas & Lineales o simples \\
\end{longtable}
}

\subsection{Implementación}\label{implementaciuxf3n-1}

\begin{Shaded}
\begin{Highlighting}[]
\CommentTok{\# Estandarizar características (importante para KNN)}
\NormalTok{scaler }\OperatorTok{=}\NormalTok{ StandardScaler()}
\NormalTok{X\_train\_scaled }\OperatorTok{=}\NormalTok{ scaler.fit\_transform(X\_train)}
\NormalTok{X\_test\_scaled }\OperatorTok{=}\NormalTok{ scaler.transform(X\_test)}

\CommentTok{\# Entrenar modelos con diferentes valores de k}
\NormalTok{k\_values }\OperatorTok{=}\NormalTok{ [}\DecValTok{1}\NormalTok{, }\DecValTok{3}\NormalTok{, }\DecValTok{5}\NormalTok{, }\DecValTok{10}\NormalTok{, }\DecValTok{15}\NormalTok{, }\DecValTok{20}\NormalTok{]}
\NormalTok{resultados\_knn }\OperatorTok{=}\NormalTok{ []}

\BuiltInTok{print}\NormalTok{(}\StringTok{"EVALUACIÓN DE KNN CON DIFERENTES VALORES DE K"}\NormalTok{)}
\BuiltInTok{print}\NormalTok{(}\StringTok{"="} \OperatorTok{*} \DecValTok{70}\NormalTok{)}
\BuiltInTok{print}\NormalTok{(}\SpecialStringTok{f"}\SpecialCharTok{\{}\StringTok{\textquotesingle{}K\textquotesingle{}}\SpecialCharTok{:\textless{}5\}}\SpecialStringTok{ }\SpecialCharTok{\{}\StringTok{\textquotesingle{}Accuracy Train\textquotesingle{}}\SpecialCharTok{:\textless{}15\}}\SpecialStringTok{ }\SpecialCharTok{\{}\StringTok{\textquotesingle{}Accuracy Test\textquotesingle{}}\SpecialCharTok{:\textless{}15\}}\SpecialStringTok{ }\SpecialCharTok{\{}\StringTok{\textquotesingle{}Diferencia\textquotesingle{}}\SpecialCharTok{:\textless{}12\}}\SpecialStringTok{"}\NormalTok{)}
\BuiltInTok{print}\NormalTok{(}\StringTok{"{-}"} \OperatorTok{*} \DecValTok{70}\NormalTok{)}

\ControlFlowTok{for}\NormalTok{ k }\KeywordTok{in}\NormalTok{ k\_values:}
    \CommentTok{\# Entrenar modelo}
\NormalTok{    modelo\_knn }\OperatorTok{=}\NormalTok{ KNeighborsClassifier(n\_neighbors}\OperatorTok{=}\NormalTok{k)}
\NormalTok{    modelo\_knn.fit(X\_train\_scaled, y\_train)}
    
    \CommentTok{\# Evaluar}
\NormalTok{    acc\_train }\OperatorTok{=}\NormalTok{ modelo\_knn.score(X\_train\_scaled, y\_train)}
\NormalTok{    acc\_test }\OperatorTok{=}\NormalTok{ modelo\_knn.score(X\_test\_scaled, y\_test)}
\NormalTok{    diff }\OperatorTok{=}\NormalTok{ acc\_train }\OperatorTok{{-}}\NormalTok{ acc\_test}
    
\NormalTok{    resultados\_knn.append(\{}
        \StringTok{\textquotesingle{}k\textquotesingle{}}\NormalTok{: k,}
        \StringTok{\textquotesingle{}acc\_train\textquotesingle{}}\NormalTok{: acc\_train,}
        \StringTok{\textquotesingle{}acc\_test\textquotesingle{}}\NormalTok{: acc\_test,}
        \StringTok{\textquotesingle{}diff\textquotesingle{}}\NormalTok{: diff}
\NormalTok{    \})}
    
    \BuiltInTok{print}\NormalTok{(}\SpecialStringTok{f"}\SpecialCharTok{\{}\NormalTok{k}\SpecialCharTok{:\textless{}5\}}\SpecialStringTok{ }\SpecialCharTok{\{}\NormalTok{acc\_train}\SpecialCharTok{:\textless{}15.4f\}}\SpecialStringTok{ }\SpecialCharTok{\{}\NormalTok{acc\_test}\SpecialCharTok{:\textless{}15.4f\}}\SpecialStringTok{ }\SpecialCharTok{\{}\NormalTok{diff}\SpecialCharTok{:\textless{}12.4f\}}\SpecialStringTok{"}\NormalTok{)}

\CommentTok{\# Encontrar mejor k}
\NormalTok{df\_knn }\OperatorTok{=}\NormalTok{ pd.DataFrame(resultados\_knn)}
\NormalTok{mejor\_k }\OperatorTok{=}\NormalTok{ df\_knn.loc[df\_knn[}\StringTok{\textquotesingle{}acc\_test\textquotesingle{}}\NormalTok{].idxmax(), }\StringTok{\textquotesingle{}k\textquotesingle{}}\NormalTok{]}
\BuiltInTok{print}\NormalTok{(}\SpecialStringTok{f"}\CharTok{\textbackslash{}n\textbackslash{}n}\SpecialStringTok{Mejor k: }\SpecialCharTok{\{}\BuiltInTok{int}\NormalTok{(mejor\_k)}\SpecialCharTok{\}}\SpecialStringTok{ (mayor accuracy en test)"}\NormalTok{)}
\end{Highlighting}
\end{Shaded}

\subsection{Visualización del efecto de
k}\label{visualizaciuxf3n-del-efecto-de-k}

\begin{Shaded}
\begin{Highlighting}[]
\CommentTok{\# Entrenar modelo con mejor k}
\NormalTok{modelo\_knn\_final }\OperatorTok{=}\NormalTok{ KNeighborsClassifier(n\_neighbors}\OperatorTok{=}\BuiltInTok{int}\NormalTok{(mejor\_k))}
\NormalTok{modelo\_knn\_final.fit(X\_train\_scaled, y\_train)}

\CommentTok{\# Graficar accuracy vs k}
\NormalTok{fig, axes }\OperatorTok{=}\NormalTok{ plt.subplots(}\DecValTok{1}\NormalTok{, }\DecValTok{2}\NormalTok{, figsize}\OperatorTok{=}\NormalTok{(}\DecValTok{15}\NormalTok{, }\DecValTok{5}\NormalTok{))}

\CommentTok{\# Accuracy}
\NormalTok{axes[}\DecValTok{0}\NormalTok{].plot(df\_knn[}\StringTok{\textquotesingle{}k\textquotesingle{}}\NormalTok{], df\_knn[}\StringTok{\textquotesingle{}acc\_train\textquotesingle{}}\NormalTok{], }\StringTok{\textquotesingle{}o{-}\textquotesingle{}}\NormalTok{, }
\NormalTok{             label}\OperatorTok{=}\StringTok{\textquotesingle{}Training\textquotesingle{}}\NormalTok{, linewidth}\OperatorTok{=}\DecValTok{2}\NormalTok{, markersize}\OperatorTok{=}\DecValTok{8}\NormalTok{)}
\NormalTok{axes[}\DecValTok{0}\NormalTok{].plot(df\_knn[}\StringTok{\textquotesingle{}k\textquotesingle{}}\NormalTok{], df\_knn[}\StringTok{\textquotesingle{}acc\_test\textquotesingle{}}\NormalTok{], }\StringTok{\textquotesingle{}s{-}\textquotesingle{}}\NormalTok{, }
\NormalTok{             label}\OperatorTok{=}\StringTok{\textquotesingle{}Test\textquotesingle{}}\NormalTok{, linewidth}\OperatorTok{=}\DecValTok{2}\NormalTok{, markersize}\OperatorTok{=}\DecValTok{8}\NormalTok{)}
\NormalTok{axes[}\DecValTok{0}\NormalTok{].axvline(x}\OperatorTok{=}\NormalTok{mejor\_k, color}\OperatorTok{=}\StringTok{\textquotesingle{}red\textquotesingle{}}\NormalTok{, linestyle}\OperatorTok{=}\StringTok{\textquotesingle{}{-}{-}\textquotesingle{}}\NormalTok{, }
\NormalTok{                label}\OperatorTok{=}\SpecialStringTok{f\textquotesingle{}Mejor k=}\SpecialCharTok{\{}\BuiltInTok{int}\NormalTok{(mejor\_k)}\SpecialCharTok{\}}\SpecialStringTok{\textquotesingle{}}\NormalTok{)}
\NormalTok{axes[}\DecValTok{0}\NormalTok{].set\_xlabel(}\StringTok{\textquotesingle{}Número de vecinos (k)\textquotesingle{}}\NormalTok{)}
\NormalTok{axes[}\DecValTok{0}\NormalTok{].set\_ylabel(}\StringTok{\textquotesingle{}Accuracy\textquotesingle{}}\NormalTok{)}
\NormalTok{axes[}\DecValTok{0}\NormalTok{].set\_title(}\StringTok{\textquotesingle{}Accuracy vs K\textquotesingle{}}\NormalTok{)}
\NormalTok{axes[}\DecValTok{0}\NormalTok{].legend()}
\NormalTok{axes[}\DecValTok{0}\NormalTok{].grid(}\VariableTok{True}\NormalTok{, alpha}\OperatorTok{=}\FloatTok{0.3}\NormalTok{)}

\CommentTok{\# Overfitting}
\NormalTok{axes[}\DecValTok{1}\NormalTok{].plot(df\_knn[}\StringTok{\textquotesingle{}k\textquotesingle{}}\NormalTok{], df\_knn[}\StringTok{\textquotesingle{}diff\textquotesingle{}}\NormalTok{], }\StringTok{\textquotesingle{}o{-}\textquotesingle{}}\NormalTok{, }
\NormalTok{             linewidth}\OperatorTok{=}\DecValTok{2}\NormalTok{, markersize}\OperatorTok{=}\DecValTok{8}\NormalTok{, color}\OperatorTok{=}\StringTok{\textquotesingle{}purple\textquotesingle{}}\NormalTok{)}
\NormalTok{axes[}\DecValTok{1}\NormalTok{].axhline(y}\OperatorTok{=}\DecValTok{0}\NormalTok{, color}\OperatorTok{=}\StringTok{\textquotesingle{}black\textquotesingle{}}\NormalTok{, linestyle}\OperatorTok{=}\StringTok{\textquotesingle{}{-}{-}\textquotesingle{}}\NormalTok{, alpha}\OperatorTok{=}\FloatTok{0.3}\NormalTok{)}
\NormalTok{axes[}\DecValTok{1}\NormalTok{].axvline(x}\OperatorTok{=}\NormalTok{mejor\_k, color}\OperatorTok{=}\StringTok{\textquotesingle{}red\textquotesingle{}}\NormalTok{, linestyle}\OperatorTok{=}\StringTok{\textquotesingle{}{-}{-}\textquotesingle{}}\NormalTok{, }
\NormalTok{                label}\OperatorTok{=}\SpecialStringTok{f\textquotesingle{}Mejor k=}\SpecialCharTok{\{}\BuiltInTok{int}\NormalTok{(mejor\_k)}\SpecialCharTok{\}}\SpecialStringTok{\textquotesingle{}}\NormalTok{)}
\NormalTok{axes[}\DecValTok{1}\NormalTok{].set\_xlabel(}\StringTok{\textquotesingle{}Número de vecinos (k)\textquotesingle{}}\NormalTok{)}
\NormalTok{axes[}\DecValTok{1}\NormalTok{].set\_ylabel(}\StringTok{\textquotesingle{}Train Accuracy {-} Test Accuracy\textquotesingle{}}\NormalTok{)}
\NormalTok{axes[}\DecValTok{1}\NormalTok{].set\_title(}\StringTok{\textquotesingle{}Diferencia Train{-}Test (Indicador de Overfitting)\textquotesingle{}}\NormalTok{)}
\NormalTok{axes[}\DecValTok{1}\NormalTok{].legend()}
\NormalTok{axes[}\DecValTok{1}\NormalTok{].grid(}\VariableTok{True}\NormalTok{, alpha}\OperatorTok{=}\FloatTok{0.3}\NormalTok{)}

\NormalTok{plt.tight\_layout()}
\NormalTok{plt.savefig(}\StringTok{\textquotesingle{}knn\_k\_selection.png\textquotesingle{}}\NormalTok{, dpi}\OperatorTok{=}\DecValTok{300}\NormalTok{, bbox\_inches}\OperatorTok{=}\StringTok{\textquotesingle{}tight\textquotesingle{}}\NormalTok{)}
\BuiltInTok{print}\NormalTok{(}\StringTok{"}\CharTok{\textbackslash{}n}\StringTok{Gráfico guardado como \textquotesingle{}knn\_k\_selection.png\textquotesingle{}"}\NormalTok{)}
\end{Highlighting}
\end{Shaded}

\subsection{Visualización de fronteras de
decisión}\label{visualizaciuxf3n-de-fronteras-de-decisiuxf3n}

\begin{Shaded}
\begin{Highlighting}[]
\CommentTok{\# Comparar diferentes valores de k}
\NormalTok{fig, axes }\OperatorTok{=}\NormalTok{ plt.subplots(}\DecValTok{2}\NormalTok{, }\DecValTok{3}\NormalTok{, figsize}\OperatorTok{=}\NormalTok{(}\DecValTok{18}\NormalTok{, }\DecValTok{12}\NormalTok{))}
\NormalTok{axes }\OperatorTok{=}\NormalTok{ axes.ravel()}

\NormalTok{h }\OperatorTok{=} \FloatTok{0.02}
\NormalTok{x\_min, x\_max }\OperatorTok{=}\NormalTok{ X\_train\_scaled[:, }\DecValTok{0}\NormalTok{].}\BuiltInTok{min}\NormalTok{() }\OperatorTok{{-}} \DecValTok{1}\NormalTok{, X\_train\_scaled[:, }\DecValTok{0}\NormalTok{].}\BuiltInTok{max}\NormalTok{() }\OperatorTok{+} \DecValTok{1}
\NormalTok{y\_min, y\_max }\OperatorTok{=}\NormalTok{ X\_train\_scaled[:, }\DecValTok{1}\NormalTok{].}\BuiltInTok{min}\NormalTok{() }\OperatorTok{{-}} \DecValTok{1}\NormalTok{, X\_train\_scaled[:, }\DecValTok{1}\NormalTok{].}\BuiltInTok{max}\NormalTok{() }\OperatorTok{+} \DecValTok{1}
\NormalTok{xx, yy }\OperatorTok{=}\NormalTok{ np.meshgrid(np.arange(x\_min, x\_max, h),}
\NormalTok{                     np.arange(y\_min, y\_max, h))}

\ControlFlowTok{for}\NormalTok{ idx, k }\KeywordTok{in} \BuiltInTok{enumerate}\NormalTok{([}\DecValTok{1}\NormalTok{, }\DecValTok{3}\NormalTok{, }\DecValTok{5}\NormalTok{, }\DecValTok{10}\NormalTok{, }\DecValTok{15}\NormalTok{, }\DecValTok{20}\NormalTok{]):}
    \CommentTok{\# Entrenar modelo}
\NormalTok{    knn }\OperatorTok{=}\NormalTok{ KNeighborsClassifier(n\_neighbors}\OperatorTok{=}\NormalTok{k)}
\NormalTok{    knn.fit(X\_train\_scaled, y\_train)}
    
    \CommentTok{\# Predecir}
\NormalTok{    Z }\OperatorTok{=}\NormalTok{ knn.predict(np.c\_[xx.ravel(), yy.ravel()])}
\NormalTok{    Z }\OperatorTok{=}\NormalTok{ Z.reshape(xx.shape)}
    
    \CommentTok{\# Graficar}
\NormalTok{    axes[idx].contourf(xx, yy, Z, alpha}\OperatorTok{=}\FloatTok{0.4}\NormalTok{, cmap}\OperatorTok{=}\StringTok{\textquotesingle{}RdYlBu\_r\textquotesingle{}}\NormalTok{)}
\NormalTok{    axes[idx].scatter(X\_train\_scaled[:, }\DecValTok{0}\NormalTok{], X\_train\_scaled[:, }\DecValTok{1}\NormalTok{], }
\NormalTok{                     c}\OperatorTok{=}\NormalTok{y\_train, cmap}\OperatorTok{=}\StringTok{\textquotesingle{}RdYlBu\_r\textquotesingle{}}\NormalTok{, }
\NormalTok{                     edgecolors}\OperatorTok{=}\StringTok{\textquotesingle{}black\textquotesingle{}}\NormalTok{, s}\OperatorTok{=}\DecValTok{50}\NormalTok{)}
    
\NormalTok{    acc }\OperatorTok{=}\NormalTok{ knn.score(X\_test\_scaled, y\_test)}
\NormalTok{    axes[idx].set\_title(}\SpecialStringTok{f\textquotesingle{}k=}\SpecialCharTok{\{}\NormalTok{k}\SpecialCharTok{\}}\SpecialStringTok{, Accuracy=}\SpecialCharTok{\{}\NormalTok{acc}\SpecialCharTok{:.3f\}}\SpecialStringTok{\textquotesingle{}}\NormalTok{)}
\NormalTok{    axes[idx].set\_xlabel(}\StringTok{\textquotesingle{}Feature 1 (escalado)\textquotesingle{}}\NormalTok{)}
\NormalTok{    axes[idx].set\_ylabel(}\StringTok{\textquotesingle{}Feature 2 (escalado)\textquotesingle{}}\NormalTok{)}
\NormalTok{    axes[idx].grid(}\VariableTok{True}\NormalTok{, alpha}\OperatorTok{=}\FloatTok{0.3}\NormalTok{)}

\NormalTok{plt.tight\_layout()}
\NormalTok{plt.savefig(}\StringTok{\textquotesingle{}knn\_comparison.png\textquotesingle{}}\NormalTok{, dpi}\OperatorTok{=}\DecValTok{300}\NormalTok{, bbox\_inches}\OperatorTok{=}\StringTok{\textquotesingle{}tight\textquotesingle{}}\NormalTok{)}
\BuiltInTok{print}\NormalTok{(}\StringTok{"Comparación guardada como \textquotesingle{}knn\_comparison.png\textquotesingle{}"}\NormalTok{)}
\end{Highlighting}
\end{Shaded}

\section{Árboles de decisión}\label{uxe1rboles-de-decisiuxf3n}

\subsection{Fundamento del algoritmo}\label{fundamento-del-algoritmo-1}

Los árboles de decisión son modelos no paramétricos que segmentan
recursivamente el espacio de características mediante reglas de decisión
simples, creando una estructura jerárquica tipo árbol. Son altamente
interpretables y capturan relaciones no lineales complejas.

\subsubsection{Estructura del árbol}\label{estructura-del-uxe1rbol}

Un árbol de decisión consiste en:

\textbf{Nodo raíz}: Contiene todos los datos

\textbf{Nodos internos}: Representan decisiones basadas en
características

\begin{itemize}
\tightlist
\item
  Test en una característica: \(X_j \leq t\)
\item
  División binaria del espacio
\end{itemize}

\textbf{Hojas (nodos terminales)}: Contienen predicciones

\begin{itemize}
\tightlist
\item
  Clasificación: Clase mayoritaria en la hoja
\item
  Probabilidad: Distribución empírica de clases
\end{itemize}

\textbf{Ramas}: Conectan nodos, representan resultados de tests

\subsubsection{Algoritmo de construcción
recursiva}\label{algoritmo-de-construcciuxf3n-recursiva}

El algoritmo CART (Classification and Regression Trees) construye el
árbol mediante particiones binarias recursivas:

\textbf{Paso 1 - Selección de división óptima}:

Para cada nodo, encontrar la mejor división \((j, t)\) que minimiza una
función de impureza:

\[
(j^*, t^*) = \arg\min_{j, t} \left[ N_L \cdot I(L) + N_R \cdot I(R) \right]
\]

donde:

\begin{itemize}
\tightlist
\item
  \(j\): Índice de característica
\item
  \(t\): Umbral de división
\item
  \(N_L, N_R\): Número de muestras en nodo izquierdo/derecho
\item
  \(I(L), I(R)\): Impureza de nodos hijo
\end{itemize}

\textbf{Paso 2 - División}:

Dividir datos según \(X_j \leq t\):

\begin{itemize}
\tightlist
\item
  Izquierda: \(\{(\mathbf{X}_i, Y_i) : X_{ij} \leq t\}\)
\item
  Derecha: \(\{(\mathbf{X}_i, Y_i) : X_{ij} > t\}\)
\end{itemize}

\textbf{Paso 3 - Recursión}:

Repetir para cada nodo hijo hasta cumplir criterio de parada.

\subsubsection{Criterios de impureza}\label{criterios-de-impureza}

\textbf{1. Índice de Gini (Gini Impurity)}:

\[
\text{Gini}(N) = 1 - \sum_{k=1}^{K} p_k^2
\]

donde \(p_k\) es la proporción de muestras de clase \(k\) en el nodo
\(N\).

\textbf{Interpretación}: Probabilidad de clasificar incorrectamente una
observación aleatoria si se etiqueta aleatoriamente según distribución
del nodo.

\textbf{Propiedades}:

\begin{itemize}
\tightlist
\item
  Mínimo: 0 (nodo puro, todas las muestras de una clase)
\item
  Máximo: \(1 - 1/K\) (distribución uniforme sobre \(K\) clases)
\item
  Para \(K=2\): Máximo en \(p=0.5\), Gini = 0.5
\end{itemize}

\textbf{2. Entropía (Entropy)}:

\[
\text{Entropy}(N) = -\sum_{k=1}^{K} p_k \log_2(p_k)
\]

(por convención, \(0 \log 0 = 0\))

\textbf{Interpretación}: Cantidad promedio de información (en bits)
necesaria para codificar la clase.

\textbf{Propiedades}:

\begin{itemize}
\tightlist
\item
  Mínimo: 0 (nodo puro)
\item
  Máximo: \(\log_2(K)\) (distribución uniforme)
\item
  Para \(K=2\): Máximo en \(p=0.5\), Entropy = 1
\end{itemize}

\textbf{Ganancia de información} (Information Gain): \[
\text{IG} = \text{Entropy}(\text{padre}) - \left[\frac{N_L}{N}\text{Entropy}(L) + \frac{N_R}{N}\text{Entropy}(R)\right]
\]

\textbf{3. Error de clasificación}:

\[
\text{Error}(N) = 1 - \max_k p_k
\]

Menos sensible que Gini y Entropy; menos usado en práctica.

\subsubsection{Comparación Gini vs
Entropy}\label{comparaciuxf3n-gini-vs-entropy}

Para clasificación binaria:

\[
\text{Gini} = 2p(1-p)
\]

\[
\text{Entropy} = -p\log_2(p) - (1-p)\log_2(1-p)
\]

\textbf{Diferencias}:

\begin{itemize}
\tightlist
\item
  Entropy penaliza más desigualdades
\item
  Gini es más rápido de calcular (no requiere logaritmo)
\item
  En práctica, producen árboles similares
\end{itemize}

\subsubsection{Criterios de parada}\label{criterios-de-parada}

El proceso de división se detiene cuando:

\begin{enumerate}
\def\labelenumi{\arabic{enumi}.}
\item
  \textbf{Pureza perfecta}: Todas las muestras en el nodo pertenecen a
  la misma clase
\item
  \textbf{Profundidad máxima} (\texttt{max\_depth}): Se alcanza
  profundidad especificada
\item
  \textbf{Muestras mínimas para dividir} (\texttt{min\_samples\_split}):
  \[
  N_{\text{nodo}} < \text{min\_samples\_split}
  \]
\item
  \textbf{Muestras mínimas en hoja} (\texttt{min\_samples\_leaf}):
  División resultaría en hoja con menos muestras que umbral
\item
  \textbf{Ganancia de información insuficiente}
  (\texttt{min\_impurity\_decrease}): \[
  \Delta \text{Impurity} < \text{threshold}
  \]
\item
  \textbf{Número de hojas máximo} (\texttt{max\_leaf\_nodes}): Se
  alcanza número especificado de hojas
\end{enumerate}

\subsubsection{Predicción}\label{predicciuxf3n}

Para clasificar una nueva observación \(\mathbf{X}_0\):

\textbf{Paso 1}: Comenzar en raíz

\textbf{Paso 2}: Seguir ramas según tests:

\begin{itemize}
\tightlist
\item
  Si \(X_{0j} \leq t\): ir a nodo izquierdo
\item
  Si \(X_{0j} > t\): ir a nodo derecho
\end{itemize}

\textbf{Paso 3}: Al llegar a hoja, predecir:

\textbf{Clase}: \[
\hat{Y}_0 = \arg\max_k p_k^{(\text{hoja})}
\]

\textbf{Probabilidad}: \[
P(Y=k|\mathbf{X}_0) = p_k^{(\text{hoja})} = \frac{N_k^{(\text{hoja})}}{N^{(\text{hoja})}}
\]

\subsubsection{Fronteras de decisión}\label{fronteras-de-decisiuxf3n-1}

Las fronteras son \textbf{hiperplanos perpendiculares a los ejes}
(axis-aligned):

\begin{itemize}
\tightlist
\item
  Cada división crea frontera perpendicular a un eje de característica
\item
  Partición recursiva del espacio en regiones rectangulares
\item
  Dentro de cada región: predicción constante
\end{itemize}

\textbf{Limitación}: No puede representar eficientemente fronteras
oblicuas (ej: \(X_1 + X_2 = 0\)).

\subsubsection{Importancia de
características}\label{importancia-de-caracteruxedsticas}

La importancia de característica \(j\) se calcula como:

\[
\text{Importance}(j) = \sum_{N : \text{usa } X_j} \Delta \text{Impurity}(N) \cdot \frac{N_{\text{samples}}(N)}{N_{\text{total}}}
\]

Suma ponderada de reducción de impureza sobre todos los nodos que usan
\(X_j\).

\textbf{Normalización}: \[
\sum_{j=1}^{p} \text{Importance}(j) = 1
\]

\subsubsection{Poda del árbol
(Pruning)}\label{poda-del-uxe1rbol-pruning}

Para controlar overfitting:

\textbf{1. Pre-pruning (early stopping)}: Detener crecimiento usando
criterios de parada estrictos

\textbf{2. Post-pruning (pruning)}: Construir árbol completo, luego
podar:

\textbf{Cost-Complexity Pruning}: \[
R_\alpha(T) = R(T) + \alpha |T|
\]

donde:

\begin{itemize}
\tightlist
\item
  \(R(T)\): Error de clasificación del árbol
\item
  \(|T|\): Número de hojas
\item
  \(\alpha \geq 0\): Parámetro de complejidad
\end{itemize}

Seleccionar \(\alpha\) que minimiza error de validación cruzada.

\subsubsection{Ventajas}\label{ventajas-2}

\begin{enumerate}
\def\labelenumi{\arabic{enumi}.}
\item
  \textbf{Interpretabilidad}: Fácil de entender y visualizar
\item
  \textbf{No requiere preprocesamiento}: No necesita estandarización ni
  normalización
\item
  \textbf{Maneja datos mixtos}: Características numéricas y categóricas
\item
  \textbf{No paramétrico}: No asume distribución de datos
\item
  \textbf{Captura no linealidades}: Automáticamente modela interacciones
  complejas
\item
  \textbf{Robustez a outliers}: En variables independientes (no en Y)
\item
  \textbf{Selección implícita de features}: Usa solo características
  informativas
\item
  \textbf{Manejo de missing values}: Algoritmos con estrategias para
  datos faltantes
\end{enumerate}

\subsubsection{Limitaciones}\label{limitaciones-2}

\begin{enumerate}
\def\labelenumi{\arabic{enumi}.}
\item
  \textbf{Overfitting}: Tendencia a ajustar demasiado datos de
  entrenamiento

  \begin{itemize}
  \tightlist
  \item
    Árboles profundos memorizan ruido
  \end{itemize}
\item
  \textbf{Inestabilidad}: Alta varianza

  \begin{itemize}
  \tightlist
  \item
    Pequeños cambios en datos → árbol completamente diferente
  \item
    No robusto
  \end{itemize}
\item
  \textbf{Sesgo con clases desbalanceadas}: Tiende a favorecer clase
  mayoritaria
\item
  \textbf{Fronteras axis-aligned}: No eficiente para fronteras oblicuas
\item
  \textbf{Greedy algorithm}: Optimización local en cada división

  \begin{itemize}
  \tightlist
  \item
    No garantiza árbol globalmente óptimo
  \end{itemize}
\item
  \textbf{Discontinuidades}: Predicciones cambian abruptamente en
  fronteras
\end{enumerate}

\subsubsection{Hiperparámetros
principales}\label{hiperparuxe1metros-principales}

{\def\LTcaptype{none} % do not increment counter
\begin{longtable}[]{@{}
  >{\raggedright\arraybackslash}p{(\linewidth - 4\tabcolsep) * \real{0.3438}}
  >{\raggedright\arraybackslash}p{(\linewidth - 4\tabcolsep) * \real{0.4062}}
  >{\raggedright\arraybackslash}p{(\linewidth - 4\tabcolsep) * \real{0.2500}}@{}}
\toprule\noalign{}
\begin{minipage}[b]{\linewidth}\raggedright
Parámetro
\end{minipage} & \begin{minipage}[b]{\linewidth}\raggedright
Descripción
\end{minipage} & \begin{minipage}[b]{\linewidth}\raggedright
Efecto
\end{minipage} \\
\midrule\noalign{}
\endhead
\bottomrule\noalign{}
\endlastfoot
\texttt{max\_depth} & Profundidad máxima & Mayor → Más complejo,
overfitting \\
\texttt{min\_samples\_split} & Muestras mínimas para dividir & Mayor →
Más simple \\
\texttt{min\_samples\_leaf} & Muestras mínimas en hoja & Mayor → Más
suave \\
\texttt{max\_features} & Features a considerar & Menor → Más
diversidad \\
\texttt{max\_leaf\_nodes} & Número máximo de hojas & Menor → Más
simple \\
\texttt{min\_impurity\_decrease} & Ganancia mínima & Mayor → Más
conservador \\
\end{longtable}
}

\subsubsection{Aplicaciones en
economía}\label{aplicaciones-en-economuxeda-1}

\begin{enumerate}
\def\labelenumi{\arabic{enumi}.}
\item
  \textbf{Credit scoring}: Reglas claras para aprobación/rechazo de
  créditos
\item
  \textbf{Segmentación de clientes}: Identificar grupos con
  características similares
\item
  \textbf{Detección de fraude}: Reglas para identificar transacciones
  sospechosas
\item
  \textbf{Predicción de churn}: Identificar clientes en riesgo de
  abandonar
\item
  \textbf{Pricing dinámico}: Reglas para ajustar precios según
  características
\item
  \textbf{Evaluación de riesgo}: Clasificación de riesgo crediticio o
  empresarial
\item
  \textbf{Decisiones de inversión}: Reglas para compra/venta de activos
\end{enumerate}

\subsubsection{Comparación con otros
métodos}\label{comparaciuxf3n-con-otros-muxe9todos}

{\def\LTcaptype{none} % do not increment counter
\begin{longtable}[]{@{}
  >{\raggedright\arraybackslash}p{(\linewidth - 6\tabcolsep) * \real{0.1765}}
  >{\raggedright\arraybackslash}p{(\linewidth - 6\tabcolsep) * \real{0.3137}}
  >{\raggedright\arraybackslash}p{(\linewidth - 6\tabcolsep) * \real{0.4118}}
  >{\raggedright\arraybackslash}p{(\linewidth - 6\tabcolsep) * \real{0.0980}}@{}}
\toprule\noalign{}
\begin{minipage}[b]{\linewidth}\raggedright
Aspecto
\end{minipage} & \begin{minipage}[b]{\linewidth}\raggedright
Decision Trees
\end{minipage} & \begin{minipage}[b]{\linewidth}\raggedright
Logistic Regression
\end{minipage} & \begin{minipage}[b]{\linewidth}\raggedright
KNN
\end{minipage} \\
\midrule\noalign{}
\endhead
\bottomrule\noalign{}
\endlastfoot
Interpretabilidad & Alta & Media & Baja \\
Captura no linealidades & Sí & No (sin transformaciones) & Sí \\
Preprocesamiento & Mínimo & Requiere estandarización & Requiere
estandarización \\
Overfitting & Alto & Bajo-Medio & Medio-Alto \\
Estabilidad & Baja & Alta & Media \\
Velocidad predicción & Rápida & Rápida & Lenta \\
\end{longtable}
}

\subsubsection{Extensiones y variantes}\label{extensiones-y-variantes-1}

\textbf{1. Árboles de regresión (CART)}: Predicción de valores continuos

\textbf{2. Random Forests}: Ensemble de árboles (siguiente sección)

\textbf{3. Gradient Boosted Trees}: Ensemble secuencial (XGBoost,
LightGBM)

\textbf{4. Conditional Inference Trees}: Usa tests estadísticos para
divisiones

\textbf{5. Árboles multivariados}: Divisiones oblicuas (no axis-aligned)

\subsection{Implementación básica}\label{implementaciuxf3n-buxe1sica}

\begin{Shaded}
\begin{Highlighting}[]
\CommentTok{\# Entrenar árbol de decisión}
\NormalTok{modelo\_tree }\OperatorTok{=}\NormalTok{ DecisionTreeClassifier(}
\NormalTok{    max\_depth}\OperatorTok{=}\DecValTok{3}\NormalTok{,}
\NormalTok{    min\_samples\_split}\OperatorTok{=}\DecValTok{10}\NormalTok{,}
\NormalTok{    min\_samples\_leaf}\OperatorTok{=}\DecValTok{5}\NormalTok{,}
\NormalTok{    random\_state}\OperatorTok{=}\DecValTok{42}
\NormalTok{)}

\NormalTok{modelo\_tree.fit(X\_train, y\_train)}

\CommentTok{\# Evaluar}
\NormalTok{y\_pred\_tree }\OperatorTok{=}\NormalTok{ modelo\_tree.predict(X\_test)}
\NormalTok{acc\_tree }\OperatorTok{=}\NormalTok{ accuracy\_score(y\_test, y\_pred\_tree)}

\BuiltInTok{print}\NormalTok{(}\StringTok{"ÁRBOL DE DECISIÓN"}\NormalTok{)}
\BuiltInTok{print}\NormalTok{(}\StringTok{"="} \OperatorTok{*} \DecValTok{70}\NormalTok{)}
\BuiltInTok{print}\NormalTok{(}\SpecialStringTok{f"Profundidad del árbol: }\SpecialCharTok{\{}\NormalTok{modelo\_tree}\SpecialCharTok{.}\NormalTok{get\_depth()}\SpecialCharTok{\}}\SpecialStringTok{"}\NormalTok{)}
\BuiltInTok{print}\NormalTok{(}\SpecialStringTok{f"Número de hojas: }\SpecialCharTok{\{}\NormalTok{modelo\_tree}\SpecialCharTok{.}\NormalTok{get\_n\_leaves()}\SpecialCharTok{\}}\SpecialStringTok{"}\NormalTok{)}
\BuiltInTok{print}\NormalTok{(}\SpecialStringTok{f"Accuracy en test: }\SpecialCharTok{\{}\NormalTok{acc\_tree}\SpecialCharTok{:.4f\}}\SpecialStringTok{ (}\SpecialCharTok{\{}\NormalTok{acc\_tree}\OperatorTok{*}\DecValTok{100}\SpecialCharTok{:.1f\}}\SpecialStringTok{\%)"}\NormalTok{)}

\CommentTok{\# Importancia de características}
\NormalTok{importancias }\OperatorTok{=}\NormalTok{ pd.DataFrame(\{}
    \StringTok{\textquotesingle{}caracteristica\textquotesingle{}}\NormalTok{: [}\StringTok{\textquotesingle{}Largo sépalo\textquotesingle{}}\NormalTok{, }\StringTok{\textquotesingle{}Ancho sépalo\textquotesingle{}}\NormalTok{],}
    \StringTok{\textquotesingle{}importancia\textquotesingle{}}\NormalTok{: modelo\_tree.feature\_importances\_}
\NormalTok{\}).sort\_values(}\StringTok{\textquotesingle{}importancia\textquotesingle{}}\NormalTok{, ascending}\OperatorTok{=}\VariableTok{False}\NormalTok{)}

\BuiltInTok{print}\NormalTok{(}\SpecialStringTok{f"}\CharTok{\textbackslash{}n}\SpecialStringTok{Importancia de características:"}\NormalTok{)}
\ControlFlowTok{for}\NormalTok{ \_, row }\KeywordTok{in}\NormalTok{ importancias.iterrows():}
    \BuiltInTok{print}\NormalTok{(}\SpecialStringTok{f"  }\SpecialCharTok{\{}\NormalTok{row[}\StringTok{\textquotesingle{}caracteristica\textquotesingle{}}\NormalTok{]}\SpecialCharTok{\}}\SpecialStringTok{: }\SpecialCharTok{\{}\NormalTok{row[}\StringTok{\textquotesingle{}importancia\textquotesingle{}}\NormalTok{]}\SpecialCharTok{:.4f\}}\SpecialStringTok{"}\NormalTok{)}
\end{Highlighting}
\end{Shaded}

\subsection{Visualización del
árbol}\label{visualizaciuxf3n-del-uxe1rbol}

\begin{Shaded}
\begin{Highlighting}[]
\CommentTok{\# Visualizar árbol}
\NormalTok{plt.figure(figsize}\OperatorTok{=}\NormalTok{(}\DecValTok{20}\NormalTok{, }\DecValTok{10}\NormalTok{))}
\NormalTok{tree.plot\_tree(}
\NormalTok{    modelo\_tree,}
\NormalTok{    feature\_names}\OperatorTok{=}\NormalTok{[}\StringTok{\textquotesingle{}Largo sépalo\textquotesingle{}}\NormalTok{, }\StringTok{\textquotesingle{}Ancho sépalo\textquotesingle{}}\NormalTok{],}
\NormalTok{    class\_names}\OperatorTok{=}\NormalTok{[}\StringTok{\textquotesingle{}No Virginica\textquotesingle{}}\NormalTok{, }\StringTok{\textquotesingle{}Virginica\textquotesingle{}}\NormalTok{],}
\NormalTok{    filled}\OperatorTok{=}\VariableTok{True}\NormalTok{,}
\NormalTok{    rounded}\OperatorTok{=}\VariableTok{True}\NormalTok{,}
\NormalTok{    fontsize}\OperatorTok{=}\DecValTok{10}
\NormalTok{)}
\NormalTok{plt.title(}\StringTok{\textquotesingle{}Árbol de Decisión {-} Clasificación de Iris\textquotesingle{}}\NormalTok{, }
\NormalTok{          fontsize}\OperatorTok{=}\DecValTok{16}\NormalTok{, fontweight}\OperatorTok{=}\StringTok{\textquotesingle{}bold\textquotesingle{}}\NormalTok{, pad}\OperatorTok{=}\DecValTok{20}\NormalTok{)}
\NormalTok{plt.tight\_layout()}
\NormalTok{plt.savefig(}\StringTok{\textquotesingle{}decision\_tree.png\textquotesingle{}}\NormalTok{, dpi}\OperatorTok{=}\DecValTok{300}\NormalTok{, bbox\_inches}\OperatorTok{=}\StringTok{\textquotesingle{}tight\textquotesingle{}}\NormalTok{)}
\BuiltInTok{print}\NormalTok{(}\StringTok{"}\CharTok{\textbackslash{}n}\StringTok{Árbol visualizado y guardado como \textquotesingle{}decision\_tree.png\textquotesingle{}"}\NormalTok{)}
\end{Highlighting}
\end{Shaded}

\subsection{Efecto de la profundidad}\label{efecto-de-la-profundidad}

\begin{Shaded}
\begin{Highlighting}[]
\CommentTok{\# Evaluar diferentes profundidades}
\NormalTok{profundidades }\OperatorTok{=} \BuiltInTok{range}\NormalTok{(}\DecValTok{1}\NormalTok{, }\DecValTok{11}\NormalTok{)}
\NormalTok{resultados\_depth }\OperatorTok{=}\NormalTok{ []}

\BuiltInTok{print}\NormalTok{(}\StringTok{"}\CharTok{\textbackslash{}n\textbackslash{}n}\StringTok{EFECTO DE LA PROFUNDIDAD DEL ÁRBOL"}\NormalTok{)}
\BuiltInTok{print}\NormalTok{(}\StringTok{"="} \OperatorTok{*} \DecValTok{70}\NormalTok{)}
\BuiltInTok{print}\NormalTok{(}\SpecialStringTok{f"}\SpecialCharTok{\{}\StringTok{\textquotesingle{}Profundidad\textquotesingle{}}\SpecialCharTok{:\textless{}12\}}\SpecialStringTok{ }\SpecialCharTok{\{}\StringTok{\textquotesingle{}Accuracy Train\textquotesingle{}}\SpecialCharTok{:\textless{}15\}}\SpecialStringTok{ }\SpecialCharTok{\{}\StringTok{\textquotesingle{}Accuracy Test\textquotesingle{}}\SpecialCharTok{:\textless{}15\}}\SpecialStringTok{ }\SpecialCharTok{\{}\StringTok{\textquotesingle{}Hojas\textquotesingle{}}\SpecialCharTok{:\textless{}10\}}\SpecialStringTok{"}\NormalTok{)}
\BuiltInTok{print}\NormalTok{(}\StringTok{"{-}"} \OperatorTok{*} \DecValTok{70}\NormalTok{)}

\ControlFlowTok{for}\NormalTok{ depth }\KeywordTok{in}\NormalTok{ profundidades:}
\NormalTok{    tree\_model }\OperatorTok{=}\NormalTok{ DecisionTreeClassifier(max\_depth}\OperatorTok{=}\NormalTok{depth, random\_state}\OperatorTok{=}\DecValTok{42}\NormalTok{)}
\NormalTok{    tree\_model.fit(X\_train, y\_train)}
    
\NormalTok{    acc\_train }\OperatorTok{=}\NormalTok{ tree\_model.score(X\_train, y\_train)}
\NormalTok{    acc\_test }\OperatorTok{=}\NormalTok{ tree\_model.score(X\_test, y\_test)}
\NormalTok{    n\_hojas }\OperatorTok{=}\NormalTok{ tree\_model.get\_n\_leaves()}
    
\NormalTok{    resultados\_depth.append(\{}
        \StringTok{\textquotesingle{}depth\textquotesingle{}}\NormalTok{: depth,}
        \StringTok{\textquotesingle{}acc\_train\textquotesingle{}}\NormalTok{: acc\_train,}
        \StringTok{\textquotesingle{}acc\_test\textquotesingle{}}\NormalTok{: acc\_test,}
        \StringTok{\textquotesingle{}n\_leaves\textquotesingle{}}\NormalTok{: n\_hojas}
\NormalTok{    \})}
    
    \BuiltInTok{print}\NormalTok{(}\SpecialStringTok{f"}\SpecialCharTok{\{}\NormalTok{depth}\SpecialCharTok{:\textless{}12\}}\SpecialStringTok{ }\SpecialCharTok{\{}\NormalTok{acc\_train}\SpecialCharTok{:\textless{}15.4f\}}\SpecialStringTok{ }\SpecialCharTok{\{}\NormalTok{acc\_test}\SpecialCharTok{:\textless{}15.4f\}}\SpecialStringTok{ }\SpecialCharTok{\{}\NormalTok{n\_hojas}\SpecialCharTok{:\textless{}10\}}\SpecialStringTok{"}\NormalTok{)}

\CommentTok{\# Graficar}
\NormalTok{df\_depth }\OperatorTok{=}\NormalTok{ pd.DataFrame(resultados\_depth)}

\NormalTok{fig, axes }\OperatorTok{=}\NormalTok{ plt.subplots(}\DecValTok{1}\NormalTok{, }\DecValTok{2}\NormalTok{, figsize}\OperatorTok{=}\NormalTok{(}\DecValTok{15}\NormalTok{, }\DecValTok{5}\NormalTok{))}

\CommentTok{\# Accuracy}
\NormalTok{axes[}\DecValTok{0}\NormalTok{].plot(df\_depth[}\StringTok{\textquotesingle{}depth\textquotesingle{}}\NormalTok{], df\_depth[}\StringTok{\textquotesingle{}acc\_train\textquotesingle{}}\NormalTok{], }
             \StringTok{\textquotesingle{}o{-}\textquotesingle{}}\NormalTok{, label}\OperatorTok{=}\StringTok{\textquotesingle{}Training\textquotesingle{}}\NormalTok{, linewidth}\OperatorTok{=}\DecValTok{2}\NormalTok{)}
\NormalTok{axes[}\DecValTok{0}\NormalTok{].plot(df\_depth[}\StringTok{\textquotesingle{}depth\textquotesingle{}}\NormalTok{], df\_depth[}\StringTok{\textquotesingle{}acc\_test\textquotesingle{}}\NormalTok{], }
             \StringTok{\textquotesingle{}s{-}\textquotesingle{}}\NormalTok{, label}\OperatorTok{=}\StringTok{\textquotesingle{}Test\textquotesingle{}}\NormalTok{, linewidth}\OperatorTok{=}\DecValTok{2}\NormalTok{)}
\NormalTok{axes[}\DecValTok{0}\NormalTok{].set\_xlabel(}\StringTok{\textquotesingle{}Profundidad máxima\textquotesingle{}}\NormalTok{)}
\NormalTok{axes[}\DecValTok{0}\NormalTok{].set\_ylabel(}\StringTok{\textquotesingle{}Accuracy\textquotesingle{}}\NormalTok{)}
\NormalTok{axes[}\DecValTok{0}\NormalTok{].set\_title(}\StringTok{\textquotesingle{}Accuracy vs Profundidad del Árbol\textquotesingle{}}\NormalTok{)}
\NormalTok{axes[}\DecValTok{0}\NormalTok{].legend()}
\NormalTok{axes[}\DecValTok{0}\NormalTok{].grid(}\VariableTok{True}\NormalTok{, alpha}\OperatorTok{=}\FloatTok{0.3}\NormalTok{)}

\CommentTok{\# Complejidad}
\NormalTok{axes[}\DecValTok{1}\NormalTok{].plot(df\_depth[}\StringTok{\textquotesingle{}depth\textquotesingle{}}\NormalTok{], df\_depth[}\StringTok{\textquotesingle{}n\_leaves\textquotesingle{}}\NormalTok{], }
             \StringTok{\textquotesingle{}o{-}\textquotesingle{}}\NormalTok{, color}\OperatorTok{=}\StringTok{\textquotesingle{}green\textquotesingle{}}\NormalTok{, linewidth}\OperatorTok{=}\DecValTok{2}\NormalTok{)}
\NormalTok{axes[}\DecValTok{1}\NormalTok{].set\_xlabel(}\StringTok{\textquotesingle{}Profundidad máxima\textquotesingle{}}\NormalTok{)}
\NormalTok{axes[}\DecValTok{1}\NormalTok{].set\_ylabel(}\StringTok{\textquotesingle{}Número de hojas\textquotesingle{}}\NormalTok{)}
\NormalTok{axes[}\DecValTok{1}\NormalTok{].set\_title(}\StringTok{\textquotesingle{}Complejidad del Árbol\textquotesingle{}}\NormalTok{)}
\NormalTok{axes[}\DecValTok{1}\NormalTok{].grid(}\VariableTok{True}\NormalTok{, alpha}\OperatorTok{=}\FloatTok{0.3}\NormalTok{)}

\NormalTok{plt.tight\_layout()}
\NormalTok{plt.savefig(}\StringTok{\textquotesingle{}tree\_depth\_analysis.png\textquotesingle{}}\NormalTok{, dpi}\OperatorTok{=}\DecValTok{300}\NormalTok{, bbox\_inches}\OperatorTok{=}\StringTok{\textquotesingle{}tight\textquotesingle{}}\NormalTok{)}
\BuiltInTok{print}\NormalTok{(}\StringTok{"}\CharTok{\textbackslash{}n}\StringTok{Análisis guardado como \textquotesingle{}tree\_depth\_analysis.png\textquotesingle{}"}\NormalTok{)}
\end{Highlighting}
\end{Shaded}

\section{Random Forests}\label{random-forests}

\subsection{Fundamento del algoritmo}\label{fundamento-del-algoritmo-2}

Random Forests (Bosques Aleatorios) es un método de ensemble que combina
múltiples árboles de decisión para producir un clasificador más robusto
y preciso. Fue propuesto por Leo Breiman en 2001 y se basa en el
principio de ``la sabiduría de las multitudes''.

\subsubsection{Concepto central: Ensemble
Learning}\label{concepto-central-ensemble-learning}

\textbf{Idea fundamental}: Combinar predicciones de múltiples modelos
débiles para crear un modelo fuerte.

\textbf{Analogía}: Consultar a múltiples expertos y promediar sus
opiniones suele dar mejor resultado que confiar en uno solo.

\subsubsection{Algoritmo Random Forest}\label{algoritmo-random-forest}

\textbf{Entrada}:

\begin{itemize}
\tightlist
\item
  Dataset \(\mathcal{D} = \{(\mathbf{X}_i, Y_i)\}_{i=1}^{n}\)
\item
  Número de árboles \(B\)
\item
  Número de características por división \(m\)
\end{itemize}

\textbf{Proceso}:

Para \(b = 1, \ldots, B\):

\textbf{Paso 1 - Bootstrap}: Crear muestra aleatoria con reemplazo \[
\mathcal{D}_b = \{\text{muestra } n \text{ observaciones de } \mathcal{D} \text{ con reemplazo}\}
\]

Aproximadamente 63.2\% de observaciones únicas en cada muestra (algunas
repetidas, otras excluidas).

\textbf{Paso 2 - Entrenar árbol con randomización adicional}:

Para cada nodo del árbol \(b\):

\begin{itemize}
\tightlist
\item
  Seleccionar aleatoriamente \(m\) características de las \(p\)
  disponibles
\item
  Encontrar mejor división usando solo estas \(m\) características
\item
  Dividir nodo
\item
  Repetir recursivamente
\end{itemize}

\textbf{Paso 3 - Crecer árbol completo}: Sin poda (o con poda mínima).
Cada árbol tiene alta varianza pero bajo sesgo.

\textbf{Predicción}:

Para nueva observación \(\mathbf{X}_0\):

\textbf{Clasificación} (votación mayoritaria): \[
\hat{Y}_0 = \text{moda}\{\hat{Y}_0^{(1)}, \hat{Y}_0^{(2)}, \ldots, \hat{Y}_0^{(B)}\}
\]

\textbf{Probabilidad}: \[
\hat{P}(Y=k|\mathbf{X}_0) = \frac{1}{B} \sum_{b=1}^{B} \mathbb{1}(\hat{Y}_0^{(b)} = k)
\]

\subsubsection{Dos fuentes de
aleatoriedad}\label{dos-fuentes-de-aleatoriedad}

Random Forests introduce dos niveles de aleatoriedad:

\textbf{1. Bootstrap (Bagging)}:

\begin{itemize}
\tightlist
\item
  Diferentes muestras de entrenamiento para cada árbol
\item
  Reduce varianza al promediar predictores correlacionados
\end{itemize}

\textbf{2. Random Feature Selection}:

\begin{itemize}
\tightlist
\item
  Diferentes subsets de características en cada división
\item
  \textbf{Decorrelaciona} los árboles
\item
  Típicamente \(m = \sqrt{p}\) para clasificación
\end{itemize}

Sin (2), los árboles serían muy similares (correlacionados), y el
promedio no reduciría tanto la varianza.

\subsubsection{Relación con Bagging}\label{relaciuxf3n-con-bagging}

\textbf{Bagging} (Bootstrap Aggregating): Caso especial de Random Forest
donde \(m = p\)

Random Forest = Bagging + Random Feature Selection

La selección aleatoria de features es la innovación clave que mejora
Bagging.

\subsubsection{Out-of-Bag (OOB) Error}\label{out-of-bag-oob-error}

Para cada árbol \(b\), aproximadamente 37\% de observaciones no están en
\(\mathcal{D}_b\).

Estas observaciones \textbf{Out-of-Bag} pueden usarse como conjunto de
validación interno:

\[
\text{OOB Error} = \frac{1}{n} \sum_{i=1}^{n} \mathbb{1}\left(Y_i \neq \frac{1}{|B_i|}\sum_{b \in B_i} \hat{Y}_i^{(b)}\right)
\]

donde \(B_i\) es el conjunto de árboles que no usaron observación \(i\).

\textbf{Ventaja}: No requiere conjunto de validación separado.

\subsubsection{Importancia de
características}\label{importancia-de-caracteruxedsticas-1}

Random Forests calcula dos medidas de importancia:

\textbf{1. Mean Decrease Impurity (MDI)}:

\[
\text{Importance}(X_j) = \frac{1}{B} \sum_{b=1}^{B} \sum_{N \in T_b : \text{usa } X_j} \Delta I(N) \cdot \frac{n_N}{n}
\]

Promedio sobre todos los árboles de la reducción de impureza cuando se
usa \(X_j\).

\textbf{2. Mean Decrease Accuracy (MDA)} (Permutation Importance):

Para cada característica \(X_j\):

\begin{enumerate}
\def\labelenumi{\arabic{enumi}.}
\tightlist
\item
  Calcular OOB error original
\item
  Permutar aleatoriamente valores de \(X_j\) en OOB
\item
  Recalcular OOB error
\item
  Importancia = Incremento en error
\end{enumerate}

\[
\text{Importance}(X_j) = \text{OOB Error}_{\text{permutado}} - \text{OOB Error}_{\text{original}}
\]

\textbf{Interpretación}: Si \(X_j\) es importante, permutarla aumentará
significativamente el error.

\subsubsection{Propiedades teóricas}\label{propiedades-teuxf3ricas-1}

\textbf{Teorema (Breiman)}:

Bajo ciertas condiciones, cuando \(B \to \infty\):

\[
\lim_{B \to \infty} \text{Error}_{RF} = \rho \cdot \bar{\sigma}^2 \cdot (1 - s^2)
\]

donde:

\begin{itemize}
\tightlist
\item
  \(\rho\): Correlación promedio entre árboles
\item
  \(\bar{\sigma}^2\): Varianza promedio de cada árbol
\item
  \(s\): Fuerza promedio de cada árbol
\end{itemize}

\textbf{Implicaciones}:

\begin{enumerate}
\def\labelenumi{\arabic{enumi}.}
\tightlist
\item
  \textbf{Aumentar B}: Siempre ayuda (no overfitting)
\item
  \textbf{Reducir ρ} (decorrelación): Ayuda → Random feature selection
\item
  \textbf{Aumentar fuerza árboles}: Ayuda → Usar más características
\end{enumerate}

\subsubsection{Hiperparámetros
principales}\label{hiperparuxe1metros-principales-1}

{\def\LTcaptype{none} % do not increment counter
\begin{longtable}[]{@{}
  >{\raggedright\arraybackslash}p{(\linewidth - 6\tabcolsep) * \real{0.2391}}
  >{\raggedright\arraybackslash}p{(\linewidth - 6\tabcolsep) * \real{0.2826}}
  >{\raggedright\arraybackslash}p{(\linewidth - 6\tabcolsep) * \real{0.3043}}
  >{\raggedright\arraybackslash}p{(\linewidth - 6\tabcolsep) * \real{0.1739}}@{}}
\toprule\noalign{}
\begin{minipage}[b]{\linewidth}\raggedright
Parámetro
\end{minipage} & \begin{minipage}[b]{\linewidth}\raggedright
Descripción
\end{minipage} & \begin{minipage}[b]{\linewidth}\raggedright
Valor típico
\end{minipage} & \begin{minipage}[b]{\linewidth}\raggedright
Efecto
\end{minipage} \\
\midrule\noalign{}
\endhead
\bottomrule\noalign{}
\endlastfoot
\texttt{n\_estimators} & Número de árboles & 100-500 & Más = mejor (con
límite) \\
\texttt{max\_features} & Features por división & \(\sqrt{p}\)
(clasificación) & Controla correlación \\
\texttt{max\_depth} & Profundidad máxima & None (completo) & Controla
sesgo-varianza \\
\texttt{min\_samples\_split} & Mínimo para dividir & 2 & Mayor = más
simple \\
\texttt{min\_samples\_leaf} & Mínimo en hoja & 1 & Mayor = más suave \\
\texttt{bootstrap} & Usar bootstrap & True & False = usar todo
dataset \\
\end{longtable}
}

\subsubsection{Ventajas}\label{ventajas-3}

\begin{enumerate}
\def\labelenumi{\arabic{enumi}.}
\item
  \textbf{Reduce overfitting}: Vs árbol único, por promediado
\item
  \textbf{Alta precisión}: Entre los mejores algoritmos out-of-the-box
\item
  \textbf{Robusto}: Menos sensible a hiperparámetros que otros métodos
\item
  \textbf{Maneja alta dimensionalidad}: Funciona bien incluso con \(p\)
  grande
\item
  \textbf{Importancia de features}: Proporciona ranking de
  características
\item
  \textbf{No requiere estandarización}: Como árboles individuales
\item
  \textbf{Paralelizable}: Árboles se entrenan independientemente
\item
  \textbf{OOB error}: Validación interna sin conjunto separado
\item
  \textbf{Maneja missing values}: Estrategias incorporadas
\item
  \textbf{Versátil}: Funciona bien en diversos problemas
\end{enumerate}

\subsubsection{Limitaciones}\label{limitaciones-3}

\begin{enumerate}
\def\labelenumi{\arabic{enumi}.}
\item
  \textbf{Menos interpretable}: Difícil entender modelo (caja negra)
\item
  \textbf{Costo computacional}:

  \begin{itemize}
  \tightlist
  \item
    Entrenamiento: \(O(B \cdot n \log n \cdot m)\)
  \item
    Predicción: \(O(B \cdot \log n)\) Más lento que un árbol único
  \end{itemize}
\item
  \textbf{Memoria}: Requiere almacenar \(B\) árboles
\item
  \textbf{Extrapolación limitada}: No predice fuera del rango de datos
  de entrenamiento
\item
  \textbf{Sesgo con clases desbalanceadas}: Puede favorecer clase
  mayoritaria
\item
  \textbf{No captura tendencias lineales simples}: A veces un modelo
  simple es mejor
\end{enumerate}

\subsubsection{Comparación con otros
métodos}\label{comparaciuxf3n-con-otros-muxe9todos-1}

{\def\LTcaptype{none} % do not increment counter
\begin{longtable}[]{@{}llll@{}}
\toprule\noalign{}
Aspecto & Random Forest & Decision Tree & Gradient Boosting \\
\midrule\noalign{}
\endhead
\bottomrule\noalign{}
\endlastfoot
Interpretabilidad & Baja & Alta & Baja \\
Accuracy & Alta & Media & Muy alta \\
Overfitting & Bajo & Alto & Medio (con tuning) \\
Velocidad entrenamiento & Media & Rápida & Lenta \\
Velocidad predicción & Media & Rápida & Media \\
Hiperparámetros & Pocos críticos & Pocos & Muchos críticos \\
Robustez & Alta & Baja & Media \\
\end{longtable}
}

\subsubsection{Selección de
hiperparámetros}\label{selecciuxf3n-de-hiperparuxe1metros}

\textbf{Estrategia típica}:

\begin{enumerate}
\def\labelenumi{\arabic{enumi}.}
\item
  \textbf{n\_estimators}: Comenzar con 100, aumentar hasta que
  performance se estabilice
\item
  \textbf{max\_features}:

  \begin{itemize}
  \tightlist
  \item
    Clasificación: \(\sqrt{p}\)
  \item
    Regresión: \(p/3\)
  \item
    Probar también \(\log_2(p)\) y \(p/2\)
  \end{itemize}
\item
  \textbf{max\_depth}: Comenzar sin límite, reducir si overfitting
\item
  \textbf{min\_samples\_split} y \textbf{min\_samples\_leaf}: Usar
  valores por defecto inicialmente
\end{enumerate}

\textbf{Grid Search} sobre:

\begin{itemize}
\tightlist
\item
  \texttt{n\_estimators}: {[}50, 100, 200, 500{]}
\item
  \texttt{max\_features}: {[}`sqrt', `log2', 0.3{]}
\item
  \texttt{max\_depth}: {[}None, 10, 20, 30{]}
\item
  \texttt{min\_samples\_split}: {[}2, 5, 10{]}
\end{itemize}

\subsubsection{Aplicaciones en
economía}\label{aplicaciones-en-economuxeda-2}

\begin{enumerate}
\def\labelenumi{\arabic{enumi}.}
\item
  \textbf{Credit scoring}: Predicción de incumplimiento con alta
  precisión
\item
  \textbf{Predicción de precios}:

  \begin{itemize}
  \tightlist
  \item
    Inmuebles
  \item
    Acciones
  \item
    Commodities
  \end{itemize}
\item
  \textbf{Detección de fraude}: Identificar transacciones fraudulentas
\item
  \textbf{Churn prediction}: Predecir abandono de clientes
\item
  \textbf{Segmentación avanzada}: Clasificar clientes en múltiples
  categorías
\item
  \textbf{Predicción de quiebra empresarial}: Clasificación de riesgo
\item
  \textbf{Trading algorítmico}: Señales de compra/venta
\item
  \textbf{Risk management}: Evaluación de riesgo de portafolios
\end{enumerate}

\subsubsection{Variantes y extensiones}\label{variantes-y-extensiones}

\textbf{1. Extremely Randomized Trees (Extra Trees)}:

\begin{itemize}
\tightlist
\item
  Divisiones completamente aleatorias (no optimizadas)
\item
  Más rápido, más diversidad
\end{itemize}

\textbf{2. Balanced Random Forest}:

\begin{itemize}
\tightlist
\item
  Balanceo de clases en cada árbol
\item
  Mejor para datasets desbalanceados
\end{itemize}

\textbf{3. Isolation Forest}:

\begin{itemize}
\tightlist
\item
  Para detección de anomalías
\item
  Aísla outliers
\end{itemize}

\textbf{4. Quantile Regression Forest}:

\begin{itemize}
\tightlist
\item
  Predice distribución completa, no solo media
\item
  Útil para intervalos de confianza
\end{itemize}

\subsection{Implementación}\label{implementaciuxf3n-2}

\begin{Shaded}
\begin{Highlighting}[]
\CommentTok{\# Entrenar Random Forest}
\NormalTok{modelo\_rf }\OperatorTok{=}\NormalTok{ RandomForestClassifier(}
\NormalTok{    n\_estimators}\OperatorTok{=}\DecValTok{100}\NormalTok{,}
\NormalTok{    max\_depth}\OperatorTok{=}\DecValTok{5}\NormalTok{,}
\NormalTok{    min\_samples\_split}\OperatorTok{=}\DecValTok{10}\NormalTok{,}
\NormalTok{    min\_samples\_leaf}\OperatorTok{=}\DecValTok{5}\NormalTok{,}
\NormalTok{    random\_state}\OperatorTok{=}\DecValTok{42}\NormalTok{,}
\NormalTok{    n\_jobs}\OperatorTok{={-}}\DecValTok{1}  \CommentTok{\# Usar todos los procesadores}
\NormalTok{)}

\NormalTok{modelo\_rf.fit(X\_train, y\_train)}

\CommentTok{\# Evaluar}
\NormalTok{y\_pred\_rf }\OperatorTok{=}\NormalTok{ modelo\_rf.predict(X\_test)}
\NormalTok{acc\_rf }\OperatorTok{=}\NormalTok{ accuracy\_score(y\_test, y\_pred\_rf)}

\BuiltInTok{print}\NormalTok{(}\StringTok{"RANDOM FOREST"}\NormalTok{)}
\BuiltInTok{print}\NormalTok{(}\StringTok{"="} \OperatorTok{*} \DecValTok{70}\NormalTok{)}
\BuiltInTok{print}\NormalTok{(}\SpecialStringTok{f"Número de árboles: }\SpecialCharTok{\{}\NormalTok{modelo\_rf}\SpecialCharTok{.}\NormalTok{n\_estimators}\SpecialCharTok{\}}\SpecialStringTok{"}\NormalTok{)}
\BuiltInTok{print}\NormalTok{(}\SpecialStringTok{f"Accuracy en test: }\SpecialCharTok{\{}\NormalTok{acc\_rf}\SpecialCharTok{:.4f\}}\SpecialStringTok{ (}\SpecialCharTok{\{}\NormalTok{acc\_rf}\OperatorTok{*}\DecValTok{100}\SpecialCharTok{:.1f\}}\SpecialStringTok{\%)"}\NormalTok{)}

\CommentTok{\# Importancia de características}
\NormalTok{importancias\_rf }\OperatorTok{=}\NormalTok{ pd.DataFrame(\{}
    \StringTok{\textquotesingle{}caracteristica\textquotesingle{}}\NormalTok{: [}\StringTok{\textquotesingle{}Largo sépalo\textquotesingle{}}\NormalTok{, }\StringTok{\textquotesingle{}Ancho sépalo\textquotesingle{}}\NormalTok{],}
    \StringTok{\textquotesingle{}importancia\textquotesingle{}}\NormalTok{: modelo\_rf.feature\_importances\_,}
    \StringTok{\textquotesingle{}std\textquotesingle{}}\NormalTok{: np.std([tree.feature\_importances\_ }
                   \ControlFlowTok{for}\NormalTok{ tree }\KeywordTok{in}\NormalTok{ modelo\_rf.estimators\_], axis}\OperatorTok{=}\DecValTok{0}\NormalTok{)}
\NormalTok{\}).sort\_values(}\StringTok{\textquotesingle{}importancia\textquotesingle{}}\NormalTok{, ascending}\OperatorTok{=}\VariableTok{False}\NormalTok{)}

\BuiltInTok{print}\NormalTok{(}\SpecialStringTok{f"}\CharTok{\textbackslash{}n}\SpecialStringTok{Importancia de características:"}\NormalTok{)}
\ControlFlowTok{for}\NormalTok{ \_, row }\KeywordTok{in}\NormalTok{ importancias\_rf.iterrows():}
    \BuiltInTok{print}\NormalTok{(}\SpecialStringTok{f"  }\SpecialCharTok{\{}\NormalTok{row[}\StringTok{\textquotesingle{}caracteristica\textquotesingle{}}\NormalTok{]}\SpecialCharTok{\}}\SpecialStringTok{: }\SpecialCharTok{\{}\NormalTok{row[}\StringTok{\textquotesingle{}importancia\textquotesingle{}}\NormalTok{]}\SpecialCharTok{:.4f\}}\SpecialStringTok{ (±}\SpecialCharTok{\{}\NormalTok{row[}\StringTok{\textquotesingle{}std\textquotesingle{}}\NormalTok{]}\SpecialCharTok{:.4f\}}\SpecialStringTok{)"}\NormalTok{)}
\end{Highlighting}
\end{Shaded}

\subsection{Efecto del número de
árboles}\label{efecto-del-nuxfamero-de-uxe1rboles}

\begin{Shaded}
\begin{Highlighting}[]
\CommentTok{\# Evaluar diferentes números de árboles}
\NormalTok{n\_trees\_list }\OperatorTok{=}\NormalTok{ [}\DecValTok{1}\NormalTok{, }\DecValTok{5}\NormalTok{, }\DecValTok{10}\NormalTok{, }\DecValTok{20}\NormalTok{, }\DecValTok{50}\NormalTok{, }\DecValTok{100}\NormalTok{, }\DecValTok{200}\NormalTok{]}
\NormalTok{resultados\_trees }\OperatorTok{=}\NormalTok{ []}

\BuiltInTok{print}\NormalTok{(}\StringTok{"}\CharTok{\textbackslash{}n\textbackslash{}n}\StringTok{EFECTO DEL NÚMERO DE ÁRBOLES"}\NormalTok{)}
\BuiltInTok{print}\NormalTok{(}\StringTok{"="} \OperatorTok{*} \DecValTok{70}\NormalTok{)}
\BuiltInTok{print}\NormalTok{(}\SpecialStringTok{f"}\SpecialCharTok{\{}\StringTok{\textquotesingle{}N Árboles\textquotesingle{}}\SpecialCharTok{:\textless{}12\}}\SpecialStringTok{ }\SpecialCharTok{\{}\StringTok{\textquotesingle{}Accuracy Train\textquotesingle{}}\SpecialCharTok{:\textless{}15\}}\SpecialStringTok{ }\SpecialCharTok{\{}\StringTok{\textquotesingle{}Accuracy Test\textquotesingle{}}\SpecialCharTok{:\textless{}15\}}\SpecialStringTok{"}\NormalTok{)}
\BuiltInTok{print}\NormalTok{(}\StringTok{"{-}"} \OperatorTok{*} \DecValTok{70}\NormalTok{)}

\ControlFlowTok{for}\NormalTok{ n\_trees }\KeywordTok{in}\NormalTok{ n\_trees\_list:}
\NormalTok{    rf }\OperatorTok{=}\NormalTok{ RandomForestClassifier(}
\NormalTok{        n\_estimators}\OperatorTok{=}\NormalTok{n\_trees,}
\NormalTok{        max\_depth}\OperatorTok{=}\DecValTok{5}\NormalTok{,}
\NormalTok{        random\_state}\OperatorTok{=}\DecValTok{42}
\NormalTok{    )}
\NormalTok{    rf.fit(X\_train, y\_train)}
    
\NormalTok{    acc\_train }\OperatorTok{=}\NormalTok{ rf.score(X\_train, y\_train)}
\NormalTok{    acc\_test }\OperatorTok{=}\NormalTok{ rf.score(X\_test, y\_test)}
    
\NormalTok{    resultados\_trees.append(\{}
        \StringTok{\textquotesingle{}n\_trees\textquotesingle{}}\NormalTok{: n\_trees,}
        \StringTok{\textquotesingle{}acc\_train\textquotesingle{}}\NormalTok{: acc\_train,}
        \StringTok{\textquotesingle{}acc\_test\textquotesingle{}}\NormalTok{: acc\_test}
\NormalTok{    \})}
    
    \BuiltInTok{print}\NormalTok{(}\SpecialStringTok{f"}\SpecialCharTok{\{}\NormalTok{n\_trees}\SpecialCharTok{:\textless{}12\}}\SpecialStringTok{ }\SpecialCharTok{\{}\NormalTok{acc\_train}\SpecialCharTok{:\textless{}15.4f\}}\SpecialStringTok{ }\SpecialCharTok{\{}\NormalTok{acc\_test}\SpecialCharTok{:\textless{}15.4f\}}\SpecialStringTok{"}\NormalTok{)}

\CommentTok{\# Graficar}
\NormalTok{df\_trees }\OperatorTok{=}\NormalTok{ pd.DataFrame(resultados\_trees)}

\NormalTok{plt.figure(figsize}\OperatorTok{=}\NormalTok{(}\DecValTok{10}\NormalTok{, }\DecValTok{6}\NormalTok{))}
\NormalTok{plt.plot(df\_trees[}\StringTok{\textquotesingle{}n\_trees\textquotesingle{}}\NormalTok{], df\_trees[}\StringTok{\textquotesingle{}acc\_train\textquotesingle{}}\NormalTok{], }
         \StringTok{\textquotesingle{}o{-}\textquotesingle{}}\NormalTok{, label}\OperatorTok{=}\StringTok{\textquotesingle{}Training\textquotesingle{}}\NormalTok{, linewidth}\OperatorTok{=}\DecValTok{2}\NormalTok{, markersize}\OperatorTok{=}\DecValTok{8}\NormalTok{)}
\NormalTok{plt.plot(df\_trees[}\StringTok{\textquotesingle{}n\_trees\textquotesingle{}}\NormalTok{], df\_trees[}\StringTok{\textquotesingle{}acc\_test\textquotesingle{}}\NormalTok{], }
         \StringTok{\textquotesingle{}s{-}\textquotesingle{}}\NormalTok{, label}\OperatorTok{=}\StringTok{\textquotesingle{}Test\textquotesingle{}}\NormalTok{, linewidth}\OperatorTok{=}\DecValTok{2}\NormalTok{, markersize}\OperatorTok{=}\DecValTok{8}\NormalTok{)}
\NormalTok{plt.xlabel(}\StringTok{\textquotesingle{}Número de árboles\textquotesingle{}}\NormalTok{)}
\NormalTok{plt.ylabel(}\StringTok{\textquotesingle{}Accuracy\textquotesingle{}}\NormalTok{)}
\NormalTok{plt.title(}\StringTok{\textquotesingle{}Accuracy vs Número de Árboles en Random Forest\textquotesingle{}}\NormalTok{)}
\NormalTok{plt.legend()}
\NormalTok{plt.grid(}\VariableTok{True}\NormalTok{, alpha}\OperatorTok{=}\FloatTok{0.3}\NormalTok{)}
\NormalTok{plt.xscale(}\StringTok{\textquotesingle{}log\textquotesingle{}}\NormalTok{)}
\NormalTok{plt.tight\_layout()}
\NormalTok{plt.savefig(}\StringTok{\textquotesingle{}random\_forest\_ntrees.png\textquotesingle{}}\NormalTok{, dpi}\OperatorTok{=}\DecValTok{300}\NormalTok{, bbox\_inches}\OperatorTok{=}\StringTok{\textquotesingle{}tight\textquotesingle{}}\NormalTok{)}
\BuiltInTok{print}\NormalTok{(}\StringTok{"}\CharTok{\textbackslash{}n}\StringTok{Gráfico guardado como \textquotesingle{}random\_forest\_ntrees.png\textquotesingle{}"}\NormalTok{)}
\end{Highlighting}
\end{Shaded}

\section{Clasificación multi-clase}\label{clasificaciuxf3n-multi-clase}

\subsection{Dataset con 3 clases}\label{dataset-con-3-clases}

\begin{Shaded}
\begin{Highlighting}[]
\CommentTok{\# Usar todas las especies de Iris (3 clases)}
\NormalTok{X\_multiclass }\OperatorTok{=}\NormalTok{ iris.data[:, :}\DecValTok{2}\NormalTok{]}
\NormalTok{y\_multiclass }\OperatorTok{=}\NormalTok{ iris.target}

\BuiltInTok{print}\NormalTok{(}\StringTok{"CLASIFICACIÓN MULTI{-}CLASE"}\NormalTok{)}
\BuiltInTok{print}\NormalTok{(}\StringTok{"="} \OperatorTok{*} \DecValTok{70}\NormalTok{)}
\BuiltInTok{print}\NormalTok{(}\SpecialStringTok{f"Número de clases: }\SpecialCharTok{\{}\BuiltInTok{len}\NormalTok{(np.unique(y\_multiclass))}\SpecialCharTok{\}}\SpecialStringTok{"}\NormalTok{)}
\BuiltInTok{print}\NormalTok{(}\SpecialStringTok{f"Clases: }\SpecialCharTok{\{}\NormalTok{iris}\SpecialCharTok{.}\NormalTok{target\_names}\SpecialCharTok{\}}\SpecialStringTok{"}\NormalTok{)}
\BuiltInTok{print}\NormalTok{(}\SpecialStringTok{f"}\CharTok{\textbackslash{}n}\SpecialStringTok{Distribución:"}\NormalTok{)}
\ControlFlowTok{for}\NormalTok{ i, nombre }\KeywordTok{in} \BuiltInTok{enumerate}\NormalTok{(iris.target\_names):}
\NormalTok{    count }\OperatorTok{=}\NormalTok{ (y\_multiclass }\OperatorTok{==}\NormalTok{ i).}\BuiltInTok{sum}\NormalTok{()}
    \BuiltInTok{print}\NormalTok{(}\SpecialStringTok{f"  }\SpecialCharTok{\{}\NormalTok{nombre}\SpecialCharTok{\}}\SpecialStringTok{: }\SpecialCharTok{\{}\NormalTok{count}\SpecialCharTok{\}}\SpecialStringTok{ muestras"}\NormalTok{)}

\CommentTok{\# Dividir datos}
\NormalTok{X\_train\_mc, X\_test\_mc, y\_train\_mc, y\_test\_mc }\OperatorTok{=}\NormalTok{ train\_test\_split(}
\NormalTok{    X\_multiclass, y\_multiclass, test\_size}\OperatorTok{=}\FloatTok{0.3}\NormalTok{, random\_state}\OperatorTok{=}\DecValTok{42}
\NormalTok{)}
\end{Highlighting}
\end{Shaded}

\subsection{Entrenar modelos
multi-clase}\label{entrenar-modelos-multi-clase}

\begin{Shaded}
\begin{Highlighting}[]
\CommentTok{\# Entrenar cada modelo}
\NormalTok{modelos\_mc }\OperatorTok{=}\NormalTok{ \{}
    \StringTok{\textquotesingle{}Logistic Regression\textquotesingle{}}\NormalTok{: LogisticRegression(random\_state}\OperatorTok{=}\DecValTok{42}\NormalTok{, max\_iter}\OperatorTok{=}\DecValTok{200}\NormalTok{),}
    \StringTok{\textquotesingle{}KNN\textquotesingle{}}\NormalTok{: KNeighborsClassifier(n\_neighbors}\OperatorTok{=}\DecValTok{5}\NormalTok{),}
    \StringTok{\textquotesingle{}Decision Tree\textquotesingle{}}\NormalTok{: DecisionTreeClassifier(max\_depth}\OperatorTok{=}\DecValTok{5}\NormalTok{, random\_state}\OperatorTok{=}\DecValTok{42}\NormalTok{),}
    \StringTok{\textquotesingle{}Random Forest\textquotesingle{}}\NormalTok{: RandomForestClassifier(n\_estimators}\OperatorTok{=}\DecValTok{100}\NormalTok{, random\_state}\OperatorTok{=}\DecValTok{42}\NormalTok{)}
\NormalTok{\}}

\NormalTok{resultados\_mc }\OperatorTok{=}\NormalTok{ \{\}}

\BuiltInTok{print}\NormalTok{(}\StringTok{"}\CharTok{\textbackslash{}n\textbackslash{}n}\StringTok{RESULTADOS CLASIFICACIÓN MULTI{-}CLASE"}\NormalTok{)}
\BuiltInTok{print}\NormalTok{(}\StringTok{"="} \OperatorTok{*} \DecValTok{70}\NormalTok{)}

\ControlFlowTok{for}\NormalTok{ nombre, modelo }\KeywordTok{in}\NormalTok{ modelos\_mc.items():}
    \CommentTok{\# Entrenar}
\NormalTok{    modelo.fit(X\_train\_mc, y\_train\_mc)}
    
    \CommentTok{\# Predecir}
\NormalTok{    y\_pred }\OperatorTok{=}\NormalTok{ modelo.predict(X\_test\_mc)}
    
    \CommentTok{\# Evaluar}
\NormalTok{    acc }\OperatorTok{=}\NormalTok{ accuracy\_score(y\_test\_mc, y\_pred)}
    
\NormalTok{    resultados\_mc[nombre] }\OperatorTok{=}\NormalTok{ \{}
        \StringTok{\textquotesingle{}modelo\textquotesingle{}}\NormalTok{: modelo,}
        \StringTok{\textquotesingle{}accuracy\textquotesingle{}}\NormalTok{: acc}
\NormalTok{    \}}
    
    \BuiltInTok{print}\NormalTok{(}\SpecialStringTok{f"}\CharTok{\textbackslash{}n}\SpecialCharTok{\{}\NormalTok{nombre}\SpecialCharTok{\}}\SpecialStringTok{:"}\NormalTok{)}
    \BuiltInTok{print}\NormalTok{(}\SpecialStringTok{f"  Accuracy: }\SpecialCharTok{\{}\NormalTok{acc}\SpecialCharTok{:.4f\}}\SpecialStringTok{ (}\SpecialCharTok{\{}\NormalTok{acc}\OperatorTok{*}\DecValTok{100}\SpecialCharTok{:.1f\}}\SpecialStringTok{\%)"}\NormalTok{)}
    
    \CommentTok{\# Matriz de confusión}
\NormalTok{    cm }\OperatorTok{=}\NormalTok{ confusion\_matrix(y\_test\_mc, y\_pred)}
    \BuiltInTok{print}\NormalTok{(}\SpecialStringTok{f"  Matriz de confusión:"}\NormalTok{)}
    \BuiltInTok{print}\NormalTok{(}\SpecialStringTok{f"  }\SpecialCharTok{\{}\NormalTok{cm}\SpecialCharTok{\}}\SpecialStringTok{"}\NormalTok{)}

\CommentTok{\# Encontrar mejor modelo}
\NormalTok{mejor\_modelo\_mc }\OperatorTok{=} \BuiltInTok{max}\NormalTok{(resultados\_mc, key}\OperatorTok{=}\KeywordTok{lambda}\NormalTok{ x: resultados\_mc[x][}\StringTok{\textquotesingle{}accuracy\textquotesingle{}}\NormalTok{])}
\BuiltInTok{print}\NormalTok{(}\SpecialStringTok{f"}\CharTok{\textbackslash{}n\textbackslash{}n}\SpecialStringTok{Mejor modelo: }\SpecialCharTok{\{}\NormalTok{mejor\_modelo\_mc}\SpecialCharTok{\}}\SpecialStringTok{"}\NormalTok{)}
\BuiltInTok{print}\NormalTok{(}\SpecialStringTok{f"Accuracy: }\SpecialCharTok{\{}\NormalTok{resultados\_mc[mejor\_modelo\_mc][}\StringTok{\textquotesingle{}accuracy\textquotesingle{}}\NormalTok{]}\SpecialCharTok{:.4f\}}\SpecialStringTok{"}\NormalTok{)}
\end{Highlighting}
\end{Shaded}

\subsection{Visualización
multi-clase}\label{visualizaciuxf3n-multi-clase}

\begin{Shaded}
\begin{Highlighting}[]
\CommentTok{\# Visualizar fronteras de decisión para cada modelo}
\NormalTok{fig, axes }\OperatorTok{=}\NormalTok{ plt.subplots(}\DecValTok{2}\NormalTok{, }\DecValTok{2}\NormalTok{, figsize}\OperatorTok{=}\NormalTok{(}\DecValTok{16}\NormalTok{, }\DecValTok{14}\NormalTok{))}
\NormalTok{axes }\OperatorTok{=}\NormalTok{ axes.ravel()}

\NormalTok{h }\OperatorTok{=} \FloatTok{0.02}
\NormalTok{x\_min, x\_max }\OperatorTok{=}\NormalTok{ X\_multiclass[:, }\DecValTok{0}\NormalTok{].}\BuiltInTok{min}\NormalTok{() }\OperatorTok{{-}} \FloatTok{0.5}\NormalTok{, X\_multiclass[:, }\DecValTok{0}\NormalTok{].}\BuiltInTok{max}\NormalTok{() }\OperatorTok{+} \FloatTok{0.5}
\NormalTok{y\_min, y\_max }\OperatorTok{=}\NormalTok{ X\_multiclass[:, }\DecValTok{1}\NormalTok{].}\BuiltInTok{min}\NormalTok{() }\OperatorTok{{-}} \FloatTok{0.5}\NormalTok{, X\_multiclass[:, }\DecValTok{1}\NormalTok{].}\BuiltInTok{max}\NormalTok{() }\OperatorTok{+} \FloatTok{0.5}
\NormalTok{xx, yy }\OperatorTok{=}\NormalTok{ np.meshgrid(np.arange(x\_min, x\_max, h),}
\NormalTok{                     np.arange(y\_min, y\_max, h))}

\ControlFlowTok{for}\NormalTok{ idx, (nombre, resultado) }\KeywordTok{in} \BuiltInTok{enumerate}\NormalTok{(resultados\_mc.items()):}
\NormalTok{    modelo }\OperatorTok{=}\NormalTok{ resultado[}\StringTok{\textquotesingle{}modelo\textquotesingle{}}\NormalTok{]}
\NormalTok{    acc }\OperatorTok{=}\NormalTok{ resultado[}\StringTok{\textquotesingle{}accuracy\textquotesingle{}}\NormalTok{]}
    
    \CommentTok{\# Predecir para malla}
\NormalTok{    Z }\OperatorTok{=}\NormalTok{ modelo.predict(np.c\_[xx.ravel(), yy.ravel()])}
\NormalTok{    Z }\OperatorTok{=}\NormalTok{ Z.reshape(xx.shape)}
    
    \CommentTok{\# Graficar}
\NormalTok{    axes[idx].contourf(xx, yy, Z, alpha}\OperatorTok{=}\FloatTok{0.4}\NormalTok{, cmap}\OperatorTok{=}\StringTok{\textquotesingle{}viridis\textquotesingle{}}\NormalTok{)}
\NormalTok{    scatter }\OperatorTok{=}\NormalTok{ axes[idx].scatter(X\_multiclass[:, }\DecValTok{0}\NormalTok{], X\_multiclass[:, }\DecValTok{1}\NormalTok{],}
\NormalTok{                               c}\OperatorTok{=}\NormalTok{y\_multiclass, cmap}\OperatorTok{=}\StringTok{\textquotesingle{}viridis\textquotesingle{}}\NormalTok{,}
\NormalTok{                               edgecolors}\OperatorTok{=}\StringTok{\textquotesingle{}black\textquotesingle{}}\NormalTok{, s}\OperatorTok{=}\DecValTok{50}\NormalTok{)}
    
\NormalTok{    axes[idx].set\_xlabel(}\StringTok{\textquotesingle{}Largo del sépalo (cm)\textquotesingle{}}\NormalTok{)}
\NormalTok{    axes[idx].set\_ylabel(}\StringTok{\textquotesingle{}Ancho del sépalo (cm)\textquotesingle{}}\NormalTok{)}
\NormalTok{    axes[idx].set\_title(}\SpecialStringTok{f\textquotesingle{}}\SpecialCharTok{\{}\NormalTok{nombre}\SpecialCharTok{\}}\CharTok{\textbackslash{}n}\SpecialStringTok{Accuracy=}\SpecialCharTok{\{}\NormalTok{acc}\SpecialCharTok{:.3f\}}\SpecialStringTok{\textquotesingle{}}\NormalTok{)}
\NormalTok{    axes[idx].grid(}\VariableTok{True}\NormalTok{, alpha}\OperatorTok{=}\FloatTok{0.3}\NormalTok{)}

\CommentTok{\# Leyenda común}
\NormalTok{plt.colorbar(scatter, ax}\OperatorTok{=}\NormalTok{axes, label}\OperatorTok{=}\StringTok{\textquotesingle{}Especie\textquotesingle{}}\NormalTok{, ticks}\OperatorTok{=}\NormalTok{[}\DecValTok{0}\NormalTok{, }\DecValTok{1}\NormalTok{, }\DecValTok{2}\NormalTok{])}

\NormalTok{plt.tight\_layout()}
\NormalTok{plt.savefig(}\StringTok{\textquotesingle{}multiclass\_comparison.png\textquotesingle{}}\NormalTok{, dpi}\OperatorTok{=}\DecValTok{300}\NormalTok{, bbox\_inches}\OperatorTok{=}\StringTok{\textquotesingle{}tight\textquotesingle{}}\NormalTok{)}
\BuiltInTok{print}\NormalTok{(}\StringTok{"}\CharTok{\textbackslash{}n}\StringTok{Comparación guardada como \textquotesingle{}multiclass\_comparison.png\textquotesingle{}"}\NormalTok{)}
\end{Highlighting}
\end{Shaded}

\section{Comparación de modelos}\label{comparaciuxf3n-de-modelos}

\subsection{Curvas ROC para cada
modelo}\label{curvas-roc-para-cada-modelo}

\begin{Shaded}
\begin{Highlighting}[]
\CommentTok{\# Calcular curvas ROC para clasificación binaria}
\KeywordTok{def}\NormalTok{ calcular\_roc\_multimodelo(modelos\_dict, X\_test, y\_test):}
    \CommentTok{"""Calcula curvas ROC para múltiples modelos"""}
    
\NormalTok{    plt.figure(figsize}\OperatorTok{=}\NormalTok{(}\DecValTok{10}\NormalTok{, }\DecValTok{8}\NormalTok{))}
    
    \ControlFlowTok{for}\NormalTok{ nombre, modelo }\KeywordTok{in}\NormalTok{ modelos\_dict.items():}
        \CommentTok{\# Obtener probabilidades}
        \ControlFlowTok{if} \BuiltInTok{hasattr}\NormalTok{(modelo, }\StringTok{\textquotesingle{}predict\_proba\textquotesingle{}}\NormalTok{):}
\NormalTok{            probs }\OperatorTok{=}\NormalTok{ modelo.predict\_proba(X\_test)[:, }\DecValTok{1}\NormalTok{]}
        \ControlFlowTok{else}\NormalTok{:}
\NormalTok{            probs }\OperatorTok{=}\NormalTok{ modelo.decision\_function(X\_test)}
        
        \CommentTok{\# Calcular curva ROC}
\NormalTok{        fpr, tpr, \_ }\OperatorTok{=}\NormalTok{ roc\_curve(y\_test, probs)}
\NormalTok{        roc\_auc }\OperatorTok{=}\NormalTok{ auc(fpr, tpr)}
        
        \CommentTok{\# Graficar}
\NormalTok{        plt.plot(fpr, tpr, linewidth}\OperatorTok{=}\DecValTok{2}\NormalTok{, }
\NormalTok{                label}\OperatorTok{=}\SpecialStringTok{f\textquotesingle{}}\SpecialCharTok{\{}\NormalTok{nombre}\SpecialCharTok{\}}\SpecialStringTok{ (AUC=}\SpecialCharTok{\{}\NormalTok{roc\_auc}\SpecialCharTok{:.3f\}}\SpecialStringTok{)\textquotesingle{}}\NormalTok{)}
    
    \CommentTok{\# Línea diagonal}
\NormalTok{    plt.plot([}\DecValTok{0}\NormalTok{, }\DecValTok{1}\NormalTok{], [}\DecValTok{0}\NormalTok{, }\DecValTok{1}\NormalTok{], }\StringTok{\textquotesingle{}k{-}{-}\textquotesingle{}}\NormalTok{, linewidth}\OperatorTok{=}\DecValTok{2}\NormalTok{, label}\OperatorTok{=}\StringTok{\textquotesingle{}Azar (AUC=0.5)\textquotesingle{}}\NormalTok{)}
    
\NormalTok{    plt.xlabel(}\StringTok{\textquotesingle{}False Positive Rate\textquotesingle{}}\NormalTok{, fontsize}\OperatorTok{=}\DecValTok{12}\NormalTok{)}
\NormalTok{    plt.ylabel(}\StringTok{\textquotesingle{}True Positive Rate\textquotesingle{}}\NormalTok{, fontsize}\OperatorTok{=}\DecValTok{12}\NormalTok{)}
\NormalTok{    plt.title(}\StringTok{\textquotesingle{}Curvas ROC {-} Comparación de Modelos\textquotesingle{}}\NormalTok{, }
\NormalTok{              fontsize}\OperatorTok{=}\DecValTok{14}\NormalTok{, fontweight}\OperatorTok{=}\StringTok{\textquotesingle{}bold\textquotesingle{}}\NormalTok{)}
\NormalTok{    plt.legend(loc}\OperatorTok{=}\StringTok{\textquotesingle{}lower right\textquotesingle{}}\NormalTok{)}
\NormalTok{    plt.grid(}\VariableTok{True}\NormalTok{, alpha}\OperatorTok{=}\FloatTok{0.3}\NormalTok{)}
\NormalTok{    plt.tight\_layout()}
\NormalTok{    plt.savefig(}\StringTok{\textquotesingle{}roc\_comparison.png\textquotesingle{}}\NormalTok{, dpi}\OperatorTok{=}\DecValTok{300}\NormalTok{, bbox\_inches}\OperatorTok{=}\StringTok{\textquotesingle{}tight\textquotesingle{}}\NormalTok{)}
    \BuiltInTok{print}\NormalTok{(}\StringTok{"Curvas ROC guardadas como \textquotesingle{}roc\_comparison.png\textquotesingle{}"}\NormalTok{)}

\CommentTok{\# Preparar modelos para comparación (clasificación binaria)}
\NormalTok{modelos\_binarios }\OperatorTok{=}\NormalTok{ \{}
    \StringTok{\textquotesingle{}Logistic Regression\textquotesingle{}}\NormalTok{: LogisticRegression(random\_state}\OperatorTok{=}\DecValTok{42}\NormalTok{),}
    \StringTok{\textquotesingle{}KNN (k=5)\textquotesingle{}}\NormalTok{: KNeighborsClassifier(n\_neighbors}\OperatorTok{=}\DecValTok{5}\NormalTok{),}
    \StringTok{\textquotesingle{}Decision Tree\textquotesingle{}}\NormalTok{: DecisionTreeClassifier(max\_depth}\OperatorTok{=}\DecValTok{5}\NormalTok{, random\_state}\OperatorTok{=}\DecValTok{42}\NormalTok{),}
    \StringTok{\textquotesingle{}Random Forest\textquotesingle{}}\NormalTok{: RandomForestClassifier(n\_estimators}\OperatorTok{=}\DecValTok{100}\NormalTok{, random\_state}\OperatorTok{=}\DecValTok{42}\NormalTok{)}
\NormalTok{\}}

\CommentTok{\# Entrenar modelos}
\ControlFlowTok{for}\NormalTok{ modelo }\KeywordTok{in}\NormalTok{ modelos\_binarios.values():}
\NormalTok{    modelo.fit(X\_train, y\_train)}

\CommentTok{\# Calcular y graficar ROC}
\NormalTok{calcular\_roc\_multimodelo(modelos\_binarios, X\_test, y\_test)}
\end{Highlighting}
\end{Shaded}

\subsection{Tabla comparativa de
métricas}\label{tabla-comparativa-de-muxe9tricas}

\begin{Shaded}
\begin{Highlighting}[]
\CommentTok{\# Compilar todas las métricas}
\BuiltInTok{print}\NormalTok{(}\StringTok{"}\CharTok{\textbackslash{}n\textbackslash{}n}\StringTok{TABLA COMPARATIVA DE MODELOS"}\NormalTok{)}
\BuiltInTok{print}\NormalTok{(}\StringTok{"="} \OperatorTok{*} \DecValTok{70}\NormalTok{)}

\NormalTok{resultados\_comparacion }\OperatorTok{=}\NormalTok{ []}

\ControlFlowTok{for}\NormalTok{ nombre, modelo }\KeywordTok{in}\NormalTok{ modelos\_binarios.items():}
    \CommentTok{\# Predicciones}
\NormalTok{    y\_pred }\OperatorTok{=}\NormalTok{ modelo.predict(X\_test)}
    
    \CommentTok{\# Probabilidades}
    \ControlFlowTok{if} \BuiltInTok{hasattr}\NormalTok{(modelo, }\StringTok{\textquotesingle{}predict\_proba\textquotesingle{}}\NormalTok{):}
\NormalTok{        probs }\OperatorTok{=}\NormalTok{ modelo.predict\_proba(X\_test)[:, }\DecValTok{1}\NormalTok{]}
    \ControlFlowTok{else}\NormalTok{:}
\NormalTok{        probs }\OperatorTok{=}\NormalTok{ modelo.decision\_function(X\_test)}
    
    \CommentTok{\# Métricas}
    \ImportTok{from}\NormalTok{ sklearn.metrics }\ImportTok{import}\NormalTok{ precision\_score, recall\_score, f1\_score}
    
\NormalTok{    acc }\OperatorTok{=}\NormalTok{ accuracy\_score(y\_test, y\_pred)}
\NormalTok{    precision }\OperatorTok{=}\NormalTok{ precision\_score(y\_test, y\_pred)}
\NormalTok{    recall }\OperatorTok{=}\NormalTok{ recall\_score(y\_test, y\_pred)}
\NormalTok{    f1 }\OperatorTok{=}\NormalTok{ f1\_score(y\_test, y\_pred)}
    
    \CommentTok{\# ROC AUC}
\NormalTok{    fpr, tpr, \_ }\OperatorTok{=}\NormalTok{ roc\_curve(y\_test, probs)}
\NormalTok{    roc\_auc }\OperatorTok{=}\NormalTok{ auc(fpr, tpr)}
    
\NormalTok{    resultados\_comparacion.append(\{}
        \StringTok{\textquotesingle{}Modelo\textquotesingle{}}\NormalTok{: nombre,}
        \StringTok{\textquotesingle{}Accuracy\textquotesingle{}}\NormalTok{: acc,}
        \StringTok{\textquotesingle{}Precision\textquotesingle{}}\NormalTok{: precision,}
        \StringTok{\textquotesingle{}Recall\textquotesingle{}}\NormalTok{: recall,}
        \StringTok{\textquotesingle{}F1{-}Score\textquotesingle{}}\NormalTok{: f1,}
        \StringTok{\textquotesingle{}AUC\textquotesingle{}}\NormalTok{: roc\_auc}
\NormalTok{    \})}

\CommentTok{\# Crear DataFrame}
\NormalTok{df\_comparacion }\OperatorTok{=}\NormalTok{ pd.DataFrame(resultados\_comparacion)}
\NormalTok{df\_comparacion }\OperatorTok{=}\NormalTok{ df\_comparacion.}\BuiltInTok{round}\NormalTok{(}\DecValTok{4}\NormalTok{)}

\BuiltInTok{print}\NormalTok{(df\_comparacion.to\_string(index}\OperatorTok{=}\VariableTok{False}\NormalTok{))}

\CommentTok{\# Guardar}
\NormalTok{df\_comparacion.to\_csv(}\StringTok{\textquotesingle{}comparacion\_modelos.csv\textquotesingle{}}\NormalTok{, index}\OperatorTok{=}\VariableTok{False}\NormalTok{)}
\BuiltInTok{print}\NormalTok{(}\StringTok{"}\CharTok{\textbackslash{}n\textbackslash{}n}\StringTok{Tabla guardada como \textquotesingle{}comparacion\_modelos.csv\textquotesingle{}"}\NormalTok{)}
\end{Highlighting}
\end{Shaded}

\section{Aplicaciones económicas}\label{aplicaciones-econuxf3micas-1}

\subsection{Aplicación 1: Clasificación de riesgo
crediticio}\label{aplicaciuxf3n-1-clasificaciuxf3n-de-riesgo-crediticio}

\begin{Shaded}
\begin{Highlighting}[]
\BuiltInTok{print}\NormalTok{(}\StringTok{"}\CharTok{\textbackslash{}n\textbackslash{}n\textbackslash{}n}\StringTok{APLICACIÓN: CLASIFICACIÓN DE RIESGO CREDITICIO"}\NormalTok{)}
\BuiltInTok{print}\NormalTok{(}\StringTok{"="} \OperatorTok{*} \DecValTok{70}\NormalTok{)}

\CommentTok{\# Generar datos sintéticos}
\NormalTok{np.random.seed(}\DecValTok{42}\NormalTok{)}
\NormalTok{n\_clientes }\OperatorTok{=} \DecValTok{2000}

\NormalTok{datos\_credito }\OperatorTok{=}\NormalTok{ pd.DataFrame(\{}
    \StringTok{\textquotesingle{}edad\textquotesingle{}}\NormalTok{: np.random.randint(}\DecValTok{18}\NormalTok{, }\DecValTok{70}\NormalTok{, n\_clientes),}
    \StringTok{\textquotesingle{}ingreso\_anual\textquotesingle{}}\NormalTok{: np.random.uniform(}\DecValTok{10000}\NormalTok{, }\DecValTok{150000}\NormalTok{, n\_clientes),}
    \StringTok{\textquotesingle{}monto\_solicitado\textquotesingle{}}\NormalTok{: np.random.uniform(}\DecValTok{5000}\NormalTok{, }\DecValTok{200000}\NormalTok{, n\_clientes),}
    \StringTok{\textquotesingle{}antiguedad\_laboral\textquotesingle{}}\NormalTok{: np.random.randint(}\DecValTok{0}\NormalTok{, }\DecValTok{30}\NormalTok{, n\_clientes),}
    \StringTok{\textquotesingle{}num\_creditos\_previos\textquotesingle{}}\NormalTok{: np.random.randint(}\DecValTok{0}\NormalTok{, }\DecValTok{10}\NormalTok{, n\_clientes),}
    \StringTok{\textquotesingle{}ratio\_deuda\_ingreso\textquotesingle{}}\NormalTok{: np.random.uniform(}\DecValTok{0}\NormalTok{, }\FloatTok{1.5}\NormalTok{, n\_clientes),}
    \StringTok{\textquotesingle{}tiene\_hipoteca\textquotesingle{}}\NormalTok{: np.random.choice([}\DecValTok{0}\NormalTok{, }\DecValTok{1}\NormalTok{], n\_clientes),}
    \StringTok{\textquotesingle{}score\_crediticio\textquotesingle{}}\NormalTok{: np.random.randint(}\DecValTok{300}\NormalTok{, }\DecValTok{850}\NormalTok{, n\_clientes)}
\NormalTok{\})}

\CommentTok{\# Variable objetivo: riesgo (0=bajo, 1=medio, 2=alto)}
\NormalTok{score\_riesgo }\OperatorTok{=}\NormalTok{ (}
    \OperatorTok{{-}}\NormalTok{datos\_credito[}\StringTok{\textquotesingle{}score\_crediticio\textquotesingle{}}\NormalTok{] }\OperatorTok{*} \FloatTok{0.003} \OperatorTok{+}
\NormalTok{    datos\_credito[}\StringTok{\textquotesingle{}ratio\_deuda\_ingreso\textquotesingle{}}\NormalTok{] }\OperatorTok{*} \DecValTok{2} \OperatorTok{+}
    \OperatorTok{{-}}\NormalTok{datos\_credito[}\StringTok{\textquotesingle{}antiguedad\_laboral\textquotesingle{}}\NormalTok{] }\OperatorTok{*} \FloatTok{0.05} \OperatorTok{+}
\NormalTok{    datos\_credito[}\StringTok{\textquotesingle{}num\_creditos\_previos\textquotesingle{}}\NormalTok{] }\OperatorTok{*} \FloatTok{0.2} \OperatorTok{+}
\NormalTok{    np.random.normal(}\DecValTok{0}\NormalTok{, }\FloatTok{0.5}\NormalTok{, n\_clientes)}
\NormalTok{)}

\CommentTok{\# Clasificar en 3 niveles}
\NormalTok{datos\_credito[}\StringTok{\textquotesingle{}riesgo\textquotesingle{}}\NormalTok{] }\OperatorTok{=}\NormalTok{ pd.cut(}
\NormalTok{    score\_riesgo,}
\NormalTok{    bins}\OperatorTok{=}\DecValTok{3}\NormalTok{,}
\NormalTok{    labels}\OperatorTok{=}\NormalTok{[}\StringTok{\textquotesingle{}Bajo\textquotesingle{}}\NormalTok{, }\StringTok{\textquotesingle{}Medio\textquotesingle{}}\NormalTok{, }\StringTok{\textquotesingle{}Alto\textquotesingle{}}\NormalTok{]}
\NormalTok{)}

\CommentTok{\# Convertir a numérico}
\NormalTok{datos\_credito[}\StringTok{\textquotesingle{}riesgo\_num\textquotesingle{}}\NormalTok{] }\OperatorTok{=}\NormalTok{ datos\_credito[}\StringTok{\textquotesingle{}riesgo\textquotesingle{}}\NormalTok{].}\BuiltInTok{map}\NormalTok{(\{}
    \StringTok{\textquotesingle{}Bajo\textquotesingle{}}\NormalTok{: }\DecValTok{0}\NormalTok{,}
    \StringTok{\textquotesingle{}Medio\textquotesingle{}}\NormalTok{: }\DecValTok{1}\NormalTok{,}
    \StringTok{\textquotesingle{}Alto\textquotesingle{}}\NormalTok{: }\DecValTok{2}
\NormalTok{\})}

\BuiltInTok{print}\NormalTok{(}\SpecialStringTok{f"}\CharTok{\textbackslash{}n}\SpecialStringTok{Dataset de crédito:"}\NormalTok{)}
\BuiltInTok{print}\NormalTok{(}\SpecialStringTok{f"  Total clientes: }\SpecialCharTok{\{}\NormalTok{n\_clientes}\SpecialCharTok{:,\}}\SpecialStringTok{"}\NormalTok{)}
\BuiltInTok{print}\NormalTok{(}\SpecialStringTok{f"}\CharTok{\textbackslash{}n}\SpecialStringTok{Distribución de riesgo:"}\NormalTok{)}
\BuiltInTok{print}\NormalTok{(datos\_credito[}\StringTok{\textquotesingle{}riesgo\textquotesingle{}}\NormalTok{].value\_counts())}

\CommentTok{\# Preparar datos}
\NormalTok{X\_credito }\OperatorTok{=}\NormalTok{ datos\_credito.drop([}\StringTok{\textquotesingle{}riesgo\textquotesingle{}}\NormalTok{, }\StringTok{\textquotesingle{}riesgo\_num\textquotesingle{}}\NormalTok{], axis}\OperatorTok{=}\DecValTok{1}\NormalTok{)}
\NormalTok{y\_credito }\OperatorTok{=}\NormalTok{ datos\_credito[}\StringTok{\textquotesingle{}riesgo\_num\textquotesingle{}}\NormalTok{]}

\NormalTok{X\_train\_cr, X\_test\_cr, y\_train\_cr, y\_test\_cr }\OperatorTok{=}\NormalTok{ train\_test\_split(}
\NormalTok{    X\_credito, y\_credito, test\_size}\OperatorTok{=}\FloatTok{0.3}\NormalTok{, random\_state}\OperatorTok{=}\DecValTok{42}\NormalTok{, stratify}\OperatorTok{=}\NormalTok{y\_credito}
\NormalTok{)}

\CommentTok{\# Estandarizar}
\NormalTok{scaler\_credito }\OperatorTok{=}\NormalTok{ StandardScaler()}
\NormalTok{X\_train\_cr\_scaled }\OperatorTok{=}\NormalTok{ scaler\_credito.fit\_transform(X\_train\_cr)}
\NormalTok{X\_test\_cr\_scaled }\OperatorTok{=}\NormalTok{ scaler\_credito.transform(X\_test\_cr)}

\CommentTok{\# Entrenar modelos}
\BuiltInTok{print}\NormalTok{(}\StringTok{"}\CharTok{\textbackslash{}n\textbackslash{}n}\StringTok{ENTRENAMIENTO DE MODELOS"}\NormalTok{)}
\BuiltInTok{print}\NormalTok{(}\StringTok{"="} \OperatorTok{*} \DecValTok{70}\NormalTok{)}

\NormalTok{modelos\_credito }\OperatorTok{=}\NormalTok{ \{}
    \StringTok{\textquotesingle{}Random Forest\textquotesingle{}}\NormalTok{: RandomForestClassifier(n\_estimators}\OperatorTok{=}\DecValTok{200}\NormalTok{, max\_depth}\OperatorTok{=}\DecValTok{10}\NormalTok{, random\_state}\OperatorTok{=}\DecValTok{42}\NormalTok{),}
    \StringTok{\textquotesingle{}Decision Tree\textquotesingle{}}\NormalTok{: DecisionTreeClassifier(max\_depth}\OperatorTok{=}\DecValTok{8}\NormalTok{, random\_state}\OperatorTok{=}\DecValTok{42}\NormalTok{),}
    \StringTok{\textquotesingle{}KNN\textquotesingle{}}\NormalTok{: KNeighborsClassifier(n\_neighbors}\OperatorTok{=}\DecValTok{10}\NormalTok{),}
    \StringTok{\textquotesingle{}Logistic Regression\textquotesingle{}}\NormalTok{: LogisticRegression(max\_iter}\OperatorTok{=}\DecValTok{1000}\NormalTok{, random\_state}\OperatorTok{=}\DecValTok{42}\NormalTok{)}
\NormalTok{\}}

\ControlFlowTok{for}\NormalTok{ nombre, modelo }\KeywordTok{in}\NormalTok{ modelos\_credito.items():}
    \CommentTok{\# Entrenar}
    \ControlFlowTok{if}\NormalTok{ nombre }\KeywordTok{in}\NormalTok{ [}\StringTok{\textquotesingle{}KNN\textquotesingle{}}\NormalTok{, }\StringTok{\textquotesingle{}Logistic Regression\textquotesingle{}}\NormalTok{]:}
\NormalTok{        modelo.fit(X\_train\_cr\_scaled, y\_train\_cr)}
\NormalTok{        y\_pred }\OperatorTok{=}\NormalTok{ modelo.predict(X\_test\_cr\_scaled)}
    \ControlFlowTok{else}\NormalTok{:}
\NormalTok{        modelo.fit(X\_train\_cr, y\_train\_cr)}
\NormalTok{        y\_pred }\OperatorTok{=}\NormalTok{ modelo.predict(X\_test\_cr)}
    
    \CommentTok{\# Evaluar}
\NormalTok{    acc }\OperatorTok{=}\NormalTok{ accuracy\_score(y\_test\_cr, y\_pred)}
    
    \BuiltInTok{print}\NormalTok{(}\SpecialStringTok{f"}\CharTok{\textbackslash{}n}\SpecialCharTok{\{}\NormalTok{nombre}\SpecialCharTok{\}}\SpecialStringTok{:"}\NormalTok{)}
    \BuiltInTok{print}\NormalTok{(}\SpecialStringTok{f"  Accuracy: }\SpecialCharTok{\{}\NormalTok{acc}\SpecialCharTok{:.4f\}}\SpecialStringTok{ (}\SpecialCharTok{\{}\NormalTok{acc}\OperatorTok{*}\DecValTok{100}\SpecialCharTok{:.1f\}}\SpecialStringTok{\%)"}\NormalTok{)}
    
    \CommentTok{\# Reporte por clase}
    \BuiltInTok{print}\NormalTok{(}\SpecialStringTok{f"}\CharTok{\textbackslash{}n}\SpecialStringTok{  Reporte de clasificación:"}\NormalTok{)}
    \BuiltInTok{print}\NormalTok{(classification\_report(}
\NormalTok{        y\_test\_cr, y\_pred,}
\NormalTok{        target\_names}\OperatorTok{=}\NormalTok{[}\StringTok{\textquotesingle{}Bajo\textquotesingle{}}\NormalTok{, }\StringTok{\textquotesingle{}Medio\textquotesingle{}}\NormalTok{, }\StringTok{\textquotesingle{}Alto\textquotesingle{}}\NormalTok{],}
\NormalTok{        digits}\OperatorTok{=}\DecValTok{3}
\NormalTok{    ))}

\CommentTok{\# Importancia de características (Random Forest)}
\BuiltInTok{print}\NormalTok{(}\StringTok{"}\CharTok{\textbackslash{}n\textbackslash{}n}\StringTok{IMPORTANCIA DE CARACTERÍSTICAS (Random Forest)"}\NormalTok{)}
\BuiltInTok{print}\NormalTok{(}\StringTok{"="} \OperatorTok{*} \DecValTok{70}\NormalTok{)}

\NormalTok{rf\_final }\OperatorTok{=}\NormalTok{ modelos\_credito[}\StringTok{\textquotesingle{}Random Forest\textquotesingle{}}\NormalTok{]}
\NormalTok{importancias }\OperatorTok{=}\NormalTok{ pd.DataFrame(\{}
    \StringTok{\textquotesingle{}caracteristica\textquotesingle{}}\NormalTok{: X\_credito.columns,}
    \StringTok{\textquotesingle{}importancia\textquotesingle{}}\NormalTok{: rf\_final.feature\_importances\_}
\NormalTok{\}).sort\_values(}\StringTok{\textquotesingle{}importancia\textquotesingle{}}\NormalTok{, ascending}\OperatorTok{=}\VariableTok{False}\NormalTok{)}

\ControlFlowTok{for}\NormalTok{ \_, row }\KeywordTok{in}\NormalTok{ importancias.iterrows():}
    \BuiltInTok{print}\NormalTok{(}\SpecialStringTok{f"  }\SpecialCharTok{\{}\NormalTok{row[}\StringTok{\textquotesingle{}caracteristica\textquotesingle{}}\NormalTok{]}\SpecialCharTok{:\textless{}25\}}\SpecialStringTok{ }\SpecialCharTok{\{}\NormalTok{row[}\StringTok{\textquotesingle{}importancia\textquotesingle{}}\NormalTok{]}\SpecialCharTok{:.4f\}}\SpecialStringTok{ }\SpecialCharTok{\{}\StringTok{\textquotesingle{}█\textquotesingle{}} \OperatorTok{*} \BuiltInTok{int}\NormalTok{(row[}\StringTok{\textquotesingle{}importancia\textquotesingle{}}\NormalTok{] }\OperatorTok{*} \DecValTok{100}\NormalTok{)}\SpecialCharTok{\}}\SpecialStringTok{"}\NormalTok{)}
\end{Highlighting}
\end{Shaded}

\subsection{Aplicación 2: Segmentación de
clientes}\label{aplicaciuxf3n-2-segmentaciuxf3n-de-clientes}

\begin{Shaded}
\begin{Highlighting}[]
\BuiltInTok{print}\NormalTok{(}\StringTok{"}\CharTok{\textbackslash{}n\textbackslash{}n\textbackslash{}n}\StringTok{APLICACIÓN: SEGMENTACIÓN DE CLIENTES"}\NormalTok{)}
\BuiltInTok{print}\NormalTok{(}\StringTok{"="} \OperatorTok{*} \DecValTok{70}\NormalTok{)}

\CommentTok{\# Generar datos de clientes}
\NormalTok{np.random.seed(}\DecValTok{42}\NormalTok{)}
\NormalTok{n\_clientes\_seg }\OperatorTok{=} \DecValTok{1500}

\NormalTok{datos\_clientes }\OperatorTok{=}\NormalTok{ pd.DataFrame(\{}
    \StringTok{\textquotesingle{}edad\textquotesingle{}}\NormalTok{: np.random.randint(}\DecValTok{18}\NormalTok{, }\DecValTok{80}\NormalTok{, n\_clientes\_seg),}
    \StringTok{\textquotesingle{}ingreso\_mensual\textquotesingle{}}\NormalTok{: np.random.uniform(}\DecValTok{1000}\NormalTok{, }\DecValTok{20000}\NormalTok{, n\_clientes\_seg),}
    \StringTok{\textquotesingle{}gasto\_mensual\textquotesingle{}}\NormalTok{: np.random.uniform(}\DecValTok{500}\NormalTok{, }\DecValTok{15000}\NormalTok{, n\_clientes\_seg),}
    \StringTok{\textquotesingle{}num\_compras\_mes\textquotesingle{}}\NormalTok{: np.random.randint(}\DecValTok{0}\NormalTok{, }\DecValTok{30}\NormalTok{, n\_clientes\_seg),}
    \StringTok{\textquotesingle{}antiguedad\_meses\textquotesingle{}}\NormalTok{: np.random.randint(}\DecValTok{1}\NormalTok{, }\DecValTok{120}\NormalTok{, n\_clientes\_seg),}
    \StringTok{\textquotesingle{}ticket\_promedio\textquotesingle{}}\NormalTok{: np.random.uniform(}\DecValTok{20}\NormalTok{, }\DecValTok{500}\NormalTok{, n\_clientes\_seg),}
    \StringTok{\textquotesingle{}visitas\_web\textquotesingle{}}\NormalTok{: np.random.randint(}\DecValTok{0}\NormalTok{, }\DecValTok{100}\NormalTok{, n\_clientes\_seg)}
\NormalTok{\})}

\CommentTok{\# Calcular características adicionales}
\NormalTok{datos\_clientes[}\StringTok{\textquotesingle{}frecuencia\_compra\textquotesingle{}}\NormalTok{] }\OperatorTok{=}\NormalTok{ datos\_clientes[}\StringTok{\textquotesingle{}num\_compras\_mes\textquotesingle{}}\NormalTok{] }\OperatorTok{/} \DecValTok{4}
\NormalTok{datos\_clientes[}\StringTok{\textquotesingle{}ratio\_gasto\_ingreso\textquotesingle{}}\NormalTok{] }\OperatorTok{=}\NormalTok{ datos\_clientes[}\StringTok{\textquotesingle{}gasto\_mensual\textquotesingle{}}\NormalTok{] }\OperatorTok{/}\NormalTok{ datos\_clientes[}\StringTok{\textquotesingle{}ingreso\_mensual\textquotesingle{}}\NormalTok{]}

\CommentTok{\# Crear segmentos (0=Bronce, 1=Plata, 2=Oro, 3=Platino)}
\NormalTok{score\_valor }\OperatorTok{=}\NormalTok{ (}
\NormalTok{    datos\_clientes[}\StringTok{\textquotesingle{}gasto\_mensual\textquotesingle{}}\NormalTok{] }\OperatorTok{*} \FloatTok{0.0003} \OperatorTok{+}
\NormalTok{    datos\_clientes[}\StringTok{\textquotesingle{}num\_compras\_mes\textquotesingle{}}\NormalTok{] }\OperatorTok{*} \FloatTok{0.1} \OperatorTok{+}
\NormalTok{    datos\_clientes[}\StringTok{\textquotesingle{}ticket\_promedio\textquotesingle{}}\NormalTok{] }\OperatorTok{*} \FloatTok{0.005} \OperatorTok{+}
\NormalTok{    np.random.normal(}\DecValTok{0}\NormalTok{, }\FloatTok{0.3}\NormalTok{, n\_clientes\_seg)}
\NormalTok{)}

\NormalTok{datos\_clientes[}\StringTok{\textquotesingle{}segmento\textquotesingle{}}\NormalTok{] }\OperatorTok{=}\NormalTok{ pd.cut(}
\NormalTok{    score\_valor,}
\NormalTok{    bins}\OperatorTok{=}\DecValTok{4}\NormalTok{,}
\NormalTok{    labels}\OperatorTok{=}\NormalTok{[}\StringTok{\textquotesingle{}Bronce\textquotesingle{}}\NormalTok{, }\StringTok{\textquotesingle{}Plata\textquotesingle{}}\NormalTok{, }\StringTok{\textquotesingle{}Oro\textquotesingle{}}\NormalTok{, }\StringTok{\textquotesingle{}Platino\textquotesingle{}}\NormalTok{]}
\NormalTok{)}

\NormalTok{datos\_clientes[}\StringTok{\textquotesingle{}segmento\_num\textquotesingle{}}\NormalTok{] }\OperatorTok{=}\NormalTok{ datos\_clientes[}\StringTok{\textquotesingle{}segmento\textquotesingle{}}\NormalTok{].}\BuiltInTok{map}\NormalTok{(\{}
    \StringTok{\textquotesingle{}Bronce\textquotesingle{}}\NormalTok{: }\DecValTok{0}\NormalTok{,}
    \StringTok{\textquotesingle{}Plata\textquotesingle{}}\NormalTok{: }\DecValTok{1}\NormalTok{,}
    \StringTok{\textquotesingle{}Oro\textquotesingle{}}\NormalTok{: }\DecValTok{2}\NormalTok{,}
    \StringTok{\textquotesingle{}Platino\textquotesingle{}}\NormalTok{: }\DecValTok{3}
\NormalTok{\})}

\BuiltInTok{print}\NormalTok{(}\SpecialStringTok{f"}\CharTok{\textbackslash{}n}\SpecialStringTok{Distribución de segmentos:"}\NormalTok{)}
\BuiltInTok{print}\NormalTok{(datos\_clientes[}\StringTok{\textquotesingle{}segmento\textquotesingle{}}\NormalTok{].value\_counts().sort\_index())}

\CommentTok{\# Preparar datos}
\NormalTok{features\_seg }\OperatorTok{=}\NormalTok{ [}\StringTok{\textquotesingle{}edad\textquotesingle{}}\NormalTok{, }\StringTok{\textquotesingle{}ingreso\_mensual\textquotesingle{}}\NormalTok{, }\StringTok{\textquotesingle{}gasto\_mensual\textquotesingle{}}\NormalTok{, }\StringTok{\textquotesingle{}num\_compras\_mes\textquotesingle{}}\NormalTok{,}
                \StringTok{\textquotesingle{}antiguedad\_meses\textquotesingle{}}\NormalTok{, }\StringTok{\textquotesingle{}ticket\_promedio\textquotesingle{}}\NormalTok{, }\StringTok{\textquotesingle{}visitas\_web\textquotesingle{}}\NormalTok{,}
                \StringTok{\textquotesingle{}frecuencia\_compra\textquotesingle{}}\NormalTok{, }\StringTok{\textquotesingle{}ratio\_gasto\_ingreso\textquotesingle{}}\NormalTok{]}

\NormalTok{X\_seg }\OperatorTok{=}\NormalTok{ datos\_clientes[features\_seg]}
\NormalTok{y\_seg }\OperatorTok{=}\NormalTok{ datos\_clientes[}\StringTok{\textquotesingle{}segmento\_num\textquotesingle{}}\NormalTok{]}

\NormalTok{X\_train\_seg, X\_test\_seg, y\_train\_seg, y\_test\_seg }\OperatorTok{=}\NormalTok{ train\_test\_split(}
\NormalTok{    X\_seg, y\_seg, test\_size}\OperatorTok{=}\FloatTok{0.3}\NormalTok{, random\_state}\OperatorTok{=}\DecValTok{42}\NormalTok{, stratify}\OperatorTok{=}\NormalTok{y\_seg}
\NormalTok{)}

\CommentTok{\# Entrenar Random Forest (mejor modelo típicamente)}
\BuiltInTok{print}\NormalTok{(}\StringTok{"}\CharTok{\textbackslash{}n\textbackslash{}n}\StringTok{MODELO DE SEGMENTACIÓN"}\NormalTok{)}
\BuiltInTok{print}\NormalTok{(}\StringTok{"="} \OperatorTok{*} \DecValTok{70}\NormalTok{)}

\NormalTok{rf\_segmentacion }\OperatorTok{=}\NormalTok{ RandomForestClassifier(}
\NormalTok{    n\_estimators}\OperatorTok{=}\DecValTok{200}\NormalTok{,}
\NormalTok{    max\_depth}\OperatorTok{=}\DecValTok{12}\NormalTok{,}
\NormalTok{    min\_samples\_split}\OperatorTok{=}\DecValTok{20}\NormalTok{,}
\NormalTok{    random\_state}\OperatorTok{=}\DecValTok{42}
\NormalTok{)}

\NormalTok{rf\_segmentacion.fit(X\_train\_seg, y\_train\_seg)}

\CommentTok{\# Evaluar}
\NormalTok{y\_pred\_seg }\OperatorTok{=}\NormalTok{ rf\_segmentacion.predict(X\_test\_seg)}
\NormalTok{acc\_seg }\OperatorTok{=}\NormalTok{ accuracy\_score(y\_test\_seg, y\_pred\_seg)}

\BuiltInTok{print}\NormalTok{(}\SpecialStringTok{f"Accuracy: }\SpecialCharTok{\{}\NormalTok{acc\_seg}\SpecialCharTok{:.4f\}}\SpecialStringTok{ (}\SpecialCharTok{\{}\NormalTok{acc\_seg}\OperatorTok{*}\DecValTok{100}\SpecialCharTok{:.1f\}}\SpecialStringTok{\%)"}\NormalTok{)}

\CommentTok{\# Matriz de confusión}
\NormalTok{cm\_seg }\OperatorTok{=}\NormalTok{ confusion\_matrix(y\_test\_seg, y\_pred\_seg)}
\BuiltInTok{print}\NormalTok{(}\SpecialStringTok{f"}\CharTok{\textbackslash{}n}\SpecialStringTok{Matriz de confusión:"}\NormalTok{)}
\BuiltInTok{print}\NormalTok{(}\SpecialStringTok{f"                    Predicho"}\NormalTok{)}
\BuiltInTok{print}\NormalTok{(}\SpecialStringTok{f"             Bronce  Plata   Oro  Platino"}\NormalTok{)}
\BuiltInTok{print}\NormalTok{(}\SpecialStringTok{f"Real Bronce    }\SpecialCharTok{\{}\NormalTok{cm\_seg[}\DecValTok{0}\NormalTok{, }\DecValTok{0}\NormalTok{]}\SpecialCharTok{:4d\}}\SpecialStringTok{   }\SpecialCharTok{\{}\NormalTok{cm\_seg[}\DecValTok{0}\NormalTok{, }\DecValTok{1}\NormalTok{]}\SpecialCharTok{:4d\}}\SpecialStringTok{  }\SpecialCharTok{\{}\NormalTok{cm\_seg[}\DecValTok{0}\NormalTok{, }\DecValTok{2}\NormalTok{]}\SpecialCharTok{:4d\}}\SpecialStringTok{    }\SpecialCharTok{\{}\NormalTok{cm\_seg[}\DecValTok{0}\NormalTok{, }\DecValTok{3}\NormalTok{]}\SpecialCharTok{:4d\}}\SpecialStringTok{"}\NormalTok{)}
\BuiltInTok{print}\NormalTok{(}\SpecialStringTok{f"     Plata     }\SpecialCharTok{\{}\NormalTok{cm\_seg[}\DecValTok{1}\NormalTok{, }\DecValTok{0}\NormalTok{]}\SpecialCharTok{:4d\}}\SpecialStringTok{   }\SpecialCharTok{\{}\NormalTok{cm\_seg[}\DecValTok{1}\NormalTok{, }\DecValTok{1}\NormalTok{]}\SpecialCharTok{:4d\}}\SpecialStringTok{  }\SpecialCharTok{\{}\NormalTok{cm\_seg[}\DecValTok{1}\NormalTok{, }\DecValTok{2}\NormalTok{]}\SpecialCharTok{:4d\}}\SpecialStringTok{    }\SpecialCharTok{\{}\NormalTok{cm\_seg[}\DecValTok{1}\NormalTok{, }\DecValTok{3}\NormalTok{]}\SpecialCharTok{:4d\}}\SpecialStringTok{"}\NormalTok{)}
\BuiltInTok{print}\NormalTok{(}\SpecialStringTok{f"     Oro       }\SpecialCharTok{\{}\NormalTok{cm\_seg[}\DecValTok{2}\NormalTok{, }\DecValTok{0}\NormalTok{]}\SpecialCharTok{:4d\}}\SpecialStringTok{   }\SpecialCharTok{\{}\NormalTok{cm\_seg[}\DecValTok{2}\NormalTok{, }\DecValTok{1}\NormalTok{]}\SpecialCharTok{:4d\}}\SpecialStringTok{  }\SpecialCharTok{\{}\NormalTok{cm\_seg[}\DecValTok{2}\NormalTok{, }\DecValTok{2}\NormalTok{]}\SpecialCharTok{:4d\}}\SpecialStringTok{    }\SpecialCharTok{\{}\NormalTok{cm\_seg[}\DecValTok{2}\NormalTok{, }\DecValTok{3}\NormalTok{]}\SpecialCharTok{:4d\}}\SpecialStringTok{"}\NormalTok{)}
\BuiltInTok{print}\NormalTok{(}\SpecialStringTok{f"     Platino   }\SpecialCharTok{\{}\NormalTok{cm\_seg[}\DecValTok{3}\NormalTok{, }\DecValTok{0}\NormalTok{]}\SpecialCharTok{:4d\}}\SpecialStringTok{   }\SpecialCharTok{\{}\NormalTok{cm\_seg[}\DecValTok{3}\NormalTok{, }\DecValTok{1}\NormalTok{]}\SpecialCharTok{:4d\}}\SpecialStringTok{  }\SpecialCharTok{\{}\NormalTok{cm\_seg[}\DecValTok{3}\NormalTok{, }\DecValTok{2}\NormalTok{]}\SpecialCharTok{:4d\}}\SpecialStringTok{    }\SpecialCharTok{\{}\NormalTok{cm\_seg[}\DecValTok{3}\NormalTok{, }\DecValTok{3}\NormalTok{]}\SpecialCharTok{:4d\}}\SpecialStringTok{"}\NormalTok{)}

\CommentTok{\# Características más importantes}
\BuiltInTok{print}\NormalTok{(}\SpecialStringTok{f"}\CharTok{\textbackslash{}n\textbackslash{}n}\SpecialStringTok{CARACTERÍSTICAS MÁS IMPORTANTES PARA SEGMENTACIÓN"}\NormalTok{)}
\BuiltInTok{print}\NormalTok{(}\StringTok{"="} \OperatorTok{*} \DecValTok{70}\NormalTok{)}

\NormalTok{importancias\_seg }\OperatorTok{=}\NormalTok{ pd.DataFrame(\{}
    \StringTok{\textquotesingle{}caracteristica\textquotesingle{}}\NormalTok{: features\_seg,}
    \StringTok{\textquotesingle{}importancia\textquotesingle{}}\NormalTok{: rf\_segmentacion.feature\_importances\_}
\NormalTok{\}).sort\_values(}\StringTok{\textquotesingle{}importancia\textquotesingle{}}\NormalTok{, ascending}\OperatorTok{=}\VariableTok{False}\NormalTok{)}

\ControlFlowTok{for}\NormalTok{ \_, row }\KeywordTok{in}\NormalTok{ importancias\_seg.iterrows():}
\NormalTok{    barra }\OperatorTok{=} \StringTok{\textquotesingle{}█\textquotesingle{}} \OperatorTok{*} \BuiltInTok{int}\NormalTok{(row[}\StringTok{\textquotesingle{}importancia\textquotesingle{}}\NormalTok{] }\OperatorTok{*} \DecValTok{100}\NormalTok{)}
    \BuiltInTok{print}\NormalTok{(}\SpecialStringTok{f"  }\SpecialCharTok{\{}\NormalTok{row[}\StringTok{\textquotesingle{}caracteristica\textquotesingle{}}\NormalTok{]}\SpecialCharTok{:\textless{}25\}}\SpecialStringTok{ }\SpecialCharTok{\{}\NormalTok{row[}\StringTok{\textquotesingle{}importancia\textquotesingle{}}\NormalTok{]}\SpecialCharTok{:.4f\}}\SpecialStringTok{ }\SpecialCharTok{\{}\NormalTok{barra}\SpecialCharTok{\}}\SpecialStringTok{"}\NormalTok{)}

\BuiltInTok{print}\NormalTok{(}\SpecialStringTok{f"}\CharTok{\textbackslash{}n\textbackslash{}n}\SpecialStringTok{APLICACIÓN PRÁCTICA:"}\NormalTok{)}
\BuiltInTok{print}\NormalTok{(}\StringTok{"Este modelo permite:"}\NormalTok{)}
\BuiltInTok{print}\NormalTok{(}\StringTok{"  1. Clasificar automáticamente nuevos clientes"}\NormalTok{)}
\BuiltInTok{print}\NormalTok{(}\StringTok{"  2. Personalizar ofertas según segmento"}\NormalTok{)}
\BuiltInTok{print}\NormalTok{(}\StringTok{"  3. Identificar clientes con potencial de upgrade"}\NormalTok{)}
\BuiltInTok{print}\NormalTok{(}\StringTok{"  4. Optimizar estrategias de retención"}\NormalTok{)}
\end{Highlighting}
\end{Shaded}

\section{Ejercicios prácticos}\label{ejercicios-pruxe1cticos}

\subsection{Ejercicio: Optimización de
hiperparámetros}\label{ejercicio-optimizaciuxf3n-de-hiperparuxe1metros}

\begin{Shaded}
\begin{Highlighting}[]
\BuiltInTok{print}\NormalTok{(}\StringTok{"}\CharTok{\textbackslash{}n\textbackslash{}n\textbackslash{}n}\StringTok{EJERCICIO: OPTIMIZACIÓN DE HIPERPARÁMETROS"}\NormalTok{)}
\BuiltInTok{print}\NormalTok{(}\StringTok{"="} \OperatorTok{*} \DecValTok{70}\NormalTok{)}

\CommentTok{\# Grid search manual para Random Forest}
\NormalTok{param\_grid }\OperatorTok{=}\NormalTok{ \{}
    \StringTok{\textquotesingle{}n\_estimators\textquotesingle{}}\NormalTok{: [}\DecValTok{50}\NormalTok{, }\DecValTok{100}\NormalTok{, }\DecValTok{200}\NormalTok{],}
    \StringTok{\textquotesingle{}max\_depth\textquotesingle{}}\NormalTok{: [}\DecValTok{5}\NormalTok{, }\DecValTok{10}\NormalTok{, }\DecValTok{15}\NormalTok{, }\VariableTok{None}\NormalTok{],}
    \StringTok{\textquotesingle{}min\_samples\_split\textquotesingle{}}\NormalTok{: [}\DecValTok{2}\NormalTok{, }\DecValTok{10}\NormalTok{, }\DecValTok{20}\NormalTok{]}
\NormalTok{\}}

\NormalTok{mejores\_params }\OperatorTok{=} \VariableTok{None}
\NormalTok{mejor\_score }\OperatorTok{=} \DecValTok{0}

\BuiltInTok{print}\NormalTok{(}\StringTok{"}\CharTok{\textbackslash{}n}\StringTok{Buscando mejores hiperparámetros..."}\NormalTok{)}
\BuiltInTok{print}\NormalTok{(}\SpecialStringTok{f"Total combinaciones: }\SpecialCharTok{\{}\BuiltInTok{len}\NormalTok{(param\_grid[}\StringTok{\textquotesingle{}n\_estimators\textquotesingle{}}\NormalTok{]) }\OperatorTok{*} \BuiltInTok{len}\NormalTok{(param\_grid[}\StringTok{\textquotesingle{}max\_depth\textquotesingle{}}\NormalTok{]) }\OperatorTok{*} \BuiltInTok{len}\NormalTok{(param\_grid[}\StringTok{\textquotesingle{}min\_samples\_split\textquotesingle{}}\NormalTok{])}\SpecialCharTok{\}}\SpecialStringTok{"}\NormalTok{)}

\NormalTok{resultados\_grid }\OperatorTok{=}\NormalTok{ []}

\ControlFlowTok{for}\NormalTok{ n\_est }\KeywordTok{in}\NormalTok{ param\_grid[}\StringTok{\textquotesingle{}n\_estimators\textquotesingle{}}\NormalTok{]:}
    \ControlFlowTok{for}\NormalTok{ max\_d }\KeywordTok{in}\NormalTok{ param\_grid[}\StringTok{\textquotesingle{}max\_depth\textquotesingle{}}\NormalTok{]:}
        \ControlFlowTok{for}\NormalTok{ min\_split }\KeywordTok{in}\NormalTok{ param\_grid[}\StringTok{\textquotesingle{}min\_samples\_split\textquotesingle{}}\NormalTok{]:}
            
            \CommentTok{\# Entrenar modelo}
\NormalTok{            rf }\OperatorTok{=}\NormalTok{ RandomForestClassifier(}
\NormalTok{                n\_estimators}\OperatorTok{=}\NormalTok{n\_est,}
\NormalTok{                max\_depth}\OperatorTok{=}\NormalTok{max\_d,}
\NormalTok{                min\_samples\_split}\OperatorTok{=}\NormalTok{min\_split,}
\NormalTok{                random\_state}\OperatorTok{=}\DecValTok{42}
\NormalTok{            )}
            
            \CommentTok{\# Validación cruzada}
\NormalTok{            scores }\OperatorTok{=}\NormalTok{ cross\_val\_score(rf, X\_train, y\_train, cv}\OperatorTok{=}\DecValTok{5}\NormalTok{)}
\NormalTok{            score\_mean }\OperatorTok{=}\NormalTok{ scores.mean()}
            
\NormalTok{            resultados\_grid.append(\{}
                \StringTok{\textquotesingle{}n\_estimators\textquotesingle{}}\NormalTok{: n\_est,}
                \StringTok{\textquotesingle{}max\_depth\textquotesingle{}}\NormalTok{: max\_d }\ControlFlowTok{if}\NormalTok{ max\_d }\ControlFlowTok{else} \StringTok{\textquotesingle{}None\textquotesingle{}}\NormalTok{,}
                \StringTok{\textquotesingle{}min\_samples\_split\textquotesingle{}}\NormalTok{: min\_split,}
                \StringTok{\textquotesingle{}cv\_score\textquotesingle{}}\NormalTok{: score\_mean}
\NormalTok{            \})}
            
            \ControlFlowTok{if}\NormalTok{ score\_mean }\OperatorTok{\textgreater{}}\NormalTok{ mejor\_score:}
\NormalTok{                mejor\_score }\OperatorTok{=}\NormalTok{ score\_mean}
\NormalTok{                mejores\_params }\OperatorTok{=}\NormalTok{ \{}
                    \StringTok{\textquotesingle{}n\_estimators\textquotesingle{}}\NormalTok{: n\_est,}
                    \StringTok{\textquotesingle{}max\_depth\textquotesingle{}}\NormalTok{: max\_d,}
                    \StringTok{\textquotesingle{}min\_samples\_split\textquotesingle{}}\NormalTok{: min\_split}
\NormalTok{                \}}

\BuiltInTok{print}\NormalTok{(}\SpecialStringTok{f"}\CharTok{\textbackslash{}n\textbackslash{}n}\SpecialStringTok{MEJORES HIPERPARÁMETROS"}\NormalTok{)}
\BuiltInTok{print}\NormalTok{(}\StringTok{"="} \OperatorTok{*} \DecValTok{70}\NormalTok{)}
\BuiltInTok{print}\NormalTok{(}\SpecialStringTok{f"  n\_estimators: }\SpecialCharTok{\{}\NormalTok{mejores\_params[}\StringTok{\textquotesingle{}n\_estimators\textquotesingle{}}\NormalTok{]}\SpecialCharTok{\}}\SpecialStringTok{"}\NormalTok{)}
\BuiltInTok{print}\NormalTok{(}\SpecialStringTok{f"  max\_depth: }\SpecialCharTok{\{}\NormalTok{mejores\_params[}\StringTok{\textquotesingle{}max\_depth\textquotesingle{}}\NormalTok{]}\SpecialCharTok{\}}\SpecialStringTok{"}\NormalTok{)}
\BuiltInTok{print}\NormalTok{(}\SpecialStringTok{f"  min\_samples\_split: }\SpecialCharTok{\{}\NormalTok{mejores\_params[}\StringTok{\textquotesingle{}min\_samples\_split\textquotesingle{}}\NormalTok{]}\SpecialCharTok{\}}\SpecialStringTok{"}\NormalTok{)}
\BuiltInTok{print}\NormalTok{(}\SpecialStringTok{f"  CV Score: }\SpecialCharTok{\{}\NormalTok{mejor\_score}\SpecialCharTok{:.4f\}}\SpecialStringTok{"}\NormalTok{)}

\CommentTok{\# Entrenar modelo final}
\NormalTok{rf\_optimo }\OperatorTok{=}\NormalTok{ RandomForestClassifier(}\OperatorTok{**}\NormalTok{mejores\_params, random\_state}\OperatorTok{=}\DecValTok{42}\NormalTok{)}
\NormalTok{rf\_optimo.fit(X\_train, y\_train)}

\NormalTok{acc\_optimo }\OperatorTok{=}\NormalTok{ rf\_optimo.score(X\_test, y\_test)}
\BuiltInTok{print}\NormalTok{(}\SpecialStringTok{f"}\CharTok{\textbackslash{}n}\SpecialStringTok{Accuracy en test: }\SpecialCharTok{\{}\NormalTok{acc\_optimo}\SpecialCharTok{:.4f\}}\SpecialStringTok{ (}\SpecialCharTok{\{}\NormalTok{acc\_optimo}\OperatorTok{*}\DecValTok{100}\SpecialCharTok{:.1f\}}\SpecialStringTok{\%)"}\NormalTok{)}

\CommentTok{\# Top 10 combinaciones}
\NormalTok{df\_grid }\OperatorTok{=}\NormalTok{ pd.DataFrame(resultados\_grid).sort\_values(}\StringTok{\textquotesingle{}cv\_score\textquotesingle{}}\NormalTok{, ascending}\OperatorTok{=}\VariableTok{False}\NormalTok{)}
\BuiltInTok{print}\NormalTok{(}\SpecialStringTok{f"}\CharTok{\textbackslash{}n\textbackslash{}n}\SpecialStringTok{Top 10 combinaciones:"}\NormalTok{)}
\BuiltInTok{print}\NormalTok{(df\_grid.head(}\DecValTok{10}\NormalTok{).to\_string(index}\OperatorTok{=}\VariableTok{False}\NormalTok{))}
\end{Highlighting}
\end{Shaded}

\section{Conclusión}\label{conclusiuxf3n}

En esta guía hemos explorado los principales algoritmos de
clasificación:

\textbf{Regresión Logística}

Modelo lineal simple y rápido. Usa función sigmoide para probabilidades.
Interpretable. Bueno como baseline.

\textbf{K-Nearest Neighbors (KNN)}

No paramétrico. Clasifica según vecinos más cercanos. Sensible a escala.
Bueno para fronteras complejas.

\textbf{Árboles de Decisión}

Reglas interpretables. No requiere normalización. Propenso a
overfitting. Base para métodos ensemble.

\textbf{Random Forests}

Ensemble de árboles. Reduce overfitting. Alta precisión. Menos
interpretable que árbol único.

\textbf{Comparación de modelos}

Importancia de evaluar múltiples algoritmos. Usar validación cruzada.
Considerar trade-offs interpretabilidad vs precisión.

\section{Próximos pasos}\label{pruxf3ximos-pasos}

En la siguiente guía (Guía 10: Métodos de ML para Regresión)
exploraremos:

\begin{itemize}
\tightlist
\item
  Regresión lineal
\item
  Ridge y Lasso
\item
  Regresión polinomial
\item
  Métricas: R², MSE, MAE
\item
  Aplicaciones económicas
\end{itemize}

\section{Recursos adicionales}\label{recursos-adicionales}

Para profundizar en clasificación:

\begin{itemize}
\tightlist
\item
  Scikit-learn Documentation: scikit-learn.org
\item
  Introduction to Statistical Learning (ISLR)
\item
  Hands-On Machine Learning (Aurélien Géron)
\item
  Papers en credit scoring y risk management
\end{itemize}

\section{Publicaciones Similares}\label{publicaciones-similares}

Si te interesó este artículo, te recomendamos que explores otros blogs y
recursos relacionados que pueden ampliar tus conocimientos. Aquí te dejo
algunas sugerencias:

\begin{enumerate}
\def\labelenumi{\arabic{enumi}.}
\tightlist
\item
  \href{https://numerus-scriptum.netlify.app/python/2020-06-19-instalacion-de-anaconda/index.pdf}{\faIcon{file-pdf}}
  \href{https://numerus-scriptum.netlify.app/python/2020-06-19-instalacion-de-anaconda}{Instalacion
  De Anaconda}
\item
  \href{https://numerus-scriptum.netlify.app/python/2020-06-20-configurar-entorno-virtual-python-anaconda/index.pdf}{\faIcon{file-pdf}}
  \href{https://numerus-scriptum.netlify.app/python/2020-06-20-configurar-entorno-virtual-python-anaconda}{Configurar
  Entorno Virtual Python Anaconda}
\item
  \href{https://numerus-scriptum.netlify.app/python/2021-04-17-01-introducion-a-la-programacion-con-python/index.pdf}{\faIcon{file-pdf}}
  \href{https://numerus-scriptum.netlify.app/python/2021-04-17-01-introducion-a-la-programacion-con-python}{01
  Introducion A La Programacion Con Python}
\item
  \href{https://numerus-scriptum.netlify.app/python/2021-05-31-02-variables-expresiones-y-statements-con-python/index.pdf}{\faIcon{file-pdf}}
  \href{https://numerus-scriptum.netlify.app/python/2021-05-31-02-variables-expresiones-y-statements-con-python}{02
  Variables Expresiones Y Statements Con Python}
\item
  \href{https://numerus-scriptum.netlify.app/python/2021-06-07-03-objetos-de-python/index.pdf}{\faIcon{file-pdf}}
  \href{https://numerus-scriptum.netlify.app/python/2021-06-07-03-objetos-de-python}{03
  Objetos De Python}
\item
  \href{https://numerus-scriptum.netlify.app/python/2021-06-14-04-ejecucion-condicional-con-python/index.pdf}{\faIcon{file-pdf}}
  \href{https://numerus-scriptum.netlify.app/python/2021-06-14-04-ejecucion-condicional-con-python}{04
  Ejecucion Condicional Con Python}
\item
  \href{https://numerus-scriptum.netlify.app/python/2021-06-21-05-iteraciones-con-python/index.pdf}{\faIcon{file-pdf}}
  \href{https://numerus-scriptum.netlify.app/python/2021-06-21-05-iteraciones-con-python}{05
  Iteraciones Con Python}
\item
  \href{https://numerus-scriptum.netlify.app/python/2021-08-16-06-funciones-con-python/index.pdf}{\faIcon{file-pdf}}
  \href{https://numerus-scriptum.netlify.app/python/2021-08-16-06-funciones-con-python}{06
  Funciones Con Python}
\item
  \href{https://numerus-scriptum.netlify.app/python/2021-08-23-07-dataframes-con-python/index.pdf}{\faIcon{file-pdf}}
  \href{https://numerus-scriptum.netlify.app/python/2021-08-23-07-dataframes-con-python}{07
  Dataframes Con Python}
\item
  \href{https://numerus-scriptum.netlify.app/python/2021-11-29-08-prediccion-y-metrica-de-performance-con-python/index.pdf}{\faIcon{file-pdf}}
  \href{https://numerus-scriptum.netlify.app/python/2021-11-29-08-prediccion-y-metrica-de-performance-con-python}{08
  Prediccion Y Metrica De Performance Con Python}
\item
  \href{https://numerus-scriptum.netlify.app/python/2021-12-06-09-metodos-de-machine-learning-para-clasificacion-con-python/index.pdf}{\faIcon{file-pdf}}
  \href{https://numerus-scriptum.netlify.app/python/2021-12-06-09-metodos-de-machine-learning-para-clasificacion-con-python}{09
  Metodos De Machine Learning Para Clasificacion Con Python}
\item
  \href{https://numerus-scriptum.netlify.app/python/2021-12-13-10-metodos-de-machine-learning-para-regresion-con-python/index.pdf}{\faIcon{file-pdf}}
  \href{https://numerus-scriptum.netlify.app/python/2021-12-13-10-metodos-de-machine-learning-para-regresion-con-python}{10
  Metodos De Machine Learning Para Regresion Con Python}
\item
  \href{https://numerus-scriptum.netlify.app/python/2022-10-31-11-validacion-cruzada-y-composicion-del-modelo-con-python/index.pdf}{\faIcon{file-pdf}}
  \href{https://numerus-scriptum.netlify.app/python/2022-10-31-11-validacion-cruzada-y-composicion-del-modelo-con-python}{11
  Validacion Cruzada Y Composicion Del Modelo Con Python}
\item
  \href{https://numerus-scriptum.netlify.app/python/2025-05-10-visualizacion-de-datos-con-python/index.pdf}{\faIcon{file-pdf}}
  \href{https://numerus-scriptum.netlify.app/python/2025-05-10-visualizacion-de-datos-con-python}{Visualizacion
  De Datos Con Python}
\end{enumerate}

Esperamos que encuentres estas publicaciones igualmente interesantes y
útiles. ¡Disfruta de la lectura!






\end{document}
